% !TEX program = xelatex
% 从张量到并矢 · 整本讲义 · round-01（2026-11-11）
\documentclass[11pt]{ctexart}
%=====================================================================
% theme-violet.tex —— 讲义精讲 · 紫花主题
% 用法：\documentclass[11pt]{ctexart}
%       \input{assets/theme-violet}
%       \begin{document} ... \end{document}
% 只用 ctexart 自带能力 + xcolor/enumitem/fancyhdr/pifont/amsmath，
% 不依赖 tikz/tcolorbox，保证 MiKTeX 可直接编译。
%=====================================================================

\usepackage[a4paper,left=2.0cm,right=2.0cm,top=2.1cm,bottom=2.3cm,%
            headsep=12pt,footskip=22pt]{geometry}
\usepackage{amsmath,amssymb}
\usepackage{needspace}
\usepackage{xcolor}
\usepackage{array}
\usepackage{enumitem}
\usepackage{pifont}
\usepackage{fancyhdr}
\usepackage[hidelinks]{hyperref}

%---------------------------------------------------------------- 配色
\definecolor{violet}{RGB}{107,78,143}      % 主色：紫
\definecolor{violetdeep}{RGB}{63,44,92}    % 深紫：标题/强调
\definecolor{petal}{RGB}{183,159,208}      % 花瓣：装饰线、花饰
\definecolor{bandviolet}{RGB}{247,243,251} % 浅紫底：档案/提示框
\definecolor{bandcream}{RGB}{252,249,242}  % 米白底：易错/注意框
\definecolor{rulepurple}{RGB}{217,207,230} % 分隔线
\definecolor{inkgray}{RGB}{92,96,110}      % 次要文字

%---------------------------------------------------------------- 字体
% 有 Noto CJK 用 Noto，没有就退回 ctex 默认（Windows 下为宋体/黑体）
\IfFontExistsTF{Noto Serif CJK SC}{\setCJKmainfont{Noto Serif CJK SC}}{}
\IfFontExistsTF{Noto Sans CJK SC}{\setCJKsansfont{Noto Sans CJK SC}}{}

%---------------------------------------------------------------- 版式
\setlist{nosep,leftmargin=1.9em}
\linespread{1.28}
\setlength{\parskip}{3.5pt}
\setlength{\parindent}{0pt}
\setlength{\tabcolsep}{5pt}
\renewcommand{\arraystretch}{1.28}

%---------------------------------------------------------------- 花饰
% \ding{95} = 五瓣小花，\ding{96} = 星形花；都来自 pifont，必定存在
\newcommand{\orn}{\textcolor{petal}{\ding{95}}}          % 行内小标
\newcommand{\ornstar}{\textcolor{petal}{\ding{96}}}      % 星形花
\newcommand{\flowerrule}{%                                % 居中花饰分隔线
  \par\vspace{7pt}%
  \noindent\makebox[\linewidth]{\color{petal}%
    \rule{0.20\linewidth}{0.6pt}\hspace{8pt}{\small\ding{95}}\hspace{4pt}%
    \raisebox{1.1pt}{\tiny\ding{95}}\hspace{4pt}{\small\ding{95}}\hspace{8pt}%
    \rule{0.20\linewidth}{0.6pt}}%
  \par\vspace{7pt}}

%---------------------------------------------------------------- 标题与节头
\newcommand{\doctitle}[3]{%   #1 主标题  #2 副标题  #3 元信息（可用 \\ 分行）
  \begin{center}
    {\color{petal}\small \orn\enspace #2\enspace \ornstar}\\[7pt]
    {\LARGE\bfseries\color{violetdeep} #1}\\[10pt]
    {\color{petal}\rule{0.55\linewidth}{1pt}}\\[4pt]
    {\small\color{inkgray} #3}
  \end{center}\vspace{2pt}}

\newcommand{\hdr}[2]{%   #1 章节号（如 §0）  #2 标题
  \par\vspace{12pt}%
  \noindent\colorbox{violetdeep}{\color{white}\bfseries\small\;#1\;}\hspace{7pt}%
  {\color{violetdeep}\large\bfseries #2}\par\nobreak
  \vspace{3pt}{\color{petal}\hrule height 1.2pt}\vspace{6pt}}

\newcommand{\subhdr}[1]{%
  \par\vspace{6pt}%
  \noindent{\color{violet}\bfseries #1}\par\vspace{1pt}}

%---------------------------------------------------------------- 提示框
\newcommand{\callout}[1]{% 浅紫底：讲义档案、读法说明
  \par\medskip\noindent
  \fcolorbox{rulepurple}{bandviolet}{\parbox{\dimexpr\linewidth-2\fboxsep-2\fboxrule}{\vspace{2pt}#1\vspace{2pt}}}%
  \par\medskip}

\newcommand{\warnbox}[1]{% 米白底：易错、边界条件
  \par\medskip\noindent
  \fcolorbox{rulepurple}{bandcream}{\parbox{\dimexpr\linewidth-2\fboxsep-2\fboxrule}{\vspace{2pt}#1\vspace{2pt}}}%
  \par\medskip}

\newcommand{\gapitem}[2]{% #1 缺口类型（紫） #2 内容
  \par\vspace{2pt}\noindent
  \hangindent=1.4em\hangafter=1
  {\color{petal}\small\ding{96}}\enspace{\bfseries\color{violetdeep}#1}\enspace #2\par}

% 题干块（\qstem）：每道例题／作业的「详解」开头必须先复述题干与出处。
% 起因：2026-10-25 第五章讲义把七道作业的题面集中放在一页、详解另起一页，
% 读者翻到「第 4 题」时看不到条件，等于一道没有题干的题目。
% 规则：详解块内自带题干，读者不翻页就能读懂题在问什么。
\newcommand{\qstem}[1]{%
  \par\medskip\noindent
  \fcolorbox{violet}{bandviolet}{%
    \parbox{\dimexpr\linewidth-2\fboxsep-2\fboxrule}{%
      \vspace{2pt}%
      {\small\bfseries\color{violet}题干}\hspace{0.8em}#1%
      \vspace{2pt}}}%
  \par\medskip}

% v2: 保护题干与开头，长题可跨页，续页补上下文。
\newcommand{\qbreak}{\needspace{6\baselineskip}}
% 独立练习：ID、来源、题干、考点、答案标签（通常 ans:ID）。
\newcommand{\practice}[5]{%
  \qbreak\subhdr{练习 #1}\qstem{（#2）#3}%
  \textbf{考点}：#4\par
  详解见附录第~\pageref{#5}~页。\par\medskip}
% 附录开头：ID、来源、条件摘要；自动产生 ans:ID 标签。
\newcommand{\answerstart}[3]{%
  \qbreak\subhdr{#1\quad 详解}\label{ans:#1}\qstem{（#2）#3}}
\newcommand{\answercontinue}[2]{%
  \needspace{5\baselineskip}\subhdr{#1\quad 续解}%
  {\small\color{inkgray}本页接续：#2}\par\medskip}
% 延迟提供默认值，让旧讲义可在导言区保留自己的 \newcommand{\step}。
\AtBeginDocument{\providecommand{\step}[1]{\needspace{3\baselineskip}\par\medskip\textbf{#1}\par}}

% 前提旁注：只放在「需要理解的那一步」后面，给陈述 + 本题代入，不重证。
\newcommand{\aside}[1]{%
  \par\vspace{1pt}\noindent
  \hangindent=1.6em\hangafter=1
  {\small\color{inkgray}%
    \hspace*{0.6em}{\color{petal}$\triangleright$}\hspace{2pt}#1}%
  \par\vspace{2pt}}

%---------------------------------------------------------------- 表格
% 本机 MiKTeX 的 colortbl 与 array 版本不匹配（\rowcolor 会报错），
% 故表格统一用「上粗紫线 + 表头 + 中细线」的三线式，不依赖 colortbl。
\newcommand{\tabtop}{\par\vspace{3pt}{\color{violet}\hrule height 1.2pt}\vspace{1pt}}
\newcommand{\tabmid}{\par\vspace{1pt}{\color{petal}\hrule height 0.5pt}\vspace{1pt}}
\newcommand{\tabbot}{\par\vspace{3pt}{\color{rulepurple}\hrule height 0.5pt}\vspace{8pt}}
\newcommand{\thead}[1]{{\bfseries\color{violetdeep}#1}}

%---------------------------------------------------------------- 页眉页脚
\pagestyle{fancy}
\fancyhf{}
\renewcommand{\headrulewidth}{0.4pt}
\renewcommand{\headrule}{\hbox to\headwidth{\color{rulepurple}\leaders\hrule height \headrulewidth\hfill}}
\fancyhead[L]{\small\color{inkgray}\songti \doctag}
\fancyhead[R]{\small\color{inkgray}\songti \coursetag}
\fancyfoot[C]{\small\color{inkgray}\orn\enspace{\color{violet}\thepage}\enspace\orn}

% 每篇讲义在导言区设这两个值
\newcommand{\doctag}{}
\newcommand{\coursetag}{}
\newcommand{\settags}[2]{\renewcommand{\doctag}{#1}\renewcommand{\coursetag}{#2}}

%---------------------------------------------------------------- 数学记号
\newcommand{\vect}[1]{\boldsymbol{#1}}
\newcommand{\dd}{\,\mathrm{d}}
\newcommand{\ddx}{\frac{\mathrm{d}}{\mathrm{d}x}}
\addtolength{\jot}{2pt}



\usepackage{tikz}
\usetikzlibrary{arrows.meta,calc}
\setlength{\headheight}{22pt}
\linespread{1.24}
\raggedbottom
\settags{专题 · 从张量到并矢}{理论力学}

% 记号：矢量粗斜体，二阶张量粗正体
\newcommand{\ten}[1]{\mathbf{#1}}
\newcommand{\unitd}{\hat{\ten{1}}}
\newcommand{\tr}{\operatorname{tr}}
\newcommand{\R}{\mathbb{R}}
\newcommand{\figcap}[1]{\par\vspace{2pt}{\small\color{inkgray}#1}\par}
\newcommand{\step}[1]{\par\smallskip\noindent{\color{violetdeep}\bfseries #1}\hspace{0.35em}}
\tikzset{
  ax/.style={inkgray,densely dotted,->},
  vv/.style={thick,->,>={Stealth[length=1.7mm]}},
  vviolet/.style={violet,thick,->,>={Stealth[length=1.7mm]}},
  >=Stealth,
}

\let\orighdr\hdr
\renewcommand{\hdr}[2]{\par\needspace{10\baselineskip}\orighdr{#1}{#2}}
\let\origsubhdr\subhdr
\renewcommand{\subhdr}[1]{\par\needspace{3\baselineskip}\origsubhdr{#1}}

\setlength{\abovedisplayskip}{6pt plus 2pt minus 2pt}
\setlength{\belowdisplayskip}{6pt plus 2pt minus 2pt}
\setlength{\emergencystretch}{1.5em}

\begin{document}

\doctitle{从张量到并矢}
{二次型主轴、并矢运算与三处力学落点}
{专题讲义\quad 2026-11-11\quad 主要来源：TM2026-Topic8-P1 p8--p12、p15；RIT Whelan notes 10 pp.2--14}

\callout{\textbf{这份材料为谁写、怎么用。}
这份讲义面向“会写矢量分量、会做点乘与叉乘、会做矩阵乘矢量，但还没有系统学过二次型、特征值、张量与并矢”的读者。点积的几何含义、叉乘的方向、定积分与质量微元、质心与角动量这些项目，学习者尚未确认熟练程度，因此都在第一次真正用到的地方补讲，不默认已经会。
\par\smallskip
紧跟正文的\textbf{教学示范}给完整演算，用来学方法；\textbf{独立练习} P1--P12 只在正文给完整题干、考点与附录页码，详解统一放在附录 A。请先自己做，再翻答案。正文里的“选读”可以跳过，跳过后的继续位置会写明。
\par\smallskip
\textbf{取材与边界。}课件 TM2026-Topic8-P1 p8--p12 给出正交基、求和约定、叉乘与对偶张量、张量积（Dyad）、左右作用、线性算子与二次型、单位张量、旋转算子及其角速度；p15 的刚体速度关系 $v=V+\boldsymbol\Omega\times r$ 作为已用前提。RIT Whelan notes 10 的 pp.2--14 用来对照惯量张量的构造、谱分解与平行轴定理。\textbf{欧拉方程、对称陀螺、应力张量、输运定理}在课件对应页没有展开，本稿明确按自己的推导补出，不伪称课件原有；课件 p13--p18 的欧拉角、E8.01/E8.02、HW8.01--HW8.03 不属于本专题，不扩入。
\par\smallskip
本稿是 DSH 在旧稿范围与已验收片段基础上按 v2.3.0 规则独立重组，不是从零生成，也不是把旧稿或片段原样拼接。旧稿与片段的可用讲解、证明和例子，只有放回全书依赖、重新核对条件与记号后才保留。}

\warnbox{\textbf{全篇只有一条主线。}
\textbf{同一个物理对象，换坐标系时分量会变，但对象本身不变。}
矩阵是“对象在一组基下的账本”，张量是“对象本身”，并矢是“二阶张量的另一种写法”。\S0 说清账本与换基；\S1 找出账本里不随基变的数与主轴；\S2 把对象写成并矢并建立运算；\S3--\S6 用它处理旋转、惯量、陀螺与应力。读每一节时都可以问一句：现在是在描述对象，还是在换一套账本？}

\hdr{§0}{起点：矢量、矩阵与换基}
\label{sec:tools}

这一节不是线性代数课，只把后面反复用到的几件事练一遍：写分量、点乘与投影、叉乘与方向、矩阵作用、换基、迹与行列式的含义，以及本轮实际用到的导数、积分与一个实根引理。每一件都在后面有明确去处，用到时会说明服务的是哪一步。

\subhdr{0.1\quad 矢量、点积、长度、正交与投影}
\label{subsec:dot}

取一组互相垂直的单位矢量 $\vect e_1,\vect e_2,\vect e_3$，叫一组\textbf{正交归一基}。任意矢量 $\vect a$ 都能唯一写成
\[
\vect a=a_1\vect e_1+a_2\vect e_2+a_3\vect e_3=\sum_{k=1}^{3}a_k\vect e_k .
\]
三个数 $a_1,a_2,a_3$ 是 $\vect a$ 在这组基下的\textbf{分量}。求法只有一条：拿基矢量去点它，$a_k=\vect a\cdot\vect e_k$。这是因为正交归一基满足 $\vect e_i\cdot\vect e_k=\delta_{ik}$，其中
\[
\delta_{ik}=\begin{cases}1,&i=k,\\0,&i\ne k,\end{cases}
\]
叫\textbf{克罗内克符号}。

\textbf{点积。}两个矢量的点积是一个数：
\[
\vect a\cdot\vect b=a_1b_1+a_2b_2+a_3b_3=\sum_{k=1}^{3}a_kb_k .
\]
它满足交换律，并且对每个因子都是线性的。点积不是“把两个矢量乘成一个矢量”；它的输出是一个数，这一点在 \S2 讲并矢时会反复对照。

\textbf{长度。}矢量自身的点积
\[
|\vect a|^2=\vect a\cdot\vect a=a_1^2+a_2^2+a_3^2
\]
是三个平方之和，所以 $|\vect a|^2\ge0$，且只有 $\vect a=\vect0$ 时才等于 $0$。定义 $|\vect a|=\sqrt{\vect a\cdot\vect a}$，这就是通常的几何长度：在三个坐标轴互相垂直的坐标系里，连续用两次勾股定理，原点到点 $(a_1,a_2,a_3)$ 的距离正好是 $\sqrt{a_1^2+a_2^2+a_3^2}$。

\textbf{点积与夹角。}设 $\vect a,\vect b$ 都非零，$\theta$ 是它们之间的夹角。在 $\vect a,\vect b$ 张成的平面里，第三边是 $\vect a-\vect b$。由平面几何的余弦定理，
\[
|\vect a-\vect b|^2=|\vect a|^2+|\vect b|^2-2|\vect a||\vect b|\cos\theta .
\]
另一方面，用点积展开：
\[
|\vect a-\vect b|^2
=(\vect a-\vect b)\cdot(\vect a-\vect b)
=|\vect a|^2+|\vect b|^2-2\,\vect a\cdot\vect b .
\]
两式比较，得到
\[
\vect a\cdot\vect b=|\vect a||\vect b|\cos\theta .
\]
这就是点积的几何含义：它把两个长度和夹角合成为一个数。特别地，若 $\vect n$ 是单位矢量，则 $\vect a\cdot\vect n=|\vect a|\cos\theta$，其中 $\theta$ 是 $\vect a$ 与 $\vect n$ 的夹角。

\textbf{垂直判据。}若 $\vect a\cdot\vect b=0$ 且 $\vect a,\vect b$ 非零，则 $|\vect a||\vect b|\ne0$，所以 $\cos\theta=0$，即 $\theta=90^\circ$，两矢量垂直（正交）。反过来，若 $\vect a\perp\vect b$，则 $\cos\theta=0$，所以 $\vect a\cdot\vect b=0$。因此
\[
\vect a\cdot\vect b=0
\quad\Longleftrightarrow\quad
\vect a\text{ 与 }\vect b\text{ 正交}\qquad(\vect a,\vect b\ne\vect0).
\]
“点积为零”不是数值碰巧为零，它正好是夹角余弦为零，所以几何上就是直角。

\textbf{具体核对。}在平面里取 $\vect a=(3,0)$、$\vect b=(3,4)$。则
\[
|\vect a|=3,\qquad |\vect b|=5,\qquad \vect a\cdot\vect b=9 .
\]
由 $\vect a\cdot\vect b=|\vect a||\vect b|\cos\theta$ 得 $\cos\theta=9/(3\cdot5)=3/5$，所以 $\theta\approx53.1^\circ$；$\vect b$ 在 $\vect a$ 方向上的投影长是 $|\vect b|\cos\theta=5\cdot(3/5)=3$，正好是 $\vect b$ 的水平分量。再取 $\vect c=(0,4)$：$\vect a\cdot\vect c=0$，所以 $\vect a\perp\vect c$，夹角是 $90^\circ$。这两个数都很小，可以口算；它们说明“点积、长度、垂直、投影”不是四套公式，而是同一件事的四种读法。

\textbf{沿单位方向的投影。}设 $\vect n$ 是单位矢量，$\vect a$ 任意。把 $\vect a$ 写成
\[
\vect a=(\vect a\cdot\vect n)\vect n+\vect r,
\qquad
\vect r=\vect a-(\vect a\cdot\vect n)\vect n .
\]
第一项沿 $\vect n$；第二项与 $\vect n$ 正交，因为
\[
\vect r\cdot\vect n
=\vect a\cdot\vect n-(\vect a\cdot\vect n)(\vect n\cdot\vect n)
=\vect a\cdot\vect n-\vect a\cdot\vect n=0 .
\]
所以 $\vect a\cdot\vect n$ 是 $\vect a$ 在 $\vect n$ 方向上的\textbf{带符号投影长度}：若 $\vect a\cdot\vect n>0$，投影与 $\vect n$ 同向；若小于 $0$，反向；它的绝对值才是投影的长度。投影矢量是 $(\vect a\cdot\vect n)\vect n$。这个“减去投影就得到正交补”的动作，\S1.4 证明 $\ten T$ 保持正交补时还会再用一次。

\textbf{小例。}取 $\vect a=(1,2,2)$，$\vect n=(1,1,1)/\sqrt3$。则
\[
\vect a\cdot\vect n=\frac{1+2+2}{\sqrt3}=\frac{5}{\sqrt3},
\qquad
(\vect a\cdot\vect n)\vect n=\frac53(1,1,1).
\]
于是
\[
\vect a-\frac53(1,1,1)=\left(-\frac23,\frac13,\frac13\right),
\]
它与 $\vect n$ 的点积为 $\frac{1}{\sqrt3}(-\frac23+\frac13+\frac13)=0$。这里 $\vect a\cdot\vect n=5/\sqrt3>0$，所以投影与 $\vect n$ 同向；剩余部分落在与 $\vect n$ 垂直的平面里。这就是“把 $\vect a$ 分成沿 $\vect n$ 的部分和垂直于 $\vect n$ 的部分”。

\subhdr{0.2\quad 叉乘：大小、方向与 $\varepsilon_{ijk}$}
\label{subsec:cross}

两个矢量的叉乘还是矢量。按分量写，
\[
\vect a\times\vect b=
\begin{pmatrix}
a_2b_3-a_3b_2\\
a_3b_1-a_1b_3\\
a_1b_2-a_2b_1
\end{pmatrix}.
\]
它满足 $\vect a\times\vect b=-\vect b\times\vect a$，且 $\vect a\cdot(\vect a\times\vect b)=0$、$\vect b\cdot(\vect a\times\vect b)=0$：叉乘结果与两个因子都垂直。大小是
\[
|\vect a\times\vect b|=|\vect a||\vect b|\sin\theta,
\]
也就是两矢量张成的平行四边形的面积。方向按右手定则：四指从 $\vect a$ 弯向 $\vect b$，拇指指向 $\vect a\times\vect b$。

\textbf{先算一个能判断方向的例子。}取 $\vect r=(1,0,0)$、$\vect F=(0,1,0)$。代入分量式：
\[
\vect r\times\vect F=
\begin{pmatrix}
0\cdot0-0\cdot1\\
0\cdot0-1\cdot0\\
1\cdot1-0\cdot0
\end{pmatrix}
=(0,0,1).
\]
所以 $\vect r\times\vect F$ 沿 $+\vect e_3$，即 $z$ 轴正方向。几何核对：$\vect r$ 沿 $x$ 轴，$\vect F$ 沿 $y$ 轴，右手四指从 $x$ 弯向 $y$，拇指正好指向 $z$。这个例子同时说明“叉乘结果垂直于原来的两个矢量”：$(0,0,1)\cdot(1,0,0)=0$，$(0,0,1)\cdot(0,1,0)=0$。

\textbf{力学里的两个叉乘。}把 $\vect F$ 换成力、$\vect r$ 换成从支点到作用点的位置矢量，$\vect M=\vect r\times\vect F$ 就是\textbf{力矩}：它衡量这个力使物体绕支点转动的倾向，方向由右手定则给出。把 $\vect F$ 换成动量 $\vect p=m\vect v$，$\vect L=\vect r\times\vect p$ 就是\textbf{角动量}。两者都是“位置矢量叉乘另一个矢量”，所以上面那个 $(1,0,0)\times(0,1,0)=(0,0,1)$ 的例子同时告诉你：沿 $x$ 轴的力臂配上沿 $y$ 轴的力，力矩沿 $z$ 轴。\S4 与 \S5 会反复用这两个量。

\textbf{再算一个一般例子。}取 $\vect a=(1,2,2)$、$\vect b=(2,1,0)$。则
\[
\vect a\times\vect b=(2\cdot0-2\cdot1,\;2\cdot2-1\cdot0,\;1\cdot1-2\cdot2)=(-2,4,-3).
\]
核对：$\vect a\cdot(\vect a\times\vect b)=-2+8-6=0$，且 $|\vect a\times\vect b|=\sqrt{4+16+9}=\sqrt{29}$。$\vect a$ 的长度是 $3$，$\vect b$ 的长度是 $\sqrt5$，所以 $\sin\theta=\sqrt{29}/(3\sqrt5)$；这也可以由 $\cos\theta=4/(3\sqrt5)$ 与 $\sin^2+\cos^2=1$ 核对。

\textbf{$\varepsilon_{ijk}$ 记号。}定义
\[
\varepsilon_{ijk}=
\begin{cases}
+1,&(i,j,k)\text{ 是 }(1,2,3)\text{ 的偶排列},\\
-1,&(i,j,k)\text{ 是 }(1,2,3)\text{ 的奇排列},\\
0,&\text{有重复指标}.
\end{cases}
\]
六个排列是 $123,231,312$（偶，取 $+1$）与 $132,213,321$（奇，取 $-1$）。于是
\[
(\vect a\times\vect b)_i=\varepsilon_{ijk}a_jb_k .
\]
这里用了求和约定：同一项里同一个指标出现两次，就对 $1,2,3$ 求和。核对第一分量：$\varepsilon_{123}a_2b_3+\varepsilon_{132}a_3b_2=a_2b_3-a_3b_2$，与上面的分量式一致。这个记号在 \S2.6 把叉乘写成“反对称张量在作用”时会再用一次。

\subhdr{0.3\quad 矩阵：线性映射的账本}
\label{subsec:matrix}

一个 $3\times3$ 实矩阵 $\ten T=(T_{ij})$ 作用在矢量 $\vect x=(x_1,x_2,x_3)$ 上，得到一个新矢量 $\ten T\vect x$，其第 $i$ 个分量是
\[
(\ten T\vect x)_i=\sum_{j=1}^{3}T_{ij}x_j=T_{ij}x_j .
\]
矩阵乘法按“行乘列”定义：
\[
(\ten A\ten B)_{ij}=\sum_{k=1}^{3}A_{ik}B_{kj}=A_{ik}B_{kj},
\]
它满足结合律 $(\ten A\ten B)\ten C=\ten A(\ten B\ten C)$。单位矩阵 $\unitd$ 满足 $\unitd\vect x=\vect x$，其分量为 $\delta_{ij}$。

\textbf{列是输入基矢量的像。}把 $\ten T\vect e_j$ 在基 $\{\vect e_i\}$ 下展开，系数记作 $T_{ij}$：
\[
\ten T\vect e_j=\sum_{i=1}^{3}T_{ij}\vect e_i
\qquad\Longleftrightarrow\qquad
(\ten T\vect x)_i=\sum_{j=1}^{3}T_{ij}x_j .
\]
把九个 $T_{ij}$ 按“第 $i$ 行、第 $j$ 列”排成方阵，就是 $\ten T$ 的矩阵。\textbf{第二个指标 $j$ 是输入基矢量的编号，第一个指标 $i$ 是输出分量的编号}；换句话说，\textbf{矩阵的第 $j$ 列就是 $\ten T(\vect e_j)$ 的三个分量}。这个约定决定了后面所有点乘的方向，混了就全错。

\textbf{数字例子。}设 $\ten T$ 把 $\vect e_1$ 变成 $2\vect e_1$，把 $\vect e_2$ 变成 $\vect e_1+\vect e_2$，把 $\vect e_3$ 变成 $3\vect e_3$：
\[
\ten T(\vect e_1)=\begin{pmatrix}2\\0\\0\end{pmatrix},\quad
\ten T(\vect e_2)=\begin{pmatrix}1\\1\\0\end{pmatrix},\quad
\ten T(\vect e_3)=\begin{pmatrix}0\\0\\3\end{pmatrix}
\quad\Longrightarrow\quad
\ten T=\begin{pmatrix}2&1&0\\0&1&0\\0&0&3\end{pmatrix}.
\]
注意这三个列矢量是按列摆进去的。取 $\vect x=(1,2,3)$：
\[
\ten T\vect x=\begin{pmatrix}2&1&0\\0&1&0\\0&0&3\end{pmatrix}\begin{pmatrix}1\\2\\3\end{pmatrix}
=\begin{pmatrix}2\cdot1+1\cdot2\\0\cdot1+1\cdot2\\0+0+3\cdot3\end{pmatrix}=\begin{pmatrix}4\\2\\9\end{pmatrix}.
\]
口算核对：$\ten T(\vect x)=1\cdot(2,0,0)+2\cdot(1,1,0)+3\cdot(0,0,3)=(4,2,9)$。两条路一致。

\textbf{矩阵乘法就是复合。}若先作用 $\ten B$ 再作用 $\ten A$，则 $(\ten A\ten B)\vect x=\ten A(\ten B\vect x)$。\textbf{转置}是行列互换：$(\ten A^{\mathsf T})_{ij}=A_{ji}$。它有一条后面反复用的判别式：
\[
(\ten A\vect x)\cdot\vect y=\vect x\cdot(\ten A^{\mathsf T}\vect y).
\]
证明只有一行：左边是 $\sum_{i,j}A_{ij}x_jy_i$，右边是 $\sum_{j,i}x_jA_{ij}y_i$，只是把数相乘的顺序换了。

\subhdr{0.4\quad 换基：账本变，机器不变}
\label{subsec:changebasis}

同一台机器，换一组基来记账，九个数字会变。这句话是本稿的枢纽，值得用一个具体例子钉死。

\textbf{二维例子。}机器 $\ten T$ 是“沿 $x$ 轴拉伸两倍”：$\ten T(x,y)=(2x,y)$。在标准基 $\vect e_1,\vect e_2$ 下
\[
\ten T=\begin{pmatrix}2&0\\0&1\end{pmatrix}.
\]
现在改用转了 $45^\circ$ 的基 $\vect n_1=(\vect e_1+\vect e_2)/\sqrt2$、$\vect n_2=(-\vect e_1+\vect e_2)/\sqrt2$。先算输入：$\ten T\vect n_1=\ten T(\frac{1}{\sqrt2},\frac{1}{\sqrt2})=(\frac{2}{\sqrt2},\frac{1}{\sqrt2})$。再把它在这组新基下分解，用 $\vect n_i\cdot$ 投影取系数：
\[
\ten T\vect n_1\cdot\vect n_1=\frac{2+1}{2}=\frac32,\qquad
\ten T\vect n_1\cdot\vect n_2=\frac{-2+1}{2}=-\frac12 .
\]
同法 $\ten T\vect n_2=(-\frac{2}{\sqrt2},\frac{1}{\sqrt2})$，
\[
\ten T\vect n_2\cdot\vect n_1=\frac{-2+1}{2}=-\frac12,\qquad
\ten T\vect n_2\cdot\vect n_2=\frac{2+1}{2}=\frac32 .
\]
于是同一台机器的账本变成
\[
\ten T'=\begin{pmatrix}3/2&-1/2\\-1/2&3/2\end{pmatrix}.
\]
两个账本长得完全不一样，可它们描述的是同一件事：沿 $\vect e_1$ 方向拉两倍、沿 $\vect e_2$ 方向不变。\textbf{判断“是不是同一台机器”，不能看数字，要看数字怎么变。}下面把这条换算规矩推出来，只有三步。

\step{第一步：每个新分量都是两次点乘}
回看上面的算法，新账本的每一个元素都长成同一个形状：
\[
T'_{ij}=\vect n_i\cdot(\ten T\vect n_j).
\]
内层那个点乘在\textbf{输入侧}：$\ten T$ 只认旧基坐标，所以得先把 $\vect n_j$ 写成旧基里的样子；外层那个点乘在\textbf{输出侧}：$\ten T$ 吐出来的仍是旧基坐标，得再投影到 $\vect n_i$ 上，才算新基下的第 $i$ 个分量。\textbf{两次翻译，就是后面那两个求和。}向量换基只乘一次，映射换基要乘两次：向量只有“分量”一个位置要翻，映射有“输入”和“输出”两个位置，各翻一次。

\step{第二步：把两次点乘按旧基写开}
记
\[
R_{ki}=\vect n_i\cdot\vect e_k,
\]
它的第 $i$ 列是 $\vect n_i$ 在旧基下的三个分量。于是 $\vect n_j=\sum_l R_{lj}\vect e_l$，喂给 $\ten T$ 之后
\[
\ten T\vect n_j=\sum_k\Bigl(\sum_l T_{kl}R_{lj}\Bigr)\vect e_k .
\]
最后拿 $\vect n_i=\sum_k R_{ki}\vect e_k$ 去点它，$R_{ki}$ 就是第二个因子：
\[
T'_{ij}=\sum_{k}R_{ki}\Bigl(\sum_{l}T_{kl}R_{lj}\Bigr)=\sum_{k,l}R_{ki}T_{kl}R_{lj}.
\]

\step{第三步：读成两个矩阵乘}
\[
\boxed{\;\ten T'=\ten R^{\mathsf T}\ten T\ten R\;},\qquad R_{ki}=\vect n_i\cdot\vect e_k .
\]
$\ten R$ 的第 $i$ 列是 $\vect n_i$ 在旧基下的分量；正交基之间 $\ten R^{\mathsf T}\ten R=\unitd$，所以 $\ten R^{-1}=\ten R^{\mathsf T}$。这个等式就是“新基也正交归一”的分量写法：$\ten R^{\mathsf T}\ten R$ 的 $(i,k)$ 元是 $\vect n_i\cdot\vect n_k=\delta_{ik}$。这个夹心结构怎么读，只看每个因子贴着哪一侧：

\begin{center}
\tabtop
\begin{tabular}{@{}p{0.24\linewidth}p{0.20\linewidth}p{0.48\linewidth}@{}}
\thead{因子} & \thead{贴在} & \thead{它在干什么} \\
\multicolumn{3}{@{}l@{}}{\color{petal}\rule{\linewidth}{0.5pt}}\\
右边的 $\ten R$ & 输入侧 & 把新基坐标翻成旧基坐标，$\ten T$ 才认 \\
中间的 $\ten T$ & 机器本身 & 它只会在旧账本上算 \\
左边的 $\ten R^{\mathsf T}$ & 输出侧 & 算完再翻回新基坐标 \\
\end{tabular}
\tabbot
\end{center}

一句话口诀：\textbf{指标在哪一侧，翻译就贴在那一侧}——$i$ 是输出指标，配左边的 $\ten R^{\mathsf T}$；$j$ 是输入指标，配右边的 $\ten R$。顺序不能对调：$\ten R\ten T\ten R^{\mathsf T}$ 是把基向相反方向转同样的角度后得到的账本，它与 $\ten T'$ 相似（迹、行列式、特征值都一样），但不是同一组基下的账本。

\textbf{矢量的新坐标是 $\ten R^{\mathsf T}$ 乘旧坐标。}这一点最容易写反，单独说清：
\[
x'_i=\vect n_i\cdot\vect x=\sum_k(\vect n_i\cdot\vect e_k)x_k=\sum_k R_{ki}x_k=(\ten R^{\mathsf T}\vect x)_i .
\]
所以\textbf{矢量贴一个 $\ten R^{\mathsf T}$，二阶张量左右各贴一个}。如果你读的书把 $\ten R$ 定义成“行是新基分量”，公式会写成 $\ten T'=\ten R\ten T\ten R^{\mathsf T}$、$\vect x'=\ten R\vect x$；两套只差一个转置，不要混用。本稿固定用“列是新基分量”的 $\ten R$。

\textbf{核对。}上面例子里 $\ten R=\frac{1}{\sqrt2}\begin{pmatrix}1&-1\\1&1\end{pmatrix}$，代入 $\ten R^{\mathsf T}\ten T\ten R$ 得 $\begin{pmatrix}3/2&-1/2\\-1/2&3/2\end{pmatrix}$，与逐项投影所得一致。这只是\textbf{核对}：规矩本身由上面三步推出，不是靠这个数值例子“试出来”的。

\subhdr{0.5\quad 迹、行列式与不随基变的数}
\label{subsec:invariants}

既然换个基数字全变，一个自然的问题是：\textbf{账本里有没有哪个组合换基后不变？}有，而且恰好是需要的那几个。

\textbf{迹。}方阵 $\ten A$ 的\textbf{迹}是对角元之和：
\[
\tr\ten A=\sum_i A_{ii}.
\]
它换基不变，理由只有一行。矩阵乘法有“循环性”：
\[
\tr(\ten A\ten B)=\sum_i\sum_j A_{ij}B_{ji}=\sum_j\sum_i B_{ji}A_{ij}=\tr(\ten B\ten A).
\]
取 $\ten B=\ten R$、$\ten A=\ten R^{\mathsf T}\ten T$，就有
\[
\tr(\ten R^{\mathsf T}\ten T\ten R)
=\tr\bigl((\ten R^{\mathsf T}\ten T)\ten R\bigr)
=\tr\bigl(\ten R(\ten R^{\mathsf T}\ten T)\bigr)
=\tr\bigl((\ten R\ten R^{\mathsf T})\ten T\bigr)
=\tr\ten T .
\]
最后一步用了 $\ten R\ten R^{\mathsf T}=\unitd$。\textbf{一行读完：迹不变，因为迹允许把因子转圈，而正交矩阵转一圈就消掉了。}数字核对：\S0.4 的两个账本都有 $\tr=3$。

\textbf{行列式。}$\det\ten T$ 是这台机器对\textbf{体积}的缩放倍数：把单位立方体喂给 $\ten T$，出来一个平行六面体，它的体积是 $|\det\ten T|$；符号表示是否翻转了手性。换坐标系不会改变“体积被放大几倍”，所以 $\det$ 也是不变量。代数上同样一句话：
\[
\det(\ten R^{\mathsf T}\ten T\ten R)=\det(\ten R^{\mathsf T})\det\ten T\det(\ten R)=\det\ten T,
\]
因为正交矩阵满足 $\det\ten R=\pm1$，平方为 $1$。$\det\ten T=0$ 意味着体积被压成零，机器不可逆；后面“某个主值为零”就是这种退化。

二维行列式的算法是 $\det\begin{pmatrix}a&b\\c&d\end{pmatrix}=ad-bc$。三维按第一行展开：
\[
\det\begin{pmatrix}a&b&c\\d&e&f\\g&h&i\end{pmatrix}
=a(ei-fh)-b(di-fg)+c(dh-eg).
\]
例如 $\det\begin{pmatrix}4&1&0\\1&4&0\\0&0&2\end{pmatrix}=4(8-0)-1(2-0)+0=30$，这个数在 P1 里会再用到。本稿只在具体数字里用到三维行列式，不要求背展开式，只要会按一行“正、负、正”交替地展开即可。

\begin{center}
\fcolorbox{violet}{bandviolet}{\parbox{0.94\linewidth}{\vspace{3pt}
\textbf{非零解判据。}设 $\ten A$ 是 $3\times3$ 实矩阵。齐次方程 $\ten A\vect x=\vect0$ 有非零解，当且仅当
\[
\det\ten A=0 .
\]
\vspace{3pt}}}
\end{center}

这条判据的“$\det\ne0$”一半可以用逆矩阵说明：若 $\det\ten A\ne0$，则 $\ten A$ 可逆，$\ten A\vect x=\vect0$ 两边左乘 $\ten A^{-1}$ 得 $\vect x=\vect0$。反过来，若 $\det\ten A=0$，用消元把 $\ten A$ 化成上三角；三个对角元不可能全非零（否则行列式非零），于是回代时出现自由未知量，可以取非零解。这里用的是简单线性代数里的消元事实，不引入秩或维数的新语言。这条判据在 \S1.3 建立特征方程、\S1.4 从特征值求特征向量时都会用到。

\textbf{特征多项式。}把“$\det$ 与 $\tr$ 都不变”这个念头推到底，就是问：$\det(\ten T-\lambda\unitd)$ 作为 $\lambda$ 的多项式是什么？它的系数是不变量，根将在 \S1.3 正式定义为特征值。二维先看清楚：$\ten T=\begin{pmatrix}a&b\\c&d\end{pmatrix}$，
\[
\det(\ten T-\lambda\unitd)=\det\begin{pmatrix}a-\lambda&b\\c&d-\lambda\end{pmatrix}
=(a-\lambda)(d-\lambda)-bc=\lambda^2-(a+d)\lambda+(ad-bc).
\]
所以二维的式子只有两个系数，分别就是 $\tr$ 与 $\det$：
\[
\lambda^2-(\tr\ten T)\lambda+\det\ten T=0 .
\]
三维照做，得到一个三次多项式：
\[
\det(\ten T-\lambda\unitd)=-\lambda^3+I_1\lambda^2-I_2\lambda+I_3 ,
\]
\[
I_1=\tr\ten T,\qquad
I_3=\det\ten T,\qquad
I_2=\tfrac12\bigl[(\tr\ten T)^2-\tr\ten T^2\bigr].
\]
三个系数 $I_1,I_2,I_3$ 都由 $\tr$ 与 $\det$ 拼成，所以三个都不随基变。

\textbf{一个容易误解的地方。}“$I_1,I_2,I_3$ 不变”说的是\textbf{换基时}不变，不是“$\ten T$ 不变”。换基时九个分量全变，只有这三个组合不变。特征值与主轴方向的关系在 \S1 讲；那里会说明：特征值不变，但主轴方向的\textbf{分量}会随基一起转。

\subhdr{0.6\quad 导数、定积分与质量微元}
\label{subsec:calculus}

这一小节只补本轮真正用到的微积分读法，不展开整门微积分。

\textbf{导数。}$\dd\vect A/\dd t$ 是矢量 $\vect A$ 在一小段时间里的改变量除以这段时间，再让时间趋于零。对矩阵逐项求导，就是对每个分量求导。矩阵乘积的求导法则与数的乘积法则相同：
\[
\frac{\dd}{\dd t}(\ten A\ten B)=\frac{\dd\ten A}{\dd t}\ten B+\ten A\frac{\dd\ten B}{\dd t},
\]
因为每个分量都是“两个数的乘积”，逐项用一次乘积法则再按矩阵乘法合并即可。\S3.4 对 $\ten S\ten S^{\mathsf T}=\unitd$ 求导时用的就是这一条。

\textbf{定积分。}定积分是把区间切成小段、把每段的贡献加起来再取极限。例如
\[
\int_0^a x^2\dd x=\frac{a^3}{3},
\]
因为 $x^3/3$ 的导数是 $x^2$。本稿的一维定积分只用这一种“先找原函数再代上下限”的算法，以及下面这种“切块求和”的读法。

\textbf{质量微元。}对一个连续体，把物体切成很多小块，每块的体积记作 $\dd V$，密度记作 $\rho$，则质量微元是
\[
\dd m=\rho\dd V .
\]
总和 $\int\dd m$ 就是总质量 $m$。若物体均匀，$\rho$ 是常数，$\int\dd m=\rho V$。三维的 $\int f\dd m$ 读作“把每个质量块的质量 $\dd m$ 乘上它对应的 $f$，再加起来”；它不是一个新的积分技巧，只是把“求和”写成了积分号。

\textbf{质心。}质心是质量加权平均位置：
\[
\vect R=\frac{1}{m}\int\vect r\dd m .
\]
把原点挪到质心，新的位置矢量 $\vect r'=\vect r-\vect R$ 满足
\[
\int\vect r'\dd m=\int\vect r\dd m-\vect R\int\dd m=m\vect R-m\vect R=\vect0 .
\]
这条性质在 \S4.5 的平行轴定理里会让交叉项整块消失。

\textbf{一维例子（后面会再用）。}均质细杆长 $L$、质量 $m$，沿 $x$ 轴放，中点在原点。线密度是 $m/L$，所以
\[
\int_{-L/2}^{L/2}x^2\dd m
=\int_{-L/2}^{L/2}x^2\frac{m}{L}\dd x
=\frac{m}{L}\cdot\frac{x^3}{3}\Big|_{-L/2}^{L/2}
=\frac{mL^2}{12}.
\]
这个结果就是细杆绕中点垂直轴的转动惯量；杆绕自身轴（$x$ 轴）的惯量为零，因为质量都在轴上。\S4.4 会直接引用它。

\subhdr{0.7\quad 奇次实多项式一定有实根}
\label{subsec:realroot}

\S1.4 需要一个存在性事实：一个三次实系数多项式至少有一个实根。下面的引理把这件事说清楚；\S1.3 定义的特征多项式正是这种多项式。

\begin{center}
\fcolorbox{violet}{bandviolet}{\parbox{0.94\linewidth}{\vspace{3pt}
\textbf{引理（奇次实多项式有实根）。}设
\[
p(t)=a_nt^n+a_{n-1}t^{n-1}+\cdots+a_0,\qquad a_n\ne0,
\]
其中 $n$ 是奇数，系数都是实数。则存在实数 $\lambda$ 使 $p(\lambda)=0$。
\vspace{3pt}}}
\end{center}

\textbf{证明。}若 $a_n<0$，把 $p$ 换成 $-p$；根不变，所以不妨设 $a_n>0$。当 $t>0$ 时，
\[
p(t)=t^n\left(a_n+\frac{a_{n-1}}{t}+\cdots+\frac{a_0}{t^n}\right).
\]
当 $t$ 足够大，括号内的量与 $a_n$ 的距离小于 $a_n/2$，括号为正，故 $p(t)>0$。当 $t<0$ 且 $|t|$ 足够大，$t^n<0$（因为 $n$ 为奇数），括号仍趋近于 $a_n>0$，故 $p(t)<0$。于是可以取到实数 $A<B$，使
\[
p(A)<0<p(B).
\]

多项式 $p$ 是连续函数。连续函数在闭区间 $[A,B]$ 上有介值性质：若两端函数值异号，则中间某点函数值为零。为看清这条性质，用二分法：令 $A_0=A$，$B_0=B$。取中点 $m_k$；若 $p(m_k)=0$，已经找到根；否则保留使 $p$ 在两端异号的那一半，记为 $[A_{k+1},B_{k+1}]$。这样得到嵌套区间
\[
[A_0,B_0]\supset[A_1,B_1]\supset[A_2,B_2]\supset\cdots,
\qquad
B_k-A_k=\frac{B-A}{2^k}\to0 .
\]
实数的完备性保证这些区间有唯一公共点 $\lambda$。因为 $p$ 连续，$p(A_k)$ 与 $p(B_k)$ 都趋于 $p(\lambda)$；而它们始终异号，所以 $p(\lambda)$ 既不能为正也不能为负，只能为 $0$。证完。

\begin{center}
\fcolorbox{rulepurple}{bandcream}{\parbox{0.94\linewidth}{\vspace{3pt}
\textbf{这条引理服务哪一步。}\S1.4 第二步要证明：$3\times3$ 实对称矩阵至少有一个实特征值。\S1.3 定义的特征多项式是三次实系数多项式，由本引理得到实根 $\lambda_1$；再由 \S0.5 的非零解判据得到特征向量。证明谱定理的存在性时，这是唯一的“分析”输入；后面的 $2\times2$ 情形完全是代数计算。
\vspace{3pt}}}
\end{center}

\hdr{§1}{二次型与主轴：椭球、特征值、谱定理}
\label{sec:quadratic}

现在回到本节的主问题。给定一个 $3\times3$ 实矩阵 $\ten T$，二次型
\[
E(\vect a)=\vect a\cdot\ten T\vect a
\]
在每一个方向 $\vect a$ 上给出一个数。我们想找一组互相垂直的坐标轴，使 $E$ 在这组轴上没有交叉项；如果找得到，这些轴就是\textbf{主轴}，轴上的系数就是\textbf{特征值}。本节的主线是：
\begin{enumerate}[label=\arabic*.,leftmargin=1.8em,itemsep=2pt]
\item 二次型只看见 $\ten T$ 的对称部分，所以可以假设 $\ten T$ 对称；
\item 在二维例子里，主轴让椭圆变正，交叉项消失，系数正好是特征值；
\item 特征值与特征向量怎么算；
\item 对任何实对称 $\ten T$，这种主轴一定存在——这就是谱定理，本节给出完整证明；
\item 用谱定理把二次型写成三个平方项，并说明怎样用三个不变量验算。
\end{enumerate}

\subhdr{1.1\quad 二次型与它的对称部分}
\label{subsec:quadratic-sym}

把 $\ten T$ 的分量记作 $T_{ik}$，$\vect a=(a_1,a_2,a_3)$。二次型写成指标和就是
\[
E(\vect a)=\vect a\cdot\ten T\vect a
=\sum_{i=1}^{3}\sum_{k=1}^{3}T_{ik}a_ia_k .
\]
写开就是
\[
E=T_{11}a_1^2+T_{22}a_2^2+T_{33}a_3^2
+(T_{12}+T_{21})a_1a_2
+(T_{13}+T_{31})a_1a_3
+(T_{23}+T_{32})a_2a_3 .
\]
请注意交叉项的系数是 $T_{12}+T_{21}$，不是无条件的一个 $2T_{12}$。只有当 $\ten T$ 对称，即 $T_{12}=T_{21}$ 时，才可以把它写成 $2T_{12}a_1a_2$。这个差别看起来只是记账，实际上决定了“二次型只与对称部分有关”这条结论。

\textbf{为什么只与对称部分有关。}在求和 $E=\sum_{i,k}T_{ik}a_ia_k$ 里，$i,k$ 都只是求和的哑指标。把两者的名字对换（只是换字母，值不变）：
\[
\sum_{i,k}T_{ik}a_ia_k
=\sum_{i,k}T_{ki}a_ka_i
=\sum_{i,k}T_{ki}a_ia_k .
\]
把原来的式子与这个式子相加再除以 $2$，得到
\[
E=\sum_{i,k}\frac{T_{ik}+T_{ki}}{2}\,a_ia_k
=\sum_{i,k}S_{ik}a_ia_k,
\qquad
S_{ik}=\frac{T_{ik}+T_{ki}}{2}.
\]
矩阵
\[
\ten S=\frac12\bigl(\ten T+\ten T^{\mathsf T}\bigr)
\]
叫 $\ten T$ 的\textbf{对称部分}。上面的计算说明：把 $\ten T$ 换成 $\ten S$，二次型 $E$ 一字不变。所以从这以后，凡讲二次型，都默认矩阵已经对称；如果题目给了一个不对称的 $\ten T$，先取 $\ten S$ 再谈主轴。

\textbf{教学小例。}取
\[
\ten T=\begin{pmatrix}1&2\\0&1\end{pmatrix},
\qquad
\vect a=\begin{pmatrix}x\\y\end{pmatrix}.
\]
直接算：
\[
E=\begin{pmatrix}x&y\end{pmatrix}
\begin{pmatrix}1&2\\0&1\end{pmatrix}
\begin{pmatrix}x\\y\end{pmatrix}
=x^2+2xy+y^2=(x+y)^2 .
\]
对称部分是
\[
\ten S=\frac12\left[
\begin{pmatrix}1&2\\0&1\end{pmatrix}
+\begin{pmatrix}1&0\\2&1\end{pmatrix}\right]
=\begin{pmatrix}1&1\\1&1\end{pmatrix}.
\]
用 $\ten S$ 算：
\[
E=\begin{pmatrix}x&y\end{pmatrix}
\begin{pmatrix}1&1\\1&1\end{pmatrix}
\begin{pmatrix}x\\y\end{pmatrix}
=x^2+xy+yx+y^2=x^2+2xy+y^2 .
\]
结果相同。这个例子特意保留了一个不对称的 $\ten T$，因为要看清：不是 $\ten T$ 的每个分量都不起作用，而是它的\textbf{反对称部分}在二次型里被抵消了。

\textbf{不要把两件事混在一起。}“二次型只看见对称部分”说的是 $E(\vect a)=\vect a\cdot\ten T\vect a$ 这个数。$\ten T$ 作为线性算子作用在矢量上时，反对称部分一般是有作用的；例如叉乘就可以写成一个反对称矩阵的作用（\S2.6）。本稿只在二次型的语境里把矩阵取成对称。

\textbf{物理上它长什么样。}刚体的转动动能是
\[
K=\frac12\vect\omega\cdot\ten J\vect\omega
=\frac12J_{ik}\omega_i\omega_k,
\]
正是一个二次型（差一个因子 $\frac12$）。这里 $\ten J$ 是 \S4 要构造的\textbf{转动惯量矩阵}；先把它当作一个实对称矩阵即可。它本来对称，所以“只取对称部分”在这里不损失任何信息。\S4 会把这个对应说完整；现在只需要记住：能量是 $\vect\omega$ 的二次函数，而二次函数的系数天然是一张 $3\times3$ 的表。

\subhdr{1.2\quad 等值面、椭球与主轴}
\label{subsec:ellipsoid}

先把主轴说准。\textbf{定义（特征值与特征向量）。}若存在非零矢量 $\vect n$ 和数 $\lambda$ 使
\[
\ten T\vect n=\lambda\vect n,
\]
就称 $\lambda$ 是 $\ten T$ 的\textbf{特征值}，$\vect n$ 是对应的\textbf{特征向量}；$\vect n$ 张成的直线叫一条\textbf{主轴方向}。零矢量不算特征向量，因为 $\ten T\vect0=\lambda\vect0$ 对任何 $\lambda$ 都成立，无法确定 $\lambda$。

几何上，$\ten T$ 把 $\vect n$ 张成的直线仍送到这条直线上：沿这条直线取带符号坐标 $t$（点是 $t\vect n$），则
\[
\ten T(t\vect n)=t\,\ten T\vect n=(\lambda t)\vect n,
\]
所以带符号坐标被乘以 $\lambda$，长度被乘以 $|\lambda|$。若 $\lambda>0$，方向不变；若 $\lambda<0$，方向反向；若 $\lambda=0$，整条直线被压到原点，$\vect n$ 仍按定义是特征向量，但 $\ten T\vect n$ 没有方向。因此准确地说：$\ten T$ 保持这条直线不变，沿它的带符号坐标乘以 $\lambda$；只有 $\lambda\ge0$ 时才能把“长度乘以 $\lambda$”当字面话，一般情况下长度乘以的是 $|\lambda|$。

现在令 $E(\vect a)=1$，得到一个曲面。先用二维看没有交叉项的情形：曲面是正放的椭圆；三维的椭球完全同理。若
\[
\ten T=\begin{pmatrix}1&0\\0&4\end{pmatrix},
\]
则
\[
E=x^2+4y^2=1
\quad\Longleftrightarrow\quad
\frac{x^2}{1^2}+\frac{y^2}{(1/2)^2}=1 .
\]
这是半轴分别为 $1$ 与 $1/2$ 的椭圆。两个对角元 $1,4$ 正好是 $\ten T$ 的两个特征值（直接看 $\ten T\vect e_1=\vect e_1$、$\ten T\vect e_2=4\vect e_2$）；沿第一个轴的半轴是 $1/\sqrt1=1$，沿第二个轴的半轴是 $1/\sqrt4=1/2$。一句话：\textbf{在 $\lambda>0$ 时，半轴等于 $1/\sqrt{\lambda}$，特征值越大，那个方向越短。}

\textbf{有交叉项的例子。}取
\[
\ten T=\begin{pmatrix}2&1\\1&2\end{pmatrix},
\qquad
E=2x^2+2xy+2y^2=1 .
\]
这个椭圆是斜的。先直接验证两个方向：
\[
\ten T\begin{pmatrix}1\\1\end{pmatrix}
=\begin{pmatrix}3\\3\end{pmatrix}
=3\begin{pmatrix}1\\1\end{pmatrix},
\qquad
\ten T\begin{pmatrix}-1\\1\end{pmatrix}
=\begin{pmatrix}-1\\1\end{pmatrix}
=1\begin{pmatrix}-1\\1\end{pmatrix}.
\]
所以 $\vect n_1=(1,1)/\sqrt2$ 对应特征值 $3$，$\vect n_2=(-1,1)/\sqrt2$ 对应特征值 $1$，两个方向正交。令
\[
u=\vect n_1\cdot(x,y)=\frac{x+y}{\sqrt2},
\qquad
v=\vect n_2\cdot(x,y)=\frac{-x+y}{\sqrt2}.
\]
反解是
\[
x=\frac{u-v}{\sqrt2},\qquad y=\frac{u+v}{\sqrt2}.
\]
代入：
\[
x^2+y^2=u^2+v^2,\qquad
xy=\frac{u^2-v^2}{2},
\]
所以
\[
E=2(x^2+y^2)+2xy
=2(u^2+v^2)+(u^2-v^2)
=3u^2+v^2 .
\]
交叉项真的消失了。在新的 $u,v$ 轴上，二次型是 $3u^2+v^2$；沿 $u$ 的半轴是 $1/\sqrt3$，沿 $v$ 的半轴是 $1$。让交叉项消失的那组坐标轴，就是\textbf{主轴}。

\begin{center}
\begin{minipage}{0.94\linewidth}
\centering
\begin{tikzpicture}[scale=2.35,>={Stealth[length=1.6mm]}]
  \begin{scope}[rotate=45]
    \fill[bandviolet,opacity=0.55] (0,0) ellipse ({0.5774} and {1});
    \draw[violet,thick] (0,0) ellipse ({0.5774} and {1});
  \end{scope}
  \draw[inkgray,densely dotted,->] (-1.25,0)--(1.30,0) node[right]{\scriptsize $x$};
  \draw[inkgray,densely dotted,->] (0,-1.20)--(0,1.30) node[above]{\scriptsize $y$};
  \draw[violet,thick,->] (0,0)--(0.72,0.72) node[above right]{\scriptsize $u$（$\vect n_1$）};
  \draw[violet,thick,->] (0,0)--(-0.86,0.86) node[above left]{\scriptsize $v$（$\vect n_2$）};
  \node[violet] at (0.56,0.29) {\scriptsize $\lambda_1=3$};
  \node[violet] at (-0.57,0.29) {\scriptsize $\lambda_2=1$};
  \fill (0,0) circle (0.7pt);
\end{tikzpicture}\\[2pt]
\figcap{\mbox{\textbf{图 1}（自绘）：}$E=2x^2+2xy+2y^2=1$ 的等值线。灰虚线是原来的 $x,y$ 轴；紫线是主轴 $u=\vect n_1$、$v=\vect n_2$。沿 $u$ 的特征值是 $3$，半轴只有 $1/\sqrt3\approx0.577$；沿 $v$ 的特征值是 $1$，半轴是 $1$。图中的椭圆长短轴与这两个数对应：\textbf{半轴 $=1/\sqrt{\lambda}$，特征值越大那一轴越短}。}
\end{minipage}
\end{center}

这张图与上面的代入逐项对应：$v=0$ 时 $(x,y)=u(1,1)/\sqrt2$，正是 $\vect n_1$ 方向；$u=0$ 时 $(x,y)=v(-1,1)/\sqrt2$，正是 $\vect n_2$ 方向。图中的 $u$ 轴沿 $45^\circ$，$v$ 轴沿 $135^\circ$，椭圆在 $v$ 方向更长，因为 $v$ 方向的半轴是 $1$，而 $u$ 方向的半轴只有 $0.577$。

\begin{center}
\fcolorbox{rulepurple}{bandcream}{\parbox{0.94\linewidth}{\vspace{3pt}
\textbf{边界情形。}上面的“半轴 $=1/\sqrt{\lambda}$”是在 $\lambda>0$ 时说的。若某个特征值为负，$E=1$ 在该方向上没有实点，曲面可能是双曲线或空集；若特征值为零，曲面退化（例如一对平行线或柱面）。但“主轴方向是特征向量、主值是特征值”仍然成立。谱定理只要求 $\ten T$ 实对称，不要求它正定；正定只保证 $E=1$ 是真椭球。惯量矩阵通常正定或半正定，应力矩阵则可能有负特征值——那时仍用主轴与主应力，但不画椭球。
\vspace{3pt}}}
\end{center}

\subhdr{1.3\quad 特征值与特征向量怎么算}
\label{subsec:eigen}

把 $\ten T\vect n=\lambda\vect n$ 移项：
\[
(\ten T-\lambda\unitd)\vect n=\vect0 .
\]
要它有非零解，由 \S0.5 的非零解判据（齐次方程有非零解当且仅当系数矩阵行列式为零），必须
\[
\det(\ten T-\lambda\unitd)=0 .
\]
左边是 $\lambda$ 的多项式，叫\textbf{特征多项式}；这个方程叫\textbf{特征方程}。对 $3\times3$ 矩阵，按第一行展开后它是三次多项式；对 $2\times2$ 矩阵，直接算得
\[
\det\begin{pmatrix}a-\lambda&b\\c&d-\lambda\end{pmatrix}
=\lambda^2-(a+d)\lambda+(ad-bc),
\]
这个二次公式在 \S1.4 的 $2\times2$ 引理里会直接用。解出特征值 $\lambda$ 后，再把它代回 $(\ten T-\lambda\unitd)\vect n=\vect0$，解出的非零 $\vect n$ 就是特征向量。

\textbf{三维教学示范。}取
\[
\ten D=\begin{pmatrix}6&2&0\\2&6&0\\0&0&1\end{pmatrix}.
\]
先算特征多项式：
\[
\det(\ten D-\lambda\unitd)
=(1-\lambda)\bigl[(6-\lambda)^2-4\bigr]
=(1-\lambda)(\lambda-4)(\lambda-8).
\]
所以特征值是
\[
\lambda_1=8,\qquad \lambda_2=4,\qquad \lambda_3=1.
\]
再求方向：
\begin{itemize}[leftmargin=1.8em,itemsep=2pt]
\item $\lambda=8$：$(\ten D-8\unitd)\vect n=\vect0$ 的前两行给出 $-2x+2y=0$，即 $x=y$；第三行给出 $-7z=0$，即 $z=0$。取
\[
\vect n_1=\frac{1}{\sqrt2}\begin{pmatrix}1\\1\\0\end{pmatrix}.
\]
\item $\lambda=4$：前两行给出 $2x+2y=0$，即 $y=-x$；第三行给出 $-3z=0$。取
\[
\vect n_2=\frac{1}{\sqrt2}\begin{pmatrix}1\\-1\\0\end{pmatrix}.
\]
\item $\lambda=1$：前两行是 $5x+2y=0$ 与 $2x+5y=0$，两式合起来只能 $x=y=0$；$z$ 任意。取
\[
\vect n_3=\begin{pmatrix}0\\0\\1\end{pmatrix}.
\]
\end{itemize}
三个方向两两正交，长度都是 $1$。把它们按“列乘行”的方式拼回去：
\[
8\,\vect n_1\vect n_1^{\mathsf T}
+4\,\vect n_2\vect n_2^{\mathsf T}
+1\,\vect n_3\vect n_3^{\mathsf T}
\]
\[
=8\cdot\frac12\begin{pmatrix}1&1&0\\1&1&0\\0&0&0\end{pmatrix}
+4\cdot\frac12\begin{pmatrix}1&-1&0\\-1&1&0\\0&0&0\end{pmatrix}
+\begin{pmatrix}0&0&0\\0&0&0\\0&0&1\end{pmatrix}
=\begin{pmatrix}6&2&0\\2&6&0\\0&0&1\end{pmatrix}.
\]
最后一步逐项相加即可。这里 $\vect n\vect n^{\mathsf T}$ 是“列乘行”得到的 $3\times3$ 矩阵，$(i,j)$ 元是 $n_in_j$；它作用在 $\vect x$ 上给出 $\vect n(\vect n\cdot\vect x)$，也就是把 $\vect x$ 投影到 $\vect n$ 方向。\S2.5 会把这个拼法写成一般公式。

\textbf{重复特征值。}例如
\[
\ten D=\operatorname{diag}(2,2,5).
\]
$\lambda=2$ 对应整个 $xy$ 平面：平面里任意非零矢量都满足 $\ten D\vect n=2\vect n$，所以任意一对正交单位矢量都可以当主轴；$\lambda=5$ 对应 $z$ 轴。重复特征值不破坏谱定理，只让“主轴的选择”不唯一。这里要区分两件事：特征值 $2$ 在特征多项式里出现两次（代数重数 $2$），而对应的特征空间是整个二维平面（几何重数 $2$）；对实对称矩阵，这两个重数总相等。若只给一个非对称矩阵，几何重数可能小于代数重数，那时就不能保证有足够多的正交主轴；本稿的谱定理正是把“实对称”这个条件用足。

\textbf{用迹和行列式抽查。}回到本节前面那个教学示范的 $\ten D=\begin{pmatrix}6&2&0\\2&6&0\\0&0&1\end{pmatrix}$，它的主对角线是 $6,6,1$，所以
\[
\tr\ten D=6+6+1=13=8+4+1,\qquad
\det\ten D=32=8\cdot4\cdot1 .
\]
这两个等式对任何矩阵都成立（特征值之和等于迹，特征值之积等于行列式；比较特征多项式按第一行展开的系数即可看出），可以用来抓符号错误。但它们只是\textbf{检查}，不是谱定理的证明；证明在 \S1.4。

\subhdr{1.4\quad 谱定理：陈述与完整证明}
\label{subsec:spectral}

\begin{center}
\fcolorbox{violet}{bandviolet}{\parbox{0.94\linewidth}{\vspace{4pt}
\textbf{谱定理（实对称矩阵，三维）。}设 $\ten T$ 是 $3\times3$ 实矩阵，且 $\ten T^{\mathsf T}=\ten T$。则：
\begin{enumerate}[label=(\roman*),leftmargin=2.0em,itemsep=3pt]
\item $\ten T$ 的三个特征值（按重数计）都是实数；
\item 属于不同特征值的特征向量互相正交；
\item 存在三个两两正交的单位矢量 $\vect n_1,\vect n_2,\vect n_3$ 和实数 $\lambda_1,\lambda_2,\lambda_3$，使
\[
\ten T\vect n_k=\lambda_k\vect n_k\quad(k=1,2,3),
\]
并且
\[
\ten T=\lambda_1\vect n_1\vect n_1^{\mathsf T}
+\lambda_2\vect n_2\vect n_2^{\mathsf T}
+\lambda_3\vect n_3\vect n_3^{\mathsf T}.
\]
等价地说：以 $\{\vect n_1,\vect n_2,\vect n_3\}$ 为基，$\ten T$ 的矩阵是对角阵 $\operatorname{diag}(\lambda_1,\lambda_2,\lambda_3)$。
\end{enumerate}
\vspace{4pt}}}
\end{center}

这里的 $\vect n_k\vect n_k^{\mathsf T}$ 就是 \S1.3 用过的“列乘行”矩阵；它不是装饰写法，而是一个真正的 $3\times3$ 矩阵。定理说：任何实对称矩阵都可以写成三个这样的投影矩阵的加权和，权重就是特征值。

\textbf{证明路线。}证明分五步：
\begin{enumerate}[label=\arabic*.,leftmargin=1.8em,itemsep=2pt]
\item 先完全解决 $2\times2$ 情形；
\item 用 \S0.7 的实根引理证明 $3\times3$ 实对称矩阵至少有一个实特征值；
\item 证明 $\ten T$ 保持这个特征向量的正交补，把问题降成 $2\times2$；
\item 用第 1 步的 $2\times2$ 引理得到另外两个正交特征向量；
\item 验证分解式，并补证“不同特征值的特征向量正交”与“没有非实特征值”。
\end{enumerate}
这样安排的原因很实际：$3\times3$ 的直接计算会遇到三次方程，而三次方程没有简单的求根公式；但一旦有了一个特征向量，剩下的平面就是二维的，$2\times2$ 情形可以用判别式完全算清楚。

\step{第一步：$2\times2$ 引理。}
\begin{center}
\fcolorbox{violet}{bandviolet}{\parbox{0.94\linewidth}{\vspace{3pt}
\textbf{引理（$2\times2$ 实对称矩阵）。}设
\[
\ten B=\begin{pmatrix}a&b\\b&d\end{pmatrix},\qquad a,b,d\in\R .
\]
则 $\ten B$ 有两个实特征值（可能相等），并且存在一组正交归一基，使 $\ten B$ 在这组基下的矩阵是对角阵。
\vspace{3pt}}}
\end{center}

\textbf{证明。}特征多项式是
\[
p_{\ten B}(\lambda)=\det(\ten B-\lambda\unitd)
=\lambda^2-(a+d)\lambda+(ad-b^2).
\]
判别式
\[
\Delta=(a+d)^2-4(ad-b^2)
=(a-d)^2+4b^2\ge0 .
\]
所以两个根
\[
\mu_\pm=\frac{(a+d)\pm\sqrt{\Delta}}{2}
\]
都是实数。

先处理 $b=0$ 的情形。此时 $\ten B=\operatorname{diag}(a,d)$，特征值是 $a,d$，对应的特征向量可以取 $\vect e_1,\vect e_2$。若 $a=d$，则 $\ten B=a\unitd$，任意单位矢量都是特征向量，任取一组正交归一基即可。

再处理 $b\ne0$ 的情形。此时 $\Delta=(a-d)^2+4b^2>0$，两个根 $\mu_+\ne\mu_-$。对每个根 $\mu$，令
\[
\vect v_\mu=\begin{pmatrix}b\\ \mu-a\end{pmatrix}.
\]
因为 $b\ne0$，$\vect v_\mu\ne\vect0$。验证它确实是特征向量：
\[
\ten B\vect v_\mu
=\begin{pmatrix}a&b\\b&d\end{pmatrix}
\begin{pmatrix}b\\ \mu-a\end{pmatrix}
=\begin{pmatrix}ab+b(\mu-a)\\ b^2+d(\mu-a)\end{pmatrix}
=\begin{pmatrix}b\mu\\ b^2+d\mu-ad\end{pmatrix}.
\]
由 $p_{\ten B}(\mu)=0$ 得
\[
\mu^2-(a+d)\mu+(ad-b^2)=0
\quad\Longrightarrow\quad
\mu^2-a\mu=d\mu-ad+b^2 .
\]
而第二个分量正好是
\[
b^2+d\mu-ad=\mu^2-a\mu=\mu(\mu-a),
\]
所以 $\ten B\vect v_\mu=\mu\vect v_\mu$。

最后看两个特征向量是否正交。利用根与系数的关系
\[
\mu_++\mu_-=a+d,\qquad
\mu_+\mu_-=ad-b^2,
\]
得到
\[
(\mu_+-a)(\mu_--a)
=\mu_+\mu_- -a(\mu_++\mu_-)+a^2
=(ad-b^2)-a(a+d)+a^2
=-b^2 .
\]
于是
\[
\vect v_+ \cdot \vect v_-
=b^2+(\mu_+-a)(\mu_--a)
=b^2-b^2=0 .
\]
把 $\vect v_\pm$ 分别单位化，就得到一组正交归一特征基。$2\times2$ 引理证完。

\step{第二步：$3\times3$ 实对称矩阵至少有一个实特征值。}
设 $\ten T$ 是 $3\times3$ 实对称矩阵。它的特征多项式
\[
p_{\ten T}(\lambda)=\det(\ten T-\lambda\unitd)
\]
按第一行展开后是三次实系数多项式，最高次项为 $-\lambda^3$。由 \S0.7 的引理，存在实数 $\lambda_1$ 使 $p_{\ten T}(\lambda_1)=0$。再由 \S0.5 的非零解判据，方程
\[
(\ten T-\lambda_1\unitd)\vect n=\vect0
\]
有非零解；取一个这样的解，单位化后记作 $\vect u_1$。于是
\[
\ten T\vect u_1=\lambda_1\vect u_1,\qquad |\vect u_1|=1 .
\]

\step{第三步：$\ten T$ 保持 $\vect u_1$ 的正交补。}
令
\[
P=\{\vect x\in\R^3:\vect x\cdot\vect u_1=0\}.
\]
这是过原点、与 $\vect u_1$ 垂直的平面。若 $\vect x\in P$，则用 \S0.3 的转置恒等式和 $\ten T^{\mathsf T}=\ten T$：
\[
(\ten T\vect x)\cdot\vect u_1
=\vect x\cdot(\ten T^{\mathsf T}\vect u_1)
=\vect x\cdot(\ten T\vect u_1)
=\lambda_1(\vect x\cdot\vect u_1)=0 .
\]
所以 $\ten T\vect x$ 仍在 $P$ 里，即 $\ten T(P)\subset P$。这一步是降维的关键：$\ten T$ 不会把平面 $P$ 里的矢量“翘出”平面。

\begin{center}
\begin{tikzpicture}[scale=1.02,>={Stealth[length=1.6mm]}]
  \fill[bandviolet,opacity=0.55] (-1.65,-0.45)--(1.20,-0.45)--(1.75,0.55)--(-1.15,0.55)--cycle;
  \draw[rulepurple] (-1.65,-0.45)--(1.20,-0.45)--(1.75,0.55)--(-1.15,0.55)--cycle;
  \draw[thick,->] (0,0)--(0,2.0) node[above]{\scriptsize $\vect u_1$};
  \draw[violet,thick,->] (0,0)--(0.95,-0.02) node[below]{\scriptsize $\vect x$};
  \draw[violet,thick,->] (0,0)--(0.30,0.38) node[above right]{\scriptsize $\ten T\vect x$};
  \node at (1.25,0.28) {\scriptsize $P=\vect u_1^{\perp}$};
  \fill (0,0) circle (0.6pt);
\end{tikzpicture}\\[2pt]
\figcap{\mbox{\textbf{图 2}（自绘，示意）：}$\vect u_1$ 与平面 $P$ 垂直；若 $\vect x\in P$，则 $\ten T\vect x$ 仍在 $P$ 内。图中只画方向与包含关系，不表示长度；它对应第三步的等式 $(\ten T\vect x)\cdot\vect u_1=0$。}
\end{center}

在 $P$ 内取一组正交归一基 $\vect e_2,\vect e_3$。这样的基一定存在：先取 $P$ 内任意非零矢量并单位化，记为 $\vect e_2$；再取 $P$ 内一个与 $\vect e_2$ 不平行的矢量 $\vect w$（平面是二维的，这样的 $\vect w$ 存在），令
\[
\vect z=\vect w-(\vect w\cdot\vect e_2)\vect e_2 .
\]
由 \S0.1 的投影公式，$\vect z\cdot\vect e_2=0$，且 $\vect z\ne\vect0$（否则 $\vect w$ 与 $\vect e_2$ 平行，矛盾）。把 $\vect z$ 单位化就得到 $\vect e_3$。因为 $\ten T(P)\subset P$，每个 $\ten T\vect e_j$ 都能写成 $\vect e_2,\vect e_3$ 的线性组合。定义
\[
B_{ij}=\vect e_i\cdot\ten T\vect e_j,\qquad i,j=2,3,
\]
则
\[
\ten T\vect e_j=\sum_{i=2}^{3}B_{ij}\vect e_i .
\]
矩阵 $\ten B=(B_{ij})$ 是 $2\times2$ 实矩阵，而且对称：
\[
B_{ij}=\vect e_i\cdot\ten T\vect e_j
=(\ten T\vect e_i)\cdot\vect e_j
=\vect e_j\cdot\ten T\vect e_i
=B_{ji},
\]
其中第二个等号用了 $\ten T$ 的对称性，第三个等号用了点积的交换律。于是 $\ten B$ 是 $2\times2$ 实对称矩阵。

\step{第四步：在平面内用 $2\times2$ 引理。}
先用一个具体对象看清这一步在做什么。取标准基 $\vect i=(1,0,0)$、$\vect j=(0,1,0)$、$\vect k=(0,0,1)$，并令
\[
\ten T=\begin{pmatrix}5&0&0\\0&2&1\\0&1&2\end{pmatrix}.
\]
它实对称，$\vect u_1=\vect i$ 是特征向量，$\ten T\vect u_1=5\vect u_1$，所以 $\lambda_1=5$，而 $P=\{\vect x:\vect x\cdot\vect u_1=0\}$ 就是 $yz$ 平面。在 $P$ 内取 $\vect e_2=\vect j$、$\vect e_3=\vect k$，则第三步的 $\ten B$ 是
\[
B_{ij}=\vect e_i\cdot\ten T\vect e_j
\quad\Longrightarrow\quad
\ten B=\begin{pmatrix}2&1\\1&2\end{pmatrix}.
\]
这个 $\ten B$ 正是 \S1.2 用过的 $2\times2$ 矩阵。$P$ 中的三维矢量 $y\vect e_2+z\vect e_3$ 用两个坐标 $(y,z)$ 表示，$\ten T$ 在 $P$ 上的作用由 $\ten B$ 完成：
\[
\ten T(y\vect e_2+z\vect e_3)
=(2y+z)\vect e_2+(y+2z)\vect e_3,
\]
右端的两个系数正是 $\ten B(y,z)^{\mathsf T}$。例如 $\ten B(1,0)^{\mathsf T}=(2,1)^{\mathsf T}$，拼回三维就是 $\ten T\vect e_2=2\vect e_2+\vect e_3$；$\ten B(1,1)^{\mathsf T}=3(1,1)^{\mathsf T}$ 则给出特征方向。$\ten B$ 的特征向量是
\[
\vect c_2=\frac{1}{\sqrt2}\begin{pmatrix}1\\1\end{pmatrix}\ (\mu_2=3),
\qquad
\vect c_3=\frac{1}{\sqrt2}\begin{pmatrix}1\\-1\end{pmatrix}\ (\mu_3=1).
\]
把二维特征向量拼回三维：
\[
\vect u_2=c_{22}\vect e_2+c_{32}\vect e_3
=\frac{1}{\sqrt2}(\vect e_2+\vect e_3)
=\frac{1}{\sqrt2}\begin{pmatrix}0\\1\\1\end{pmatrix},
\]
\[
\vect u_3=c_{23}\vect e_2+c_{33}\vect e_3
=\frac{1}{\sqrt2}(\vect e_2-\vect e_3)
=\frac{1}{\sqrt2}\begin{pmatrix}0\\1\\-1\end{pmatrix}.
\]
直接核对：
\[
\ten T\vect u_2=3\vect u_2,\qquad
\ten T\vect u_3=\vect u_3 .
\]
所以 $\vect u_1,\vect u_2,\vect u_3$ 是一组正交归一特征基，对应的特征值是 $5,3,1$。这个例子对应一般证明的第四步：$\ten B$ 的两个特征向量 $\vect c_2,\vect c_3$ 给出两个坐标方向，按 $\vect u_k=c_{2k}\vect e_2+c_{3k}\vect e_3$ 拼回 $P$ 中的三维矢量，就得到 $\ten T$ 的另外两个特征向量。它只演示对应关系，不替代下面的字母论证。

现在回到一般情形。对 $\ten B$ 应用第一步的引理：存在 $2\times2$ 的正交归一特征基 $\vect c_2,\vect c_3$ 和实数 $\mu_2,\mu_3$，使
\[
\ten B\vect c_k=\mu_k\vect c_k,\qquad k=2,3 .
\]
把 $\vect c_k=(c_{2k},c_{3k})$ 重新拼成三维矢量
\[
\vect u_k=c_{2k}\vect e_2+c_{3k}\vect e_3,\qquad k=2,3 .
\]
因为 $\vect c_2,\vect c_3$ 正交归一，$\vect u_2,\vect u_3$ 也在 $P$ 内正交归一。并且
\[
\ten T\vect u_k
=c_{2k}\ten T\vect e_2+c_{3k}\ten T\vect e_3
=\sum_{i=2}^{3}(\ten B\vect c_k)_i\vect e_i
=\mu_k\vect u_k .
\]
所以 $\vect u_2,\vect u_3$ 是 $\ten T$ 的特征向量，特征值分别是 $\mu_2,\mu_3$。为记号统一，把 $\mu_2,\mu_3$ 重新记作 $\lambda_2,\lambda_3$。又因为 $\vect u_2,\vect u_3\in P$，它们都与 $\vect u_1$ 正交。于是 $\vect u_1,\vect u_2,\vect u_3$ 是一组正交归一特征基。

在写最后的换基计算之前，先用已建立的分解核对一件事：只要特征值有一个实特征向量，它就属于 $\{\lambda_1,\mu_2,\mu_3\}$。设 $\vect x\ne\vect0$ 是 $\ten T$ 的一个实特征向量，特征值为 $\lambda$，即 $\ten T\vect x=\lambda\vect x$。把 $\vect x$ 沿 $\vect u_1$ 与平面 $P$ 分解：
\[
\vect x=\alpha\vect u_1+\vect p,\qquad
\alpha=\vect x\cdot\vect u_1,\qquad
\vect p=\vect x-\alpha\vect u_1\in P .
\]
因为 $\ten T\vect u_1=\lambda_1\vect u_1$，且第三步已证 $\ten T\vect p\in P$，
\[
\ten T\vect x=\alpha\lambda_1\vect u_1+\ten T\vect p .
\]
左边又等于 $\lambda\vect x=\lambda\alpha\vect u_1+\lambda\vect p$。比较两边的 $\vect u_1$ 分量和 $P$ 内分量，得到
\[
\alpha\lambda_1=\lambda\alpha,\qquad
\ten T\vect p=\lambda\vect p .
\]
若 $\alpha\ne0$，第一式给出 $\lambda=\lambda_1$。若 $\alpha=0$，则 $\vect x=\vect p\in P$ 且 $\vect p\ne\vect0$；把 $\vect p$ 写成 $\vect p=c_2\vect e_2+c_3\vect e_3$，第二式就是 $\ten B\vect c=\lambda\vect c$，其中 $\vect c=(c_2,c_3)\ne\vect0$，所以 $\lambda$ 是 $\ten B$ 的特征值，由第一步的 $2\times2$ 引理只能是 $\mu_2$ 或 $\mu_3$。

上面的检查只覆盖“有实特征向量”的特征值：它说明构造出的三个数不会漏掉任何实特征值。但特征多项式还可能有非实根，因此要得到 (i)，还需要把 $p_{\ten T}$ 与刚构造的特征基直接对上。令 $\ten Q$ 为以 $\vect u_1,\vect u_2,\vect u_3$ 为列的正交矩阵。因为 $\vect u_1,\vect u_2,\vect u_3$ 是正交归一特征基，$\ten Q^{\mathsf T}\ten Q=\unitd$，并且
\[
\ten Q^{\mathsf T}\ten T\ten Q
=\operatorname{diag}(\lambda_1,\lambda_2,\lambda_3)=:\ten D,
\qquad\text{即}\qquad
\ten T=\ten Q\ten D\ten Q^{\mathsf T}.
\]
于是对任意 $\lambda$，
\[
\ten T-\lambda\unitd=\ten Q(\ten D-\lambda\unitd)\ten Q^{\mathsf T}.
\]
这里用一次 $3\times3$ 行列式的乘法性质：对任意两个 $3\times3$ 矩阵 $\ten X,\ten Y$，$\det(\ten X\ten Y)=\det\ten X\,\det\ten Y$。本条只按 $3\times3$ 使用这条标准性质（行列式对矩阵乘法的基本性质），不在这里重证它。又因为 $\det\ten Q^{\mathsf T}\det\ten Q=\det(\ten Q^{\mathsf T}\ten Q)=\det\unitd=1$，所以
\[
\det(\ten T-\lambda\unitd)
=\det\ten Q\,\det(\ten D-\lambda\unitd)\,\det\ten Q^{\mathsf T}
=\det(\ten D-\lambda\unitd).
\]
而 $\ten D-\lambda\unitd$ 是对角阵，按第一行展开就是
\[
\det(\ten D-\lambda\unitd)
=(\lambda_1-\lambda)(\lambda_2-\lambda)(\lambda_3-\lambda).
\]
所以
\[
p_{\ten T}(\lambda)=(\lambda_1-\lambda)(\lambda_2-\lambda)(\lambda_3-\lambda).
\]
右边是三个一次因子的乘积，它的根（按重数计）正是 $\lambda_1,\lambda_2,\lambda_3$，全是实数。因此 (i) 成立。

\step{第五步：正交性与分解式。}
先补证 (ii)。设 $\ten T\vect u=\lambda\vect u$，$\ten T\vect v=\mu\vect v$，且 $\lambda\ne\mu$。则
\[
\lambda(\vect u\cdot\vect v)
=(\ten T\vect u)\cdot\vect v
=\vect u\cdot(\ten T^{\mathsf T}\vect v)
=\vect u\cdot(\ten T\vect v)
=\mu(\vect u\cdot\vect v).
\]
第二个等号是 \S0.3 的转置恒等式，第三个等号用了 $\ten T^{\mathsf T}=\ten T$。移项得
\[
(\lambda-\mu)(\vect u\cdot\vect v)=0 .
\]
因为 $\lambda\ne\mu$，只能 $\vect u\cdot\vect v=0$。这就证明了 (ii)。注意这一步没有用到“特征向量已经单位化”，只用了对称性。

最后证明 (iii) 的分解式。对每个基矢量 $\vect u_j$，右边作用上去：
\[
\left(\sum_{k=1}^{3}\lambda_k\vect u_k\vect u_k^{\mathsf T}\right)\vect u_j
=\sum_{k=1}^{3}\lambda_k\vect u_k(\vect u_k\cdot\vect u_j)
=\lambda_j\vect u_j,
\]
因为 $\vect u_k\cdot\vect u_j=\delta_{kj}$。左边 $\ten T$ 作用在 $\vect u_j$ 上也给出 $\lambda_j\vect u_j$。两个线性映射在基 $\vect u_1,\vect u_2,\vect u_3$ 上作用相同，所以它们相等：
\[
\ten T=\sum_{k=1}^{3}\lambda_k\vect u_k\vect u_k^{\mathsf T}.
\]
谱定理证完。

\begin{center}
\fcolorbox{rulepurple}{bandcream}{\parbox{0.94\linewidth}{\vspace{3pt}
\textbf{这个证明实际用了什么。}它只用了：$\ten T$ 实对称；三维特征多项式是三次实系数多项式；\S0.5 的非零解判据；\S0.7 的实根存在性；$2\times2$ 的判别式计算；以及第四步末尾一次 $3\times3$ 行列式的乘法性质 $\det(\ten X\ten Y)=\det\ten X\det\ten Y$（只用于把 $p_{\ten T}$ 写成 $\prod_{k=1}^{3}(\lambda_k-\lambda)$）。它没有用正定性，没有用 $\ten T$ 可逆，也没有用多元微积分或拉格朗日乘子。因此它对惯量矩阵（通常正定或半正定）和应力矩阵（可能有负特征值）同样适用。
\vspace{3pt}}}
\end{center}

\textbf{重复特征值时会发生什么。}若三个特征值全相等，$\ten T=\lambda_1\unitd$：由分解式，
\[
\ten T=\lambda_1(\vect u_1\vect u_1^{\mathsf T}
+\vect u_2\vect u_2^{\mathsf T}
+\vect u_3\vect u_3^{\mathsf T})
=\lambda_1\unitd,
\]
所以任意正交归一基都是特征基。若两个特征值相等，对应的特征空间是一个平面，平面内任意一组正交归一基都可以当主轴。谱定理保证的是“存在一组”，不是“只有一组”；说“主轴方向”时，若特征值重复，必须说明是在哪个特征空间里选。对实对称矩阵，代数重数与几何重数相等，这一点由上面的分解式直接看出：每个特征值在 $\operatorname{diag}(\lambda_1,\lambda_2,\lambda_3)$ 里出现的次数，就是它在特征空间里能贡献的正交方向数。

\subhdr{1.5\quad 谱定理的三个用法}
\label{subsec:uses}

谱定理的价值不只在“存在三个主轴”，更在于它把矩阵、二次型和换基三件事同时说清楚。下面三个用法后面都会反复出现。

\textbf{用法一：换到主轴坐标系，矩阵变成对角阵。}把三个主轴按列拼成矩阵
\[
\ten R=\begin{pmatrix}\vect u_1&\vect u_2&\vect u_3\end{pmatrix}.
\]
因为三个列矢量正交归一，$\ten R^{\mathsf T}\ten R=\unitd$。由 $\ten T\vect u_k=\lambda_k\vect u_k$，
\[
\ten T\ten R
=\begin{pmatrix}\ten T\vect u_1&\ten T\vect u_2&\ten T\vect u_3\end{pmatrix}
=\begin{pmatrix}\lambda_1\vect u_1&\lambda_2\vect u_2&\lambda_3\vect u_3\end{pmatrix}
=\ten R\operatorname{diag}(\lambda_1,\lambda_2,\lambda_3).
\]
两边左乘 $\ten R^{\mathsf T}$，利用 $\ten R^{\mathsf T}\ten R=\unitd$，得到
\[
\ten R^{\mathsf T}\ten T\ten R
=\operatorname{diag}(\lambda_1,\lambda_2,\lambda_3).
\]
这就是“换到主轴坐标系，账本变成对角阵”的准确意思。注意 $\ten R$ 的列是新基在旧基下的分量；$\ten R^{\mathsf T}\ten T\ten R$ 正是同一线性映射在新基下的矩阵（\S0.4）。

\textbf{正交换基保持对称。}这里 $\ten R$ 的列是一组正交归一基，所以 $\ten R^{\mathsf T}\ten R=\unitd$。若 $\ten T$ 对称，则
\[
(\ten R^{\mathsf T}\ten T\ten R)^{\mathsf T}
=\ten R^{\mathsf T}\ten T^{\mathsf T}\ten R
=\ten R^{\mathsf T}\ten T\ten R,
\]
所以 $\ten R^{\mathsf T}\ten T\ten R$ 仍是对称矩阵。这个计算放在这里，是因为直到定义了 $\ten R$ 才说得清“换基后的矩阵”是什么；它不改变 \S1.4 的证明。

\textbf{用法二：二次型变成三个平方项。}任意矢量可以沿主轴展开：
\[
\vect a=\sum_{k=1}^{3}c_k\vect u_k,
\qquad
c_k=\vect u_k\cdot\vect a .
\]
因为 $\ten T\vect u_k=\lambda_k\vect u_k$，
\[
\ten T\vect a=\sum_{k=1}^{3}\lambda_k c_k\vect u_k,
\]
再与 $\vect a$ 作点积，利用 $\vect u_j\cdot\vect u_k=\delta_{jk}$：
\[
E(\vect a)=\vect a\cdot\ten T\vect a
=\sum_{k=1}^{3}\lambda_k c_k^2
=\lambda_1(\vect u_1\cdot\vect a)^2
+\lambda_2(\vect u_2\cdot\vect a)^2
+\lambda_3(\vect u_3\cdot\vect a)^2 .
\]
这就是“没有交叉项”的一般形式：交叉项消失，只剩三个投影长度的平方，各乘一个主值。

\textbf{用法三：瑞利商给出上下界。}对任意非零矢量 $\vect a$，定义
\[
\rho(\vect a)=\frac{\vect a\cdot\ten T\vect a}{\vect a\cdot\vect a}.
\]
把 $\vect a=\sum_k c_k\vect u_k$ 代入，分母是 $\sum_k c_k^2$，分子是 $\sum_k\lambda_k c_k^2$，所以
\[
\rho(\vect a)
=\frac{\sum_k\lambda_k c_k^2}{\sum_k c_k^2}
\]
是一个以 $c_k^2/\sum_j c_j^2$ 为权重的加权平均。权重非负且和为 $1$，因此
\[
\lambda_{\min}\le\rho(\vect a)\le\lambda_{\max}.
\]
其中 $\lambda_{\min}$ 与 $\lambda_{\max}$ 分别是三个特征值中最小的和最大的。这个不等式在 \S4.2 说明“绕任意方向的转动惯量介于最小与最大主转动惯量之间”时直接可用。

\textbf{半轴公式。}若 $\ten T$ 正定（三个特征值都大于零），$E(\vect a)=1$ 是椭球。沿主轴 $\vect u_k$ 取 $\vect a=t\vect u_k$，其中 $t$ 是带符号的轴向坐标。则
\[
E=\lambda_k t^2=1
\quad\Longrightarrow\quad
t=\pm\frac{1}{\sqrt{\lambda_k}} .
\]
两个解是椭圆在该轴上的两个端点；半轴长是 $|t|=1/\sqrt{\lambda_k}$。若把 $t$ 就取成非负的半轴长，则 $t=1/\sqrt{\lambda_k}$，但这是先选定了“从中心沿正方向量”的那一半，不是方程自动只给出正根。若 $\lambda_k\le0$，方程没有实解，这个公式不能给出实半轴；此时曲面不是椭球，但“主轴方向仍是特征向量”不变。

\subhdr{1.6\quad 三个不变量、重根与验算}
\label{subsec:invariants2}

由 \S0.5，三个系数是
\[
I_1=\tr\ten T=\lambda_1+\lambda_2+\lambda_3,
\qquad
I_3=\det\ten T=\lambda_1\lambda_2\lambda_3,
\]
\[
I_2=\tfrac12\bigl[(\tr\ten T)^2-\tr\ten T^2\bigr]
=\lambda_1\lambda_2+\lambda_2\lambda_3+\lambda_3\lambda_1 .
\]
右边那两个用 $\lambda$ 表示的等号不需要另证：在主轴系里 $\ten T=\operatorname{diag}(\lambda_1,\lambda_2,\lambda_3)$，而 $\tr$ 与 $\det$ 都不随基变（\S0.5），直接代进去就是。

\textbf{用法：做完一道特征值题，用这三个数各查一遍。}以 \S1.3 的 $\ten D$ 为例：$\tr\ten D=13$、$\det\ten D=32$、$\tr\ten D^2=81$，所以 $I_2=\tfrac12(169-81)=44$。而算出的特征值是 $8,4,1$：
\[
8+4+1=13,\qquad 8\cdot4+4\cdot1+1\cdot8=44,\qquad 8\cdot4\cdot1=32,
\]
三样都对上。这个“和、两两和、积”的检查几乎不花时间，却能把大部分符号错误当场抓住。它只是检查，不是证明；谱定理保证“一定有三个实特征值和正交主轴”，不变量只保证“如果你算对了，它们应当满足这些关系”。

\textbf{重根时的几何条件不能丢。}重根只让主轴的选择不唯一，不会让谱定理失效；但如果把“重根”误读成“只有一个特征向量”，就会漏掉整个平面。正确的读法是：先解出特征值，再解 $(\ten T-\lambda\unitd)\vect n=\vect0$，看这个齐次方程有几个独立方向；对实对称矩阵，这个数目正好等于该特征值的重数。选主轴时，在这个特征空间里任取一组正交归一基即可。\S4.3 的圆柱和 \S5.3 的对称陀螺都会用到这一点。

\hdr{§2}{张量与并矢：对象、写法、运算}
\label{sec:tensor}

这一节把“给一个矢量吐一个矢量”的对象正式写成二阶张量，并建立“并矢—指标—矩阵”三套记号的字典。核心只有两条作用规则；其余都是推论。读的时候请一直想着 \S0 的主线：并矢是对象，矩阵是账本。

\subhdr{2.1\quad 阶：从“吃矢量吐数”到二阶张量}
\label{subsec:order}

“张量”这个词让人紧张，但它的定义方式很朴素：\textbf{按它需要几个矢量才能给出一个数，来分阶。}

\begin{center}
\tabtop
\begin{tabular}{@{}p{0.12\linewidth}p{0.40\linewidth}p{0.40\linewidth}@{}}
\thead{阶} & \thead{吃什么、吐什么} & \thead{例子} \\
\multicolumn{3}{@{}l@{}}{\color{petal}\rule{\linewidth}{0.5pt}}\\
0 阶 & 不吃，直接是一个数 & 温度、质量、动能 \\
1 阶 & 吃一个矢量，吐一个数；且对输入线性 & 力在方向 $\vect u$ 上的投影 $\vect F\cdot\vect u$ \\
2 阶 & 吃两个矢量，吐一个数；且对每个输入都线性 & 应力：$\vect u\cdot\ten\sigma\cdot\vect v$；惯量：$\vect u\cdot\ten J\cdot\vect v$ \\
\end{tabular}
\tabbot
\end{center}

表中的 $\ten\sigma$、$\ten J$ 在 \S6、\S4 才正式定义；这里只用它们说明“二阶”是吃两个矢量。\textbf{“对每个输入都线性”叫多重线性}。对二阶就是要
\[
\ten T(\alpha\vect u_1+\beta\vect u_2,\;\vect v)=\alpha\ten T(\vect u_1,\vect v)+\beta\ten T(\vect u_2,\vect v),
\]
对第二个位置同样要求（两个位置各要求一次，合起来叫“双线性”）。

\textbf{为什么二阶张量有两个身份。}把双线性形式记作 $B(\vect u,\vect v)$。本稿固定\emph{第二个}位置：固定 $\vect v$ 后，$\vect u\mapsto B(\vect u,\vect v)$ 是“吃 $\vect u$ 吐数”的线性函数。在通常的三维空间里（有长度和点乘），每个这样的线性函数都能写成 $\vect w\cdot\vect u$，其中 $\vect w$ 唯一：取 $\vect w$ 为这个线性函数在三个基矢量上的值拼成的矢量即可。于是
\[
\vect v\;\longmapsto\;\bigl[\vect u\mapsto B(\vect u,\vect v)\bigr]
\]
是“吃 $\vect v$ 吐矢量 $\vect w$”的线性映射。反过来，给定线性映射 $\ten T$，令
\[
B(\vect u,\vect v)\stackrel{\text{def}}{=}(\ten T\vect v)\cdot\vect u ,
\]
它对两个位置都线性，所以是双线性形式。两者互逆，于是
\[
\boxed{\;\text{吃两个矢量吐一个数}\quad\Longleftrightarrow\quad\text{吃一个矢量吐一个矢量}\;}
\]
这两个说法是同一件事。课件 p9 写的是右边这个（$\vect b=\ten T\vect a$）；\S1 用的是左边这个（$E=\vect a\cdot\ten T\vect a$）。\textbf{看起来是两套东西，其实是同一个对象的两个面。}

\textbf{槽位约定与转置。}上面特意固定\emph{第二}槽，是因为本稿的并矢右作用恰好这样做。若改固定\emph{第一}槽，得到的是同一台机器的转置。用一个非对称例子把两种识别都算一遍：取 $\vect a=\vect e_1$、$\vect b=\vect e_2$，
\[
B(\vect u,\vect v)=(\vect a\otimes\vect b)(\vect u,\vect v)=(\vect a\cdot\vect u)(\vect b\cdot\vect v).
\]
按本稿“固定第二槽”，对固定 $\vect v$，
\[
\vect u\mapsto(\vect a\cdot\vect u)(\vect b\cdot\vect v)=(\vect b\cdot\vect v)(\vect a\cdot\vect u)
=\bigl[\vect a(\vect b\cdot\vect v)\bigr]\cdot\vect u ,
\]
所以对应的机器是 $\vect v\mapsto\vect a(\vect b\cdot\vect v)=(\vect a\otimes\vect b)\cdot\vect v$。取输入 $\vect v=\vect e_1$，它给出 $\vect a(\vect b\cdot\vect e_1)=\vect0$。若改固定第一槽，则对固定 $\vect u$，$\vect v\mapsto(\vect b\cdot\vect v)(\vect a\cdot\vect u)$ 对应的机器是 $\vect u\mapsto\vect b(\vect a\cdot\vect u)=(\vect b\otimes\vect a)\cdot\vect u$；同样输入 $\vect u=\vect e_1$，它给出 $\vect b$。两者相差一个转置（$\vect b\otimes\vect a=(\vect a\otimes\vect b)^{\mathsf T}$）。所以“输入 $\vect e_1$ 输出 $\vect e_2$ 还是 $\vect0$”取决于选了哪一槽；本稿之后一律固定第二槽，和 \S2.2 的右作用规则一致，\S2.3 的左作用则对应转置那一侧。

\subhdr{2.2\quad 张量积（并矢）：定义、作用、分量}
\label{subsec:dyad}

现在造一个具体的双线性对象。给两个矢量 $\vect a,\vect b$，规定
\[
(\vect a\otimes\vect b)(\vect u,\vect v)\stackrel{\text{def}}{=}(\vect a\cdot\vect u)(\vect b\cdot\vect v).
\]
右边是两个数的乘积，所以确实吐出一个数；对 $\vect u$ 线性、对 $\vect v$ 线性，都满足。这个对象记作 $\vect a\otimes\vect b$，叫 $\vect a$ 与 $\vect b$ 的\textbf{张量积}。\textbf{课件 p9 给它起的名字是 Dyad，中文叫并矢。}单个并矢是最基本的积木；一般二阶张量是并矢的和。

它跟“喂一个矢量”的那个身份怎么对上？把第二个位置定死为 $\vect v$（这就是 \S2.1 的槽位约定）：
\[
(\vect a\otimes\vect b)\cdot\vect v\;\stackrel{\text{def}}{=}\;\vect a\,(\vect b\cdot\vect v).
\]
这就是全部的作用规则，一句话：\textbf{先把 $\vect v$ 投影到 $\vect b$ 上，得到一个数，再拿这个数去放大 $\vect a$。}另一边 $\vect v\cdot(\vect a\otimes\vect b)=(\vect v\cdot\vect a)\vect b$ 是另一条规则；它对应 \S2.1 里“固定第一槽”的转置识别，二者一般不同，\S2.3 专门讲。

\textbf{分量就是分量的乘积。}取正交基 $\{\vect e_i\}$，把 $i,j$ 各取 $1,2,3$，得到九个对象 $\vect e_i\otimes\vect e_j$。它们互相独立（不存在一组不全为零的系数使这些基并矢的组合为零），所以任何二阶张量都能唯一展开成
\[
\ten T=\sum_{i=1}^{3}\sum_{j=1}^{3}T_{ij}\,\vect e_i\otimes\vect e_j .
\]
“唯一”可以直接验证：把求和指标改名为 $i,l$，$\ten T=\sum_{i,l}T_{il}\,\vect e_i\otimes\vect e_l$，再作用在基矢量对 $(\vect e_j,\vect e_k)$ 上，只有 $i=j$、$l=k$ 那一项留下，得到 $T_{jk}$。所以九个系数就是九个分量，一个不多一个不少。特别地
\[
(\vect a\otimes\vect b)_{ij}=a_ib_j
\qquad\Longrightarrow\qquad
\vect a\otimes\vect b=
\begin{pmatrix}
a_1b_1&a_1b_2&a_1b_3\\
a_2b_1&a_2b_2&a_2b_3\\
a_3b_1&a_3b_2&a_3b_3
\end{pmatrix}.
\]
矩阵的样子是“\textbf{列乘行}”，也叫外积。取 $\vect a=(1,2,0)$、$\vect b=(0,1,3)$，则
\[
\vect a\otimes\vect b=\begin{pmatrix}0&1&3\\0&2&6\\0&0&0\end{pmatrix}.
\]
\textbf{注意它不对称}：$(1,2)$ 元是 $1$，$(2,1)$ 元是 $0$。这就是“顺序有意义”的直接证据。再取 $\vect x=(1,1,1)$ 核对作用：
\[
(\vect a\otimes\vect b)\cdot\vect x=\vect a(\vect b\cdot\vect x)=\vect a\cdot4=(4,8,0),
\]
而矩阵乘 $\vect x$ 的三行分别是 $0+1+3=4$、$0+2+6=8$、$0$，一致。

\begin{center}
\begin{minipage}{\linewidth}\centering
\begin{tikzpicture}[scale=1.55,>={Stealth[length=1.7mm]}]
  \coordinate (O) at (0,0);
  \coordinate (A) at (1.0,1.2);
  \coordinate (B) at (1.0,0.0);
  \coordinate (X) at (0.4,1.0);
  \coordinate (R) at (0.4,0.48);
  \draw[inkgray,thin] (0,0)--(1.15,0);
  \draw[vv] (O)--(B) node[below]{\scriptsize $\vect b$};
  \draw[vv] (O)--(A) node[above left]{\scriptsize $\vect a$};
  \draw[vv] (O)--(X) node[left]{\scriptsize $\vect x$};
  \draw[inkgray,densely dashed] (X)--(0.4,0);
  \draw[inkgray] (0.4,0)--(0.46,0)--(0.46,0.06)--(0.4,0.06);
  \fill (0.4,0) circle (1pt);
  \node[inkgray,below] at (0.4,0) {\scriptsize $\vect b\cdot\vect x=0.4$};
  \draw[vviolet] (O)--(R) node[right]{\scriptsize $(\vect a\otimes\vect b)\cdot\vect x$};
  \draw[violet,densely dotted] (R)--(0.4,0);
\end{tikzpicture}\\[2pt]
\figcap{\mbox{\textbf{图 3}（自绘）：}并矢的右作用。先把 $\vect x$ 投影到 $\vect b$ 上，得到数 $\vect b\cdot\vect x=0.4$；再沿 $\vect a$ 放大同样的倍数，得到紫色箭头 $0.4\vect a$。图里 $\vect b$ 取单位矢量，所以投影长度就是点乘的值；一般情形只是多一个 $|\vect b|$ 的因子。}
\end{minipage}
\end{center}

\subhdr{2.3\quad 左作用、转置与对称性}
\label{subsec:leftright}

并矢有两套作用规则，方向相反：
\[
\boxed{\;(\vect a\otimes\vect b)\cdot\vect x=\vect a\,(\vect b\cdot\vect x)\;}
\qquad\text{右作用：先投到 }\vect b\text{，再沿 }\vect a\text{ 放大};
\]
\[
\boxed{\;\vect x\cdot(\vect a\otimes\vect b)=(\vect x\cdot\vect a)\,\vect b\;}
\qquad\text{左作用：先投到 }\vect a\text{，再沿 }\vect b\text{ 放大}.
\]
\textbf{两者一般不同。}取 $\vect a=\vect e_1$、$\vect b=\vect e_2$、$\vect x=\vect e_1+2\vect e_2+3\vect e_3$：
\[
(\vect e_1\otimes\vect e_2)\cdot\vect x=\vect e_1(\vect e_2\cdot\vect x)=2\vect e_1,
\qquad
\vect x\cdot(\vect e_1\otimes\vect e_2)=(\vect x\cdot\vect e_1)\vect e_2=\vect e_2 .
\]
一个指向 $\vect e_1$，一个指向 $\vect e_2$，差别一眼可见。

\textbf{什么时候两者相同：当 $\ten D$ 对称时。}对一般二阶张量 $\ten D$，分量为
\[
(\ten D\cdot\vect x)_i=D_{ij}x_j,\qquad
(\vect x\cdot\ten D)_j=x_iD_{ij}=D_{ij}x_i=(D^{\mathsf T}\vect x)_j .
\]
中间那个等号只是把两个数相乘的顺序换了一下；最后一步用了 $(D^{\mathsf T})_{ji}=D_{ij}$，即 $(\ten D^{\mathsf T}\vect x)_j=(D^{\mathsf T})_{ji}x_i=D_{ij}x_i$。所以 $\ten D\cdot\vect x=\vect x\cdot\ten D$ 对\textbf{所有} $\vect x$ 成立，当且仅当 $\ten D=\ten D^{\mathsf T}$（逐分量就是 $D_{ij}=D_{ji}$）。力学里的 $\ten J$、$\ten\sigma$ 都对称，所以“碰巧”怎么写都对；但 \S3.2 的旋转算子 $S$ \textbf{一般不对称}，左右一混立刻出错。

\textbf{转置。}
\[
(\vect a\otimes\vect b)^{\mathsf T}=\vect b\otimes\vect a,
\qquad
(\ten A\cdot\ten B)^{\mathsf T}=\ten B^{\mathsf T}\cdot\ten A^{\mathsf T},
\qquad
(\ten D^{\mathsf T})_{ij}=D_{ji}.
\]
转置的判别式很好用：\textbf{对任意 $\vect x,\vect y$ 有 $(\ten D\cdot\vect x)\cdot\vect y=\vect x\cdot(\ten D^{\mathsf T}\cdot\vect y)$}。\S1.4 证特征向量正交时用的就是这一条。

\textbf{对称分解。}任何并矢都能唯一拆成对称与反对称两半：
\[
\ten D=\underbrace{\tfrac12(\ten D+\ten D^{\mathsf T})}_{\text{对称}}\;+\;
\underbrace{\tfrac12(\ten D-\ten D^{\mathsf T})}_{\text{反对称}} .
\]
力学里这两半各有明确身份：\textbf{对称半}给惯量、应力、应变；\textbf{反对称半}等价于一个矢量，正是叉乘（\S2.6 马上说）。注意这里的“对称分解”是算子层面的；\S1.1 的“二次型只看见对称部分”是二次型层面的。两者结论相近，但不是同一句话：前者说算子可以拆，后者说二次型对反对称半的贡献为零。

\subhdr{2.4\quad 缩并：点乘、双点积、迹}
\label{subsec:contraction}

两个并矢之间可以“乘”，乘掉几个指标就叫几点积：
\[
(\vect a\otimes\vect b)\cdot(\vect c\otimes\vect d)=(\vect b\cdot\vect c)\,\vect a\otimes\vect d
\qquad\text{（缩并掉一个指标，剩二阶）},
\]
\[
(\vect a\otimes\vect b):(\vect c\otimes\vect d)=(\vect a\cdot\vect c)(\vect b\cdot\vect d)
\qquad\text{（缩并掉两个指标，剩标量）}.
\]
第二条叫\textbf{双点积}，写成指标就是 $\ten A:\ten B=\sum_{i,j}A_{ij}B_{ij}$，与矩阵写法 $\tr(\ten A^{\mathsf T}\ten B)$ 等价：因为 $(\ten A^{\mathsf T}\ten B)_{jj}=A_{ij}B_{ij}$，再对 $j$ 求和。

\textbf{迹。}$\tr\ten D=\sum_iD_{ii}$。用基并矢写，$\tr(\vect a\otimes\vect b)=\vect a\cdot\vect b$；于是
\[
\unitd:\ten D=\tr\ten D,\qquad
\ten A:(\vect u\otimes\vect v)=\vect u\cdot\ten A\cdot\vect v .
\]
\textbf{数字例子。}取 $\ten D=\vect e_1\otimes\vect e_1+2\vect e_1\otimes\vect e_2+3\vect e_2\otimes\vect e_2$，矩阵是 $\begin{pmatrix}1&2\\0&3\end{pmatrix}$（二维示意，三维同理）。则
\[
\tr\ten D=1+3=4,\qquad
\ten D:\ten D=1^2+2^2+3^2=14,\qquad
\unitd:\ten D=4 .
\]
最后一条可以当公式用：$\unitd$ 与谁双点积就取谁的迹。

\subhdr{2.5\quad 单位并矢与投影并矢}
\label{subsec:projection}

\textbf{单位并矢}就是克罗内克符号那个张量：
\[
\unitd=\sum_{k=1}^{3}\vect e_k\otimes\vect e_k
\quad\Longrightarrow\quad
\unitd\cdot\vect x=\sum_k\vect e_k(\vect e_k\cdot\vect x)=\vect x
\qquad\text{（对任意正交基成立）}.
\]
最后一个等号就是“把 $\vect x$ 在三个正交方向上的投影拼回来”。$\unitd$ 换基后还是 $\unitd$，\S2.7 会用变换律再证一次。

\textbf{投影并矢。}设 $\vect n$ 是单位矢量，令 $\ten P=\vect n\otimes\vect n$，则
\[
\ten P\cdot\vect x=\vect n(\vect n\cdot\vect x)
\quad\text{（$\vect x$ 在 }\vect n\text{ 上的投影）},
\qquad
\ten P^2=\ten P\ \text{（再投一次不变）}.
\]
而 $\unitd-\vect n\otimes\vect n$ 是“投影到垂直于 $\vect n$ 的平面”的算子。两者的迹分别是 $1$ 与 $2$：$\ten P$ 的迹为 $1$，对应它\textbf{保留}沿 $\vect n$ 的 $1$ 个方向、去掉垂直平面里的 $2$ 个；$\unitd-\ten P$ 的迹为 $2$，对应它\textbf{保留}垂直平面里的 $2$ 个方向、去掉沿 $\vect n$ 的 $1$ 个。在三维里，一个投影算子的迹正好等于它保留的方向数（把“投影掉几个自由度”说反，就会把保留维数和消去维数颠倒）。

\begin{center}
\begin{tikzpicture}[scale=1.8,>={Stealth[length=1.7mm]}]
  \coordinate (O) at (0,0);
  \coordinate (N) at (0.866,0.5);
  \coordinate (X) at (0.2,1.1);
  \coordinate (P) at (0.6263,0.3616);
  \draw[inkgray,densely dashed] (-0.1,-0.058)--(1.05,0.606);
  \draw[vv] (O)--(N) node[above right]{\scriptsize $\vect n$};
  \draw[vv] (O)--(X) node[left]{\scriptsize $\vect x$};
  \draw[vviolet] (O)--(P) node[below right]{\scriptsize $\ten P\cdot\vect x$};
  \draw[inkgray,thick,->,>={Stealth[length=1.7mm]}] (P)--(X) node[right]{\scriptsize $(\unitd-\ten P)\cdot\vect x$};
  \draw[inkgray,thin] (0.6783,0.3916)--(0.6483,0.4436)--(0.5963,0.4136);
  \fill (O) circle (0.7pt);
\end{tikzpicture}\\[2pt]
\figcap{\mbox{\textbf{图 4}（自绘）：}投影并矢 $\ten P=\vect n\otimes\vect n$ 把 $\vect x$ 压到 $\vect n$ 所在的直线上，得到紫箭头 $\ten P\cdot\vect x=\vect n(\vect n\cdot\vect x)$；剩下的灰箭头是 $(\unitd-\ten P)\cdot\vect x=\vect x-\vect n(\vect n\cdot\vect x)$，它落在与 $\vect n$ 垂直的平面里（图中直角标记）。两个分量之和就是 $\vect x$ 本身——这就是 $\ten P+(\unitd-\ten P)=\unitd$ 的几何意思。\textbf{这里的 $\unitd$ 是单位并矢（二阶张量），不是数 $1$}。}
\end{center}

\subhdr{2.6\quad 对偶张量与叉乘}
\label{subsec:dual}

\S0.2 已经给出叉乘的分量式。现在把它写成一个二阶张量在作用。对任意矢量 $\vect a$，定义 $\vect a^{\times}$ 的分量为
\[
(\vect a^{\times})_{ij}=\varepsilon_{ikj}a_k,
\qquad
\vect a^{\times}=
\begin{pmatrix}
0&-a_3&a_2\\
a_3&0&-a_1\\
-a_2&a_1&0
\end{pmatrix}.
\]
它的作用正好是叉乘：
\[
(\vect a^{\times}\cdot\vect x)_i=\varepsilon_{ikj}a_kx_j=\varepsilon_{ijk}a_jx_k=(\vect a\times\vect x)_i,
\]
即 $\vect a^{\times}\cdot\vect x=\vect a\times\vect x$。反过来，若 $\ten A$ 反对称，则存在唯一的矢量 $\vect a$ 使 $\ten A=\vect a^{\times}$。直接由矩阵元素读出
\[
a_1=A_{32}=-A_{23},\qquad a_2=A_{13}=-A_{31},\qquad a_3=A_{21}=-A_{12},
\]
再代回 $\vect a^{\times}$ 就逐项还原 $\ten A$；所以这个对应是一一的。等价的反解式是 $a_k=-\frac12\varepsilon_{kij}A_{ij}$。这就是课件 p8 的对偶张量。\textbf{一句话：叉乘不是“另一种乘法”，而是一个反对称并矢在作用。}“为什么叉乘不满足结合律”也就通了：不满足结合律的是叉乘这个运算；而反对称张量的\textbf{合成}（矩阵乘法）是满足结合律的。

\textbf{符号约定要看清楚。}课件 p8 写的是 $a^{\times}_{ik}=\varepsilon_{ijk}a_j$，与本稿 $(\vect a^{\times})_{ij}=\varepsilon_{ikj}a_k$ 是同一个矩阵（把哑指标改个名字即可）；RIT Whelan notes 10 的 p7 用的是 $\ast A\cdot\vect B=-\vect A\times\vect B$，即它的 $\ast A$ 等于 $-\vect A^{\times}$。读不同书时先认矩阵的三个元素，别只记一个公式。

\subhdr{2.7\quad 张量变换律与商法则}
\label{subsec:transformation}

随便给我九个排成方阵的数，它是不是张量？\textbf{不一定。}张量的定义不在“有九个分量”，而在“换基时这九个分量怎么变”。规则就是 \S0.4 那条：
\[
\boxed{\;T'_{ij}=R_{ki}R_{lj}T_{kl}\;}
\qquad\text{（$R_{ki}=\vect n_i\cdot\vect e_k$ 是新基在旧基下的分量）}
\]
凡按这条规则变的九元数组，才是二阶张量；不按这条变的，只是九个数字。注意这条规则和 \S0.4 的 $\ten T'=\ten R^{\mathsf T}\ten T\ten R$ 是同一件事，只是写成了指标。

实践中更常用的是反过来：手上有一个物理关系，想证明里面那个九元数组是张量。这就是\textbf{商法则}：

\begin{center}
\fcolorbox{violet}{bandviolet}{\parbox{0.94\linewidth}{\vspace{3pt}
\textbf{商法则（三维，实用版）。}若对\textbf{任意}矢量 $\vect v$ 都有 $\vect u=\ten T\vect v$（即 $u_i=T_{ij}v_j$），且已知 $\vect u$、$\vect v$ 都是矢量（换基时按 $\vect u'=\ten R^{\mathsf T}\vect u$、$\vect v'=\ten R^{\mathsf T}\vect v$ 变），则 $\ten T$ 必是二阶张量，即 $\ten T'=\ten R^{\mathsf T}\ten T\ten R$。
\vspace{3pt}}}
\end{center}

证明只有三步。在带撇的坐标系里同一条定律必须同样成立，$\vect u'=\ten T'\vect v'$；把 $\vect u'=\ten R^{\mathsf T}\vect u$、$\vect v'=\ten R^{\mathsf T}\vect v$ 代进去：
\[
\ten R^{\mathsf T}\vect u=\ten T'\ten R^{\mathsf T}\vect v
\quad\Longrightarrow\quad
\vect u=\ten R\ten T'\ten R^{\mathsf T}\vect v .
\]
与原来的 $\vect u=\ten T\vect v$ 对任意 $\vect v$ 比较，得 $\ten T=\ten R\ten T'\ten R^{\mathsf T}$，两边左乘 $\ten R^{\mathsf T}$、右乘 $\ten R$（用 $\ten R^{\mathsf T}\ten R=\unitd$）就得到 $\ten T'=\ten R^{\mathsf T}\ten T\ten R$。\textbf{所以“它是张量”不是假设，是被这条定律逼出来的。}\S6.1 会用它：面力 $\vect t$ 与法向 $\vect n$ 都是矢量，而 $\vect t(\vect n)=\ten\sigma\cdot\vect n$ 对任意 $\vect n$ 成立，于是 $\ten\sigma$ 一定是二阶张量。

\textbf{一个真的张量：$\delta_{ij}$。}把 $T_{kl}=\delta_{kl}$ 代进变换律：
\[
\delta'_{ij}=R_{ki}R_{lj}\delta_{kl}=R_{ki}R_{kj}=(\ten R^{\mathsf T}\ten R)_{ij}=\delta_{ij}.
\]
所以单位张量在任何正交基下的分量都是 $\delta_{ij}$——这正是“单位矩阵换基后还是单位矩阵”的深层原因，也是 \S2.5 的 $\unitd=\sum_k\vect e_k\otimes\vect e_k$ 对任意正交基都成立的理由。

\textbf{一个不是张量的东西：位置矢量。}三个数 $(x_1,x_2,x_3)$ 在\textbf{旋转}下确实按矢量规则变（所以转动问题里可以当矢量用），但一旦允许\textbf{平移}坐标系，它会多出一个常数项，不再是齐次的（“齐次”就是分量只含坐标的一次项、不含常数项）。所以“是不是张量”这句话必须连着\textbf{允许哪些变换}说。本稿（以及整门理论力学）只允许旋转与反射，即\textbf{正交变换}，在这个群里位置矢量是矢量。

\subhdr{2.8\quad 三记号字典与缩并表（查阅）}
\label{subsec:dictionary}

下表把同一个对象在三种写法下并排列出。\textbf{它不是要你背，而是你读到任何一本书时可以横向对照。}

\begin{center}
\begin{minipage}{\linewidth}
\small
\renewcommand{\arraystretch}{1.12}
\tabtop
\begin{tabular}{@{}p{0.28\linewidth}p{0.20\linewidth}p{0.36\linewidth}@{}}
\thead{Gibbs 并矢} & \thead{指标} & \thead{矩阵} \\
\multicolumn{3}{@{}l@{}}{\color{petal}\rule{\linewidth}{0.5pt}}\\
$\vect a\vect b$ 或 $\vect a\otimes\vect b$ & $a_ib_j$ & $\vect a\,\vect b^{\mathsf T}$（列乘行） \\
$\ten D\cdot\vect x$ & $D_{ij}x_j$ & $\ten D\,\vect x$ \\
$\vect x\cdot\ten D$ & $x_iD_{ij}$ & $\vect x^{\mathsf T}\ten D$ \\
$\ten A\cdot\ten B$ & $A_{ik}B_{kj}$ & $\ten A\ten B$ \\
$\ten A:\ten B$ & $A_{ij}B_{ij}$ & $\tr(\ten A^{\mathsf T}\ten B)$ \\
$\ten D^{\mathsf T}$ & $D_{ji}$ & $\ten D^{\mathsf T}$ \\
$\unitd$ & $\delta_{ij}$ & 单位矩阵 \\
$\tr\ten D$ & $D_{ii}$ & 对角元之和 \\
$\vect a^{\times}$ & $\varepsilon_{ikj}a_k$ & 反对称矩阵（\S2.6） \\
$\ten D=\sum_k\lambda_k\vect n_k\vect n_k^{\mathsf T}$ & $\sum_k\lambda_k n_{ki}n_{kj}$ & $\ten R\operatorname{diag}(\lambda)\ten R^{\mathsf T}$ \\
\end{tabular}
\tabbot
\end{minipage}
\end{center}

\textbf{表的最后一行值得单独看。}它就是“对角化”的全部内容：找一组正交基，使矩阵变成对角阵。矩阵写法要写三个东西（$\ten R$、对角线、$\ten R^{\mathsf T}$），并矢写法只写一项求和，而物理含义在并矢写法里是\textbf{看得见}的：三个投影并矢各自乘上主值。这就是为什么力学书偏爱并矢。

\textbf{但并矢不是矩阵。}矩阵是“坐标下的账本”，并矢是不依赖坐标系的对象。换基时账本变、对象不变（\S0.4 的主线）。你可以在任意基下把并矢写成九个分量，但并矢本身不是那九个数——\S2.7 已经说明，只有按变换律变的九个数才配叫张量。

\begin{center}
\tabtop
\begin{tabular}{@{}lll@{}}
\thead{写法} & \thead{结果} & \thead{分量} \\
\multicolumn{3}{@{}l@{}}{\color{petal}\rule{\linewidth}{0.5pt}}\\
$\ten A\cdot\ten B$ & 二阶张量 & $A_{ik}B_{kj}$ \\
$\ten A\cdot\vect x$ & 矢量 & $A_{ij}x_j$ \\
$\vect x\cdot\ten A$ & 矢量 & $x_iA_{ij}$ \\
$\ten A:\ten B$ & 标量 & $A_{ij}B_{ij}$ \\
$\tr\ten A$ & 标量 & $A_{ii}$ \\
$(\vect a\otimes\vect b)\cdot(\vect c\otimes\vect d)$ & 二阶张量 & $(\vect b\cdot\vect c)\,\vect a\otimes\vect d$ \\
$\vect a\cdot(\vect b\otimes\vect c)$ & 矢量 & $(\vect a\cdot\vect b)\vect c$ \\
$(\vect a\otimes\vect b):(\vect c\otimes\vect d)$ & 标量 & $(\vect a\cdot\vect c)(\vect b\cdot\vect d)$ \\
\end{tabular}
\tabbot
\end{center}

逐行读法：前五行的规律是“每点一下吃掉一个指标”，所以二阶$\cdot$二阶留二阶、二阶$\cdot$矢量留一阶、双点积吃掉两个留零阶。最后三行要小心，它们是整表最容易错的地方：
\begin{itemize}
\item 第六行：两个并矢点乘，缩并的是\textbf{右并矢的左因子}与\textbf{左并矢的右因子}，即 $\vect b$ 与 $\vect c$。这一条是 \S3.3 证明 $S^{\mathsf T}S=\unitd$ 的关键一步。
\item 第七行：矢量点并矢，$\vect a\cdot(\vect b\otimes\vect c)=(\vect a\cdot\vect b)\vect c$，结果是矢量，方向沿 $\vect c$。它和第六行差在“谁先被缩并”，不要凭形状猜。
\item 第八行：双点积把两个并矢的四个因子两两配对，结果是一个数；它等于矩阵写法 $\tr(\ten A^{\mathsf T}\ten B)$。
\end{itemize}

\warnbox{\textbf{选读（可跳过）：上下标是从哪来的。}
你以后会看到 $T^{i}{}_{j}$、$T_{ij}$、$T^{ij}$ 三种写法并存。在\textbf{正交}基里，三种写法数值完全相同，所以随便挑一种。但在\textbf{非正交}基（斜坐标、曲线坐标）里，矢量有两套分量：把矢量投影到基矢量上得到一套（协变，下标记），把矢量沿基矢量方向分解的系数得到另一套（逆变，上标记）；二者由一个叫度规的 $g_{ij}$ 互相转换，$g_{ij}$ 就是基矢量的点乘表 $g_{ij}=\vect e_i\cdot\vect e_j$。只有当 $g_{ij}=\delta_{ij}$（正交归一）时，两套才重合，上下标才可以不分。
\textbf{结论：}理论力学的刚体部分全程正交基，所以本稿一律用下标，不会错。等到连续介质力学用曲线坐标、或相对论用四维时空时，上下标就必须写清楚——那时再回来补这半页即可。跳过后从 \S3 继续。}

\hdr{§3}{旋转算子与角速度：课件 p10--p12 的并矢读法}
\label{sec:rotation}

这一节把课件的旋转算子 $S$ 讲成“把旧基送到新基的那台机器”，再从它的导数得到角速度。它同时是并矢运算的一次综合练习：张量积、单位并矢、投影并矢、对偶张量都会出现。这里的\textbf{刚体}指任意两点距离不随时间变的物体模型；本稿只用它的运动学含义，不讨论内力与变形。

\subhdr{3.1\quad 两个基与一个算子}
\label{subsec:bases}

取固定系的正交基 $\vect e_1,\vect e_2,\vect e_3$（课件 p10 记作 $\vect n_1,\vect n_2,\vect n_3$），再取随刚体转动的正交基 $\vect n'_1(t),\vect n'_2(t),\vect n'_3(t)$，并约定
\[
\vect n'_k(0)=\vect e_k,\qquad k=1,2,3 .
\]
刚体上任意一点的位置矢量可以写成
\[
\vect r(t)=x'_1\vect n'_1(t)+x'_2\vect n'_2(t)+x'_3\vect n'_3(t)=x'_k\vect n'_k(t),
\]
其中 $x'_k$ 是\textbf{随体坐标}，对刚体上的点不随时间变（刚体上任意两点的距离不变，所以随体坐标不变）。在 $t=0$ 时 $\vect r(0)=x'_k\vect e_k$。

我们要找一台“机器” $S(t)$，它把 $t=0$ 的位置送到 $t$ 时刻的位置：
\[
\vect r(t)=S(t)\vect r(0),\qquad\text{并且}\quad \vect n'_k(t)=S(t)\vect e_k .
\]
第二句是第一句的“基矢量版本”：如果 $S$ 把每个旧基矢量都送到对应的新基矢量，那么由线性，它就把任意矢量送到正确的位置。

\subhdr{3.2\quad $S(t)=\sum_k\vect n'_k\otimes\vect e_k$}
\label{subsec:rotationS}

课件 p10 直接给出
\[
S(t)=\sum_{k=1}^{3}\vect n'_k(t)\otimes\vect e_k .
\]
为什么是它？把它作用在旧基矢量 $\vect e_j$ 上：
\[
S(t)\vect e_j=\sum_k\vect n'_k(\vect e_k\cdot\vect e_j)=\vect n'_j .
\]
正好把 $\vect e_j$ 送到 $\vect n'_j$。再作用在 $\vect r(0)=x'_k\vect e_k$ 上：
\[
S(t)\vect r(0)=x'_kS(t)\vect e_k=x'_k\vect n'_k(t)=\vect r(t).
\]
所以 $S$ 的定义不是猜的，是“要求 $S\vect e_k=\vect n'_k$”唯一确定下来的；写成并矢和只是把这个要求一次性满足。

\textbf{矩阵读法。}$S$ 在固定基下的矩阵，第 $k$ 列就是 $\vect n'_k$ 的三个分量（因为 $S\vect e_k=\vect n'_k$）。二维绕 $\vect e_3$ 转 $\varphi$ 的例子：
\[
\vect n'_1=\cos\varphi\,\vect e_1+\sin\varphi\,\vect e_2,\qquad
\vect n'_2=-\sin\varphi\,\vect e_1+\cos\varphi\,\vect e_2,
\]
\[
S(\varphi;\vect e_3)=
\begin{pmatrix}
\cos\varphi&-\sin\varphi\\
\sin\varphi&\cos\varphi
\end{pmatrix},
\qquad
S\vect e_1=\vect n'_1,\quad S\vect e_2=\vect n'_2 .
\]
这正是 \S0.4 里那个 $\ten R$（列是新基分量的矩阵）。所以 \S0.4 的换基矩阵和这里的旋转算子是同一个东西；只是那里用它换账本，这里用它转矢量。下面统一用 $S$ 表示“作用在矢量上的旋转算子”，用 $\ten R$ 表示“换基矩阵”；当新基由 $S$ 给出时，两者矩阵相同。

\begin{center}
\begin{tikzpicture}[scale=1.75,>={Stealth[length=1.7mm]}]
  \coordinate (O) at (0,0);
  \coordinate (E1) at (1,0);
  \coordinate (E2) at (0,1);
  \coordinate (N1) at (0.8192,0.5736);
  \coordinate (N2) at (-0.5736,0.8192);
  \draw[ax] (O)--(1.25,0) node[right]{\scriptsize $\vect e_1$};
  \draw[ax] (O)--(0,1.25) node[above]{\scriptsize $\vect e_2$};
  \draw[vviolet] (O)--(N1) node[right]{\scriptsize $\vect n'_1$};
  \draw[vviolet] (O)--(N2) node[above left]{\scriptsize $\vect n'_2$};
  \draw[inkgray,thin] (0.48,0) arc (0:35:0.48);
  \node[inkgray] at (0.60,0.16) {\scriptsize $\varphi$};
  \draw[inkgray,thin] (0,0) circle (0.075);
  \fill (0,0) circle (0.025);
  \node[inkgray,below right] at (0.04,-0.02) {\scriptsize $\vect n$（出纸面）};
\end{tikzpicture}\\[2pt]
\figcap{\mbox{\textbf{图 5}（据课件 p10 的定义自绘）：}在垂直于转轴 $\vect n$ 的平面里，旧基 $\vect e_1,\vect e_2$（灰虚线）被 $S$ 送到新基 $\vect n'_1,\vect n'_2$（紫线），转角 $\varphi$。转轴 $\vect n$ 指向纸外（$\odot$）。$S$ 的矩阵第 $1$ 列是 $\vect n'_1$ 的分量，第 $2$ 列是 $\vect n'_2$ 的分量；$S\vect e_k=\vect n'_k$ 是这张图的公式版本。}
\end{center}

\subhdr{3.3\quad $S$ 是正交算子：长度与点乘都不变}
\label{subsec:orthogonal}

用 \S2.8 的并矢点乘规则算 $S^{\mathsf T}S$。先写
\[
S=\sum_k\vect n'_k\otimes\vect e_k,\qquad
S^{\mathsf T}=\sum_l\vect e_l\otimes\vect n'_l .
\]
于是
\[
S^{\mathsf T}S=\sum_{k,l}(\vect e_l\otimes\vect n'_l)\cdot(\vect n'_k\otimes\vect e_k)
=\sum_{k,l}(\vect n'_l\cdot\vect n'_k)\,\vect e_l\otimes\vect e_k
=\sum_k\vect e_k\otimes\vect e_k=\unitd .
\]
中间那步用了 $\vect n'_l\cdot\vect n'_k=\delta_{lk}$（新基也是正交归一的）。同法
\[
SS^{\mathsf T}=\sum_{k,l}(\vect n'_k\otimes\vect e_k)\cdot(\vect e_l\otimes\vect n'_l)
=\sum_{k,l}(\vect e_k\cdot\vect e_l)\,\vect n'_k\otimes\vect n'_l
=\sum_k\vect n'_k\otimes\vect n'_k=\unitd .
\]
所以 $S^{-1}=S^{\mathsf T}$。它保持点乘：
\[
(S\vect x)\cdot(S\vect y)=\vect x\cdot(S^{\mathsf T}S\vect y)=\vect x\cdot\vect y .
\]
长度、夹角都不变，这正是“刚体转动”该有的性质。又因为 $\det(S^{\mathsf T}S)=(\det S)^2=1$，所以 $\det S=\pm1$；$S(t)$ 随 $t$ 连续变化且 $S(0)=\unitd$ 的 $\det$ 为 $+1$，而 $\det$ 是矩阵元素的连续函数，不能从 $+1$ 跳到 $-1$ 而不经过 $0$，所以\textbf{连续转动给 $\det S=+1$}（$\det=-1$ 对应镜像反射，比如把右手变成左手，不是刚体转动）。

\subhdr{3.4\quad 绕轴转 $\varphi$ 的显式公式}
\label{subsec:axisformula}

课件 p11 给出
\[
S(\varphi;\vect n)=\vect n\otimes\vect n+(\unitd-\vect n\otimes\vect n)\cos\varphi+\vect n^{\times}\sin\varphi .
\]
逐项验证。把任意矢量 $\vect x$ 拆成“沿轴”与“垂直轴”两部分：
\[
\vect x=\vect x_\parallel+\vect x_\perp,\qquad
\vect x_\parallel=(\vect n\cdot\vect x)\vect n=(\vect n\otimes\vect n)\vect x,\qquad
\vect x_\perp=\vect x-\vect x_\parallel .
\]
于是
\[
(\vect n\otimes\vect n)\vect x=\vect x_\parallel,\qquad
(\unitd-\vect n\otimes\vect n)\vect x=\vect x_\perp,\qquad
\vect n^{\times}\vect x=\vect n\times\vect x=\vect n\times\vect x_\perp
\]
（最后一步因为 $\vect n\times\vect x_\parallel=0$）。代入公式：
\[
S\vect x=\vect x_\parallel+\vect x_\perp\cos\varphi+(\vect n\times\vect x_\perp)\sin\varphi .
\]
读法：\textbf{沿轴分量不动；垂直轴分量在垂直于 $\vect n$ 的平面里转了 $\varphi$（$\cos$ 项是原方向，$\sin$ 项是转 $90^\circ$ 后的方向）。}这就是右手定则的转动。

\textbf{两个核对。}取 $\vect n=\vect e_3$、$\varphi=90^\circ$、$\vect x=\vect e_1$：
\[
S\vect e_1=(\vect e_3\otimes\vect e_3)\vect e_1+0\cdot(\unitd-\vect e_3\otimes\vect e_3)\vect e_1+\vect e_3^{\times}\vect e_1=\vect e_2,
\]
与“绕 $z$ 轴把 $x$ 转到 $y$”一致。取 $\varphi=0$：$S=\vect n\otimes\vect n+(\unitd-\vect n\otimes\vect n)=\unitd$，也对。

\textbf{注意 $S$ 一般不对称。}$\vect n\otimes\vect n$ 与 $\unitd-\vect n\otimes\vect n$ 都对称，但 $\vect n^{\times}$ 反对称，所以 $S$ 的转置是
\[
S^{\mathsf T}(\varphi;\vect n)=\vect n\otimes\vect n+(\unitd-\vect n\otimes\vect n)\cos\varphi-\vect n^{\times}\sin\varphi=S(-\varphi;\vect n).
\]
因此 $S\cdot\vect x$ 与 $\vect x\cdot S$ 一般不同（\S2.3 的警告在这里兑现）：前者是“把矢量转 $+\varphi$”，后者是“把矢量转 $-\varphi$”。

\subhdr{3.5\quad 无穷小旋转与角速度}
\label{subsec:angularvelocity}

把 $\varphi$ 取成无穷小 $\dd\varphi$，用 $\cos\dd\varphi\approx1$、$\sin\dd\varphi\approx\dd\varphi$（只保留 $\dd\varphi$ 的一次项，忽略 $\dd\varphi^2$ 及更高次）：
\[
S(\dd\varphi;\vect n)=\vect n\otimes\vect n+(\unitd-\vect n\otimes\vect n)+\vect n^{\times}\dd\varphi
=\unitd+\dd\varphi\,\vect n^{\times}=\unitd+\dd{\boldsymbol\varphi}^{\times},
\]
其中 $\dd{\boldsymbol\varphi}=\dd\varphi\,\vect n$ 是“转轴乘转角”的矢量。两个无穷小旋转相乘：
\[
S_1S_2=(\unitd+\dd{\boldsymbol\varphi}_1^{\times})(\unitd+\dd{\boldsymbol\varphi}_2^{\times})
=\unitd+\dd{\boldsymbol\varphi}_1^{\times}+\dd{\boldsymbol\varphi}_2^{\times}+O(\dd\varphi^2),
\]
一阶项对易（相乘顺序可以交换），所以\textbf{无穷小转动可交换}；有限转动一般不可交换（本稿不用这个结论）。

现在对 $S(t)S^{\mathsf T}(t)=\unitd$ 求导：
\[
\frac{\dd S}{\dd t}S^{\mathsf T}+S\frac{\dd S^{\mathsf T}}{\dd t}=0 .
\]
记
\[
\ten A=\frac{\dd S}{\dd t}S^{\mathsf T}.
\]
则
\[
\ten A^{\mathsf T}=\Bigl(\frac{\dd S}{\dd t}S^{\mathsf T}\Bigr)^{\mathsf T}
=S\frac{\dd S^{\mathsf T}}{\dd t}=-\ten A,
\]
所以 $\ten A$ 反对称。由 \S2.6，反对称二阶张量恰好对应一个矢量，记作 $\boldsymbol\Omega$，使
\[
\boxed{\;\boldsymbol\Omega^{\times}=\frac{\dd S}{\dd t}S^{\mathsf T}\;}
\]
这就是课件 p12 的角速度定义。$\boldsymbol\Omega$ 叫\textbf{角速度矢量}；它不是“九个分量”，而是反对称张量的三个独立分量（\S2.6），所以是一个矢量。

\textbf{为什么 $\dd\vect n'_k/\dd t=\boldsymbol\Omega\times\vect n'_k$。}由 $\ten A=\boldsymbol\Omega^{\times}$ 与 $S^{\mathsf T}S=\unitd$，有
\[
\frac{\dd S}{\dd t}=\boldsymbol\Omega^{\times}S .
\]
两边作用在 $\vect e_k$ 上，并用 $S\vect e_k=\vect n'_k$：
\[
\frac{\dd\vect n'_k}{\dd t}=\frac{\dd}{\dd t}(S\vect e_k)=\frac{\dd S}{\dd t}\vect e_k=\boldsymbol\Omega^{\times}S\vect e_k=\boldsymbol\Omega\times\vect n'_k .
\]
这正是“随体基矢量以角速度 $\boldsymbol\Omega$ 转动”的微分版本。作为核对，也可以直接用并矢算：
\[
\frac{\dd S}{\dd t}S^{\mathsf T}
=\sum_k\frac{\dd\vect n'_k}{\dd t}\otimes\vect n'_k,
\qquad
\boldsymbol\Omega^{\times}=\sum_k(\boldsymbol\Omega\times\vect n'_k)\otimes\vect n'_k,
\]
两式逐项相同，所以 $\dd\vect n'_k/\dd t=\boldsymbol\Omega\times\vect n'_k$（第二个等号用了 $\sum_k\vect n'_k\otimes\vect n'_k=\unitd$）。

\textbf{无穷小情形对一下。}取 $S=\unitd+\dd{\boldsymbol\varphi}^{\times}$，则 $\dd S/\dd t=\dot{\boldsymbol\varphi}^{\times}$、$S^{\mathsf T}\approx\unitd$，所以
\[
\boldsymbol\Omega^{\times}=\dot{\boldsymbol\varphi}^{\times}
\quad\Longrightarrow\quad
\boldsymbol\Omega=\frac{\dd{\boldsymbol\varphi}}{\dd t}.
\]
角速度就是“转角矢量对时间的变化率”；方向是瞬时转轴，大小是转动的快慢。\S5 的输运定理与欧拉方程都要用这个矢量。

\subhdr{3.6\quad 刚体速度关系：课件 p15 的已用前提}
\label{subsec:rigidvelocity}

课件 p15 总结：刚体上任意一点的速度可以写成
\[
\vect v=\vect V+\boldsymbol\Omega\times\vect r,
\]
其中 $\vect V$ 是所选动系原点的速度（平动速度），$\boldsymbol\Omega$ 是角速度，$\vect r$ 是从动系原点到该点的位置矢量。本稿从 \S4 起把它当作已用前提；这里用旋转算子解释它为什么成立，以免把它当成凭空公式。

设动系原点 $O'$ 相对固定系原点 $O$ 的位置是 $\vect R$，刚体上一点相对 $O'$ 的位置是 $\vect r$。若刚体只绕 $O'$ 转动，则 $\vect r(t)=S(t)\vect r(0)$，由 \S3.5 的 $\dd S/\dd t=\boldsymbol\Omega^{\times}S$，
\[
\frac{\dd\vect r}{\dd t}=\frac{\dd S}{\dd t}\vect r(0)=\boldsymbol\Omega^{\times}S\vect r(0)=\boldsymbol\Omega\times\vect r .
\]
再加上动系原点自己的平动速度 $\vect V=\dd\vect R/\dd t$，就得到 $\vect v=\vect V+\boldsymbol\Omega\times\vect r$。这不是另加的实验规律，而是“刚体上任意两点距离不变”加“随体基正交”的运动学结果。

\textbf{换基点时 $\boldsymbol\Omega$ 不变。}若把动系原点从 $O'$ 挪到 $O''$，位移记作 $\vect a$，同一刚体点的位置从 $\vect r$ 变成 $\vect r'=\vect r-\vect a$。代入速度式：
\[
\vect v=\vect V+\boldsymbol\Omega\times(\vect a+\vect r')
=(\vect V+\boldsymbol\Omega\times\vect a)+\boldsymbol\Omega\times\vect r' .
\]
所以新的平动速度是 $\vect V'=\vect V+\boldsymbol\Omega\times\vect a$，而角速度仍是同一个 $\boldsymbol\Omega$。这就是课件 p15 的第二条总结；它说明“角速度”不是某个基点私有的量。

\hdr{§4}{落点一：转动惯量张量}
\label{sec:inertia}

到这里工具已经齐了：线性映射与换基（\S0）、二次型与主轴（\S1）、并矢运算（\S2）、旋转算子与角速度（\S3）。从这一节起，每个结论都只用前面出现过的东西。

\subhdr{4.1\quad 从 $\vect L=\int\vect r\times\vect v\dd m$ 推出惯量张量}
\label{subsec:inertia-construct}

取刚体绕定点 $O$ 转动（或取质心为基点），角速度为 $\boldsymbol\omega$。由 \S3.6 的刚体速度关系，取动系原点在 $O$ 且 $\vect V=\vect0$，刚体上位置为 $\vect r$ 的点速度为
\[
\vect v=\boldsymbol\omega\times\vect r .
\]
角动量（对 $O$ 点）定义为
\[
\vect L=\int\vect r\times\vect v\dd m
=\int\vect r\times(\boldsymbol\omega\times\vect r)\dd m .
\]
用三矢量叉乘公式
\[
\vect a\times(\vect b\times\vect c)=\vect b(\vect a\cdot\vect c)-\vect c(\vect a\cdot\vect b)
\]
展开内层：
\[
\vect r\times(\boldsymbol\omega\times\vect r)
=\boldsymbol\omega(\vect r\cdot\vect r)-\vect r(\vect r\cdot\boldsymbol\omega)
=r^2\boldsymbol\omega-\vect r(\vect r\cdot\boldsymbol\omega).
\]
这个叉乘恒等式是标准代数恒等式，把两边按分量展开即可核对；它的几何读法是“先减去沿 $\boldsymbol\omega$ 的分量，再补上垂直方向的分量”。现在把两项都写成“某个张量作用在 $\boldsymbol\omega$ 上”：
\[
r^2\boldsymbol\omega=r^2\unitd\cdot\boldsymbol\omega,
\qquad
\vect r(\vect r\cdot\boldsymbol\omega)=(\vect r\otimes\vect r)\cdot\boldsymbol\omega .
\]
两个式子都作用在同一个 $\boldsymbol\omega$ 上，所以可以先把算子加起来再作用（这就是“$\ten J$ 是一台机器”的具体含义）。对 $\boldsymbol\omega$ 线性，所以可以提到积分外面：
\[
\boxed{\;\vect L=\ten J\cdot\boldsymbol\omega,\qquad
\ten J=\int\bigl(r^2\unitd-\vect r\otimes\vect r\bigr)\dd m\;}
\]
这就是\textbf{转动惯量张量}。它的结构值得读一遍：$r^2\unitd$ 让所有方向一视同仁，$\vect r\otimes\vect r$ 把沿 $\vect r$ 的那一份减掉。正因为有第二项，$\vect L$ 才一般\textbf{不平行}于 $\boldsymbol\omega$（\S4.6 给具体数值）。

\textbf{离散核对。}对一个位于 $\vect r$、质量为 $m$ 的\textbf{质点}（质量集中在一个点的模型），
\[
\ten J=m(r^2\unitd-\vect r\otimes\vect r),
\qquad
\ten J\cdot\boldsymbol\omega=m\bigl[r^2\boldsymbol\omega-\vect r(\vect r\cdot\boldsymbol\omega)\bigr]=m\,\vect r\times(\boldsymbol\omega\times\vect r),
\]
与直接算 $m\vect r\times\vect v$ 一致。连续体就是把许多这样的质点加起来。

\subhdr{4.2\quad 分量、几何意义、为什么对称与半正定}
\label{subsec:inertia-components}

把 $\ten J$ 逐分量写出来（用 $\delta_{ij}$ 表示单位并矢的分量）：
\[
J_{ij}=\int\bigl(r^2\delta_{ij}-x_ix_j\bigr)\dd m
\qquad\Longrightarrow\qquad
\ten J=
\begin{pmatrix}
\int(y^2+z^2)\dd m & -\int xy\dd m & -\int xz\dd m\\[3pt]
-\int xy\dd m & \int(x^2+z^2)\dd m & -\int yz\dd m\\[3pt]
-\int xz\dd m & -\int yz\dd m & \int(x^2+y^2)\dd m
\end{pmatrix}.
\]
读法：对角元是\textbf{绕对应坐标轴的转动惯量}；非对角元是\textbf{惯性积的负值}。为什么非对角元带负号？它就是 $-\vect r\otimes\vect r$ 那一项贡献的 $-x_ix_j$，没有别的理由。

\textbf{对称性是白送的}：$r^2\delta_{ij}-x_ix_j$ 对 $i,j$ 对称，所以 $J_{ij}=J_{ji}$。
\textbf{半正定性也是白送的}（半正定就是：对任意 $\boldsymbol\omega$，二次型 $\boldsymbol\omega\cdot\ten J\cdot\boldsymbol\omega\ge0$）：对任意 $\boldsymbol\omega$，
\[
\boldsymbol\omega\cdot\ten J\cdot\boldsymbol\omega
=\int\bigl[r^2\omega^2-(\vect r\cdot\boldsymbol\omega)^2\bigr]\dd m
=\int|\vect r\times\boldsymbol\omega|^2\dd m\ge0,
\]
因为 $|\vect r\times\boldsymbol\omega|^2=r^2\omega^2-(\vect r\cdot\boldsymbol\omega)^2$。等号只在“所有质量都在过原点且平行 $\boldsymbol\omega$ 的直线上”时成立；细杆沿自身方向时会出现这个退化（那根方向的转动惯量为零），一般三维刚体不会。所以惯量张量通常正定；但不管正定与否，它总是实对称，\S1.4 的谱定理都可用：三个实特征值、三根正交主轴（正定只额外保证等值面是椭球）。

\textbf{动能}是同一个张量给出的二次型：
\[
K=\tfrac12\int v^2\dd m
=\tfrac12\int(\boldsymbol\omega\times\vect r)\cdot(\boldsymbol\omega\times\vect r)\dd m
=\tfrac12\boldsymbol\omega\cdot\ten J\cdot\boldsymbol\omega .
\]
这里出现了 \S1.1 那句“课件 p9 的二次型”：$K$ 就是 $\boldsymbol\omega$ 的二次型，系数表就是 $\ten J$。

\textbf{绕任意方向的转动惯量。}对单位矢量 $\vect n$，把 $\boldsymbol\omega=\vect n$ 代入动能式，得到
\[
I(\vect n)=\vect n\cdot\ten J\cdot\vect n=\int|\vect r\times\vect n|^2\dd m,
\]
这就是“绕 $\vect n$ 轴的转动惯量”。它是 \S1.5 的二次型，因此由瑞利商，介于最小与最大主转动惯量之间。这个读法在 P4 里直接用。

\subhdr{4.3\quad 主轴 = 惯量主轴，主值 = 主转动惯量}
\label{subsec:principal}

在主轴系里 $\ten J=\operatorname{diag}(A,B,C)$（$A,B,C$ 就是三个\textbf{主转动惯量}），于是
\[
\vect L=(A\omega_1,\;B\omega_2,\;C\omega_3),
\qquad
K=\tfrac12\bigl(A\omega_1^2+B\omega_2^2+C\omega_3^2\bigr).
\]
三个主值满足三角形不等式
\[
A+B\ge C\quad\text{（以及另外两个循环）}.
\]
验证只需一行：
\[
A+B-C=\int\bigl[(y^2+z^2)+(x^2+z^2)-(x^2+y^2)\bigr]\dd m=2\int z^2\dd m\ge0 .
\]
这条不等式是很便宜的验算工具：算出来的三个主值若不满足，一定哪一步错了。它是任何质量分布都必须满足的\textbf{必要条件}；本稿只把它当检查线索，不讨论反过来的实现问题。

\textbf{对称轴的方向就是主轴。}如果刚体有一个几何对称轴，那么绕该轴转 $180^\circ$ 会把刚体映回自身，惯量张量在这个转动下不变。取该轴为 $\vect e_3$，绕 $\vect e_3$ 转 $180^\circ$ 把 $(x,y,z)$ 送到 $(-x,-y,z)$，质量分布不变，所以
\[
J_{13}=-\int xz\dd m
\quad\text{与}\quad
-\int(-x)z\dd m=+\int xz\dd m
\]
必须相等，只能 $J_{13}=0$；同理 $J_{23}=0$。于是 $\ten J\vect e_3$ 没有 $\vect e_1,\vect e_2$ 分量，只能沿 $\vect e_3$，所以 $\vect e_3$ 是一个特征方向；剩下两个主轴在垂直平面内由同样的谱定理给出（把 \S1.4 的结论用在垂直平面内的 $2\times2$ 块上）。这条判断能省掉大量特征值计算：
\begin{itemize}
\item 均质长方体：三根棱方向都是主轴（每根轴转 $180^\circ$ 都不变）；
\item 正圆锥、圆柱：对称轴是主轴；另两根在垂直于它的平面内任意取（那两个主值相等）；
\item 均质正四面体：三根主轴过对棱中点连线，三个主值相等；
\item 均质球：三个主值相等，任意方向都是主轴。
\end{itemize}

\begin{center}
\begin{tikzpicture}[scale=1.75,>={Stealth[length=1.7mm]}]
  \fill[bandviolet] (0,0)--(2.4,0)--(3.0,1.1)--(0.6,1.1)--cycle;
  \draw[violet,thick] (0,0)--(2.4,0)--(3.0,1.1)--(0.6,1.1)--cycle;
  \coordinate (C) at (1.5,0.55);
  \draw[vv] (0.45,0.55)--(2.85,0.55) node[right]{\scriptsize $J_1$};
  \draw[vv] (1.2845,0.1549)--(1.7155,0.9451) node[above right]{\scriptsize $J_2$};
  \draw[inkgray,thick] (C) circle (0.09);
  \fill (C) circle (0.02);
  \node[inkgray,above left] at (1.46,0.64) {\scriptsize $J_3$};
\end{tikzpicture}\\[2pt]
\figcap{\mbox{\textbf{图 6}（自绘）：}一块均质矩形\textbf{薄板}（厚度忽略，质量分布在板面内）的三根惯量主轴。$J_1$、$J_2$ 沿两条棱（在板面内），$J_3$ 垂直板面（画成 $\odot$）。斜投影下 $J_1$ 与 $J_2$ 看起来不垂直，实际在板面内是垂直的。绕 $J_3$、$J_1$ 或 $J_2$ 转 $180^\circ$，板都回到自己身上，所以三根都是主轴。垂直轴定理 $J_3=J_1+J_2$ \textbf{只对这样的平面薄板成立}（\S4.4）。}
\end{center}

\subhdr{4.4\quad 常用结果：细杆、薄板、长方体}
\label{subsec:standard}

先算一个一维积分，后面两个公式都从它来。均质细杆长 $L$、质量 $m$，沿 $x$ 轴放，中点在原点。线密度是 $m/L$，所以
\[
\int_{-L/2}^{L/2}x^2\dd m
=\int_{-L/2}^{L/2}x^2\frac{m}{L}\dd x
=\frac{m}{L}\cdot\frac{x^3}{3}\Big|_{-L/2}^{L/2}
=\frac{mL^2}{12}.
\]
这就是细杆绕中点垂直轴的转动惯量；杆绕自身轴（$x$ 轴）的惯量为零，因为质量都在轴上。

\textbf{矩形薄板}（$a\times b$，板面为 $xy$ 面，厚度忽略）：
\[
J_{xx}=\int y^2\dd m=\frac{mb^2}{12},\qquad
J_{yy}=\int x^2\dd m=\frac{ma^2}{12},\qquad
J_{zz}=\int(x^2+y^2)\dd m=\frac{m(a^2+b^2)}{12}.
\]
最后一条就是\textbf{垂直轴定理} $J_{zz}=J_{xx}+J_{yy}$：$z$ 方向没有厚度，所以绕 $z$ 的惯量恰好是绕 $x$、$y$ 的惯量之和。它在并矢语言里不是新定理，只是一行读数。

\textbf{长方体}（边长 $a,b,c$，质量 $m$，质心为原点，棱沿坐标轴）：
\[
\ten J=\frac{m}{12}\operatorname{diag}\bigl(b^2+c^2,\;a^2+c^2,\;a^2+b^2\bigr).
\]
推导就是把 $\int(y^2+z^2)\dd m=\int y^2\dd m+\int z^2\dd m$ 分开，各用上面的杆公式；$y$ 方向长度 $b$ 给 $mb^2/12$，$z$ 方向长度 $c$ 给 $mc^2/12$。

\subhdr{4.5\quad 平行轴定理}
\label{subsec:parallel}

把原点从质心 $C$ 挪到 $O$，设 $\vect R$ 是从 $O$ 指向 $C$ 的矢量，质元相对 $O$ 的位置是 $\vect r=\vect R+\vect r'$（$\vect r'$ 相对质心）。由质心定义 $\int\vect r'\dd m=\vect0$，展开
\[
\begin{aligned}
\ten J_O&=\int\bigl[(\vect R+\vect r')^2\unitd-(\vect R+\vect r')\otimes(\vect R+\vect r')\bigr]\dd m\\
&=\int\bigl(r'^2\unitd-\vect r'\otimes\vect r'\bigr)\dd m
+m\bigl(R^2\unitd-\vect R\otimes\vect R\bigr)\\
&\quad+\int\bigl[2(\vect R\cdot\vect r')\unitd-\vect R\otimes\vect r'-\vect r'\otimes\vect R\bigr]\dd m .
\end{aligned}
\]
最后一行三个积分的被积函数都含 $\vect r'$ 的一次因子，而 $\int\vect r'\dd m=\vect0$（质心定义），所以整行消失；剩下的第一项是 $\ten J_C$，第二项是 $m(R^2\unitd-\vect R\otimes\vect R)$。于是
\[
\boxed{\;\ten J_O=\ten J_C+m\bigl(R^2\unitd-\vect R\otimes\vect R\bigr)\;}
\]
\textbf{注意第二项是并矢，不能写成 $mR^2$ 一个标量}——绕不同方向的轴平移，增量不同。标量版 $J=J_C+md^2$ 只是它在“轴平行且过 $O$”这一个特例上的读数：取单位轴 $\vect n$，则
\[
\vect n\cdot\bigl[m(R^2\unitd-\vect R\otimes\vect R)\bigr]\cdot\vect n
=m\bigl(R^2-(\vect n\cdot\vect R)^2\bigr)=m d^2,
\]
其中 $d$ 是质心到轴的垂直距离。这一步把“并矢增量”与常见的标量平行轴定理接上了：后者不是另一条定理，而是前者沿某根轴取二次型。

\textbf{惯量椭球（选读，可跳过）。}若 $\ten J$ 正定，令 $\vect n\cdot\ten J\cdot\vect n=1$，得到惯量椭球；它的三个半轴是 $1/\sqrt{A},1/\sqrt{B},1/\sqrt{C}$，方向就是主轴。这个图像与 \S1.2 的椭球完全一样，只是把“矩阵”换成了“惯量张量”。跳过它不影响 \S4.6 与 \S5。

\subhdr{4.6\quad 教学示范：$\vect L$ 到底偏多少}
\label{subsec:inertia-example}

取均质长方体 $a=1$、$b=2$、$c=3$ m，质量 $m=12$ kg，质心在原点。由 \S4.4，
\[
\ten J=\frac{12}{12}\operatorname{diag}(2^2+3^2,\;1^2+3^2,\;1^2+2^2)
=\operatorname{diag}(13,\;10,\;5)\ \text{kg·m}^2 .
\]
让它绕一个与三条棱都成等角的方向转，$\boldsymbol\omega=\omega_0(1,1,1)/\sqrt3$，取 $\omega_0=1$ rad/s。注意这\textbf{不是}这个长方体的空间对角线；从 $(0,0,0)$ 到 $(1,2,3)$ 的空间对角线方向是 $(1,2,3)/\sqrt{14}$。这里选 $(1,1,1)/\sqrt3$，是因为它与三条棱的夹角相等、计算干净，而下面的结论只用到“$\boldsymbol\omega$ 与三根主轴都不平行”。则
\[
\vect L=\ten J\cdot\boldsymbol\omega=\frac{1}{\sqrt3}(13,\;10,\;5)\approx(7.506,\;5.774,\;2.887)\ \text{kg·m}^2/\text{s}.
\]
$\vect L$ 与 $\boldsymbol\omega\propto(1,1,1)$ 明显不平行。夹角由
\[
\cos\theta=\frac{\vect L\cdot\boldsymbol\omega}{|\vect L|\,|\boldsymbol\omega|}
=\frac{28/3}{\sqrt{294/3}}
=\frac{28}{3\sqrt{98}}=\frac{4}{3\sqrt2}=\frac{2\sqrt2}{3}\approx0.9428
\quad\Longrightarrow\quad
\theta\approx19.47^\circ .
\]
核对：$|\vect L|=\sqrt{13^2+10^2+5^2}/\sqrt3=\sqrt{294/3}=\sqrt{98}=7\sqrt2$，$\vect L\cdot\boldsymbol\omega=(13+10+5)/3=28/3$，两者相除时 $\sqrt3$ 与公因子都约掉；$\sin\theta=\sqrt{1-8/9}=1/3$。

\textbf{若强行让它不晃，需要多大力矩？}若 $\boldsymbol\omega$ 在体坐标系里保持不动，则 $\vect L$ 在体坐标系里也不动，$\dd\vect L/\dd t|_{\text{体}}=\vect0$，固定系里的变化只剩“整体转动”那一项（\S5.1）：
\[
\vect M=\boldsymbol\omega\times\vect L
=\frac{1}{3}\begin{pmatrix}1\\1\\1\end{pmatrix}\times\begin{pmatrix}13\\10\\5\end{pmatrix}
=\frac{1}{3}\begin{pmatrix}5-10\\13-5\\10-13\end{pmatrix}
=\frac{1}{3}\begin{pmatrix}-5\\8\\-3\end{pmatrix},
\]
\[
|\vect M|=\frac{1}{3}\sqrt{25+64+9}=\frac{\sqrt{98}}{3}=\frac{7\sqrt2}{3}\approx3.30\ \text{N·m}.
\]
方向垂直于 $\boldsymbol\omega$ 与 $\vect L$ 张成的平面（三者的共面关系见图 8）。注意这不是“把 $\vect L$ 扳向 $\boldsymbol\omega$”：那样做是缩小两者的夹角，转轴是 $\vect L\times\boldsymbol\omega$，方向正好与 $\vect M$ 相反；沿 $\vect M$ 加力矩的效果是让 $\vect L$ 绕着 $\boldsymbol\omega$ 进动。\textbf{自由转动时没有这个力矩，所以 $\boldsymbol\omega$ 会自己动起来}——这就是“刚体转起来会晃”的数学原因，\S5.4 会把这个力矩与欧拉方程接上。

\hdr{§5}{落点二：输运定理、欧拉方程与对称陀螺}
\label{sec:transport}

\subhdr{5.1\quad 输运定理：为什么不能直接对 $\vect L$ 求导}
\label{subsec:transport}

初学者最常写的一步错是
\[
\frac{\dd\vect L}{\dd t}=\ten J\cdot\frac{\dd\boldsymbol\omega}{\dd t}
\qquad\text{（错，除非 }\ten J\text{ 在该系里是常量）}.
\]
$\ten J$ 是“跟着刚体转”的张量。在固定系里看，它的九个分量随时间变；在体坐标系里看，它才是不变的。两个系里的时间导数由一个公式联系，叫\textbf{输运定理}。

\begin{center}
\begin{tikzpicture}[scale=1.10,>={Stealth[length=1.7mm]}]
  % 左：固定系
  \begin{scope}
    \draw[ax] (-1.5,0)--(1.65,0) node[right]{\scriptsize $\vect e_1$};
    \draw[ax] (0,-1.15)--(0,1.6) node[above]{\scriptsize $\vect e_2$};
    \begin{scope}[rotate=32]
      \draw[violet,thick,fill=bandviolet] (0,0) ellipse (1.05 and 0.42);
    \end{scope}
    \draw[violet,densely dashed,->] (-1.06,-0.66)--(1.06,0.66);
    \draw[violet,densely dashed,->] (0.53,-0.85)--(-0.53,0.85);
    \node[violet] at (1.20,0.80) {\scriptsize $\vect n'_1$};
    \node[violet] at (-0.72,0.92) {\scriptsize $\vect n'_2$};
    \draw[inkgray,->] (0.34,1.16) arc (72:38:1.18);
    \node[inkgray] at (0.92,1.36) {\scriptsize $\boldsymbol\Omega$};
    \node[align=center] at (0,-1.70) {\scriptsize 固定系：$\vect n'_k(t)$ 在转，\\ $\ten J$ 的九个数随时间变};
  \end{scope}
  % 右：体系
  \begin{scope}[xshift=6.4cm]
    \draw[ax] (-1.36,0.848)--(1.62,-1.012) node[above right]{\scriptsize $\vect e_1$};
    \draw[ax] (-0.795,-1.272)--(0.901,1.442) node[above right]{\scriptsize $\vect e_2$};
    \draw[violet,thick,fill=bandviolet] (0,0) ellipse (1.05 and 0.42);
    \draw[violet,thick,->] (0,0)--(1.45,0) node[above]{\scriptsize $\vect n'_1$};
    \draw[violet,thick,->] (0,0)--(0,1.25) node[above right]{\scriptsize $\vect n'_2$};
    \node[align=center] at (0,-1.70) {\scriptsize 体坐标系：$\vect n'_k$ 定死（画面上是正的），\\ $\ten J=\operatorname{diag}(A,B,C)$ 不变};
  \end{scope}
\end{tikzpicture}\\[2pt]
\figcap{\mbox{\textbf{图 7}（自绘）：}同一块刚体，两种看法。左边在固定系里看：主轴 $\vect n'_k$ 在转，所以 $\ten J$ 的账本随时间变，“$\dd\vect L/\dd t=\ten J\cdot\dd\boldsymbol\omega/\dd t$”是错的。右边改用体坐标系：把 $\vect n'_k$ 当成不动的坐标轴（灰虚线的 $\vect e_k$ 反过来在转），$\ten J$ 的账本就是常量 $\operatorname{diag}(A,B,C)$，代价是多出与转动有关的项（陀螺项）。}
\end{center}

\textbf{矢量版。}设随体基 $\vect n'_k$ 以角速度 $\boldsymbol\Omega$ 转动（\S3.5 的结论：$\dd\vect n'_k/\dd t=\boldsymbol\Omega\times\vect n'_k$）。把矢量写成 $\vect A=A_k\vect n'_k$，其中 $A_k$ 是体坐标分量。求导：
\[
\Bigl(\frac{\dd\vect A}{\dd t}\Bigr)_{\text{固定}}
=\dot A_k\vect n'_k+A_k\frac{\dd\vect n'_k}{\dd t}
=\Bigl(\frac{\dd\vect A}{\dd t}\Bigr)_{\text{体}}+\boldsymbol\Omega\times\vect A .
\]
第一项是“体坐标分量变化”的贡献，第二项是“基矢量自己在转”的贡献。这就是相对运动里熟悉的输运公式。

\textbf{并矢版。}设 $\ten D=D_{ij}\vect n'_i\otimes\vect n'_j$（$D_{ij}$ 是体坐标分量）。逐项求导时三项都不能漏：
\[
\begin{aligned}
\Bigl(\frac{\dd\ten D}{\dd t}\Bigr)_{\text{固定}}
&=\dot D_{ij}\,\vect n'_i\otimes\vect n'_j
+D_{ij}\Bigl(\frac{\dd\vect n'_i}{\dd t}\Bigr)\otimes\vect n'_j
+D_{ij}\,\vect n'_i\otimes\Bigl(\frac{\dd\vect n'_j}{\dd t}\Bigr)\\
&=\Bigl(\frac{\dd\ten D}{\dd t}\Bigr)_{\text{体}}
+D_{ij}(\boldsymbol\Omega\times\vect n'_i)\otimes\vect n'_j
+D_{ij}\,\vect n'_i\otimes(\boldsymbol\Omega\times\vect n'_j).
\end{aligned}
\]
第一项是体分量自己的变化率，就是 $(\dd\ten D/\dd t)_{\text{体}}$。后两项各把一个基矢量换掉；为把“换基”写成一台机器，用 \S2.6 的对偶张量 $\boldsymbol\Omega^{\times}$，并规定它对并矢的作用为
\[
\boldsymbol\Omega^{\times}\cdot(\vect a\otimes\vect b)\stackrel{\text{def}}{=}(\boldsymbol\Omega\times\vect a)\otimes\vect b,
\qquad
(\vect a\otimes\vect b)\cdot\boldsymbol\Omega^{\times}\stackrel{\text{def}}{=}\vect a\otimes(\vect b\cdot\boldsymbol\Omega^{\times})
=-\vect a\otimes(\boldsymbol\Omega\times\vect b).
\]
第二个定义用了 $(\boldsymbol\Omega^{\times})^{\mathsf T}=-\boldsymbol\Omega^{\times}$，所以右边的负号不是另加的。按线性，后两项正好是 $\boldsymbol\Omega^{\times}\cdot\ten D$ 与 $-\ten D\cdot\boldsymbol\Omega^{\times}$，于是
\[
\boxed{\;\Bigl(\frac{\dd\ten D}{\dd t}\Bigr)_{\text{固定}}
=\Bigl(\frac{\dd\ten D}{\dd t}\Bigr)_{\text{体}}
+\boldsymbol\Omega^{\times}\cdot\ten D-\ten D\cdot\boldsymbol\Omega^{\times}\;}
\]
这里两个乘积都是\textbf{普通矩阵乘法}：$\boldsymbol\Omega^{\times}\cdot\ten D$ 和 $\ten D\cdot\boldsymbol\Omega^{\times}$ 各把张量的一个指标“转一次”，两个指标地位相反，所以一正一负。

\textbf{单并矢核对。}取 $\ten D=\vect a\otimes\vect b$，按定义
\[
\boldsymbol\Omega^{\times}\cdot\ten D=(\boldsymbol\Omega\times\vect a)\otimes\vect b,
\qquad
\ten D\cdot\boldsymbol\Omega^{\times}=-\vect a\otimes(\boldsymbol\Omega\times\vect b).
\]
所以右边两项合起来是 $(\boldsymbol\Omega\times\vect a)\otimes\vect b+\vect a\otimes(\boldsymbol\Omega\times\vect b)$，正好是乘积求导的两项。用 $\vect a=\vect e_1$、$\vect b=\vect e_2$、$\boldsymbol\Omega=\vect e_3$ 核验：
\[
(\boldsymbol\Omega\times\vect e_1)\otimes\vect e_2+\vect e_1\otimes(\boldsymbol\Omega\times\vect e_2)
=\vect e_2\otimes\vect e_2-\vect e_1\otimes\vect e_1
=\operatorname{diag}(-1,1,0),
\]
而
\[
\boldsymbol\Omega^{\times}\cdot(\vect e_1\otimes\vect e_2)-(\vect e_1\otimes\vect e_2)\cdot\boldsymbol\Omega^{\times}
=\vect e_2\otimes\vect e_2-\vect e_1\otimes\vect e_1,
\]
两边相同。这个例子既非零又非对称，能暴露转置或符号错误。

\textbf{单位张量的检查。}取 $\ten D=\unitd$，左边为 $\vect0$（单位张量永远是单位张量），右边
\[
\boldsymbol\Omega^{\times}\cdot\unitd-\unitd\cdot\boldsymbol\Omega^{\times}
=\boldsymbol\Omega^{\times}-\boldsymbol\Omega^{\times}=\vect0 .
\]

\subhdr{5.2\quad 体坐标系里的角动量定理与欧拉方程}
\label{subsec:euler}

固定系里角动量定理是
\[
\Bigl(\frac{\dd\vect L}{\dd t}\Bigr)_{\text{固定}}=\vect M,
\]
其中 $\vect M$ 是外力对定点（或质心）的力矩。力矩就是“力对点的矩” $\vect M=\int\vect r\times\dd\vect F$。这条角动量定理不是新的独立假设，而是牛顿第二定律对质点系求和的结果：对第 $i$ 个质点，$\vect r_i\times m_i\vect a_i=\vect r_i\times\vect F_i$；对 $i$ 求和，左边是
\[
\sum_i\vect r_i\times m_i\vect a_i
=\frac{\dd}{\dd t}\sum_i\vect r_i\times m_i\vect v_i
=\frac{\dd\vect L}{\dd t},
\]
右边是外力矩之和加上内力矩之和。内力成对出现、大小相等方向相反，且沿两质点的连线，所以每一对内力矩 $\vect r_i\times\vect F_{ij}+\vect r_j\times\vect F_{ji}=(\vect r_i-\vect r_j)\times\vect F_{ij}$ 为零；内力矩整体相消，只剩外力矩 $\vect M$。条件是惯性系、对定点取矩；对质心取矩时把 $\vect r_i$ 换成相对质心的位置，平动项在求和中抵消，结论同样成立。换成体坐标系（体坐标系以 $\boldsymbol\omega$ 转动），用矢量版输运公式：
\[
\Bigl(\frac{\dd\vect L}{\dd t}\Bigr)_{\text{体}}+\boldsymbol\omega\times\vect L=\vect M .
\]
体坐标系里 $\ten J=\operatorname{diag}(A,B,C)$ 是常量，$\vect L=(A\omega_1,B\omega_2,C\omega_3)$，于是 $(\dd\vect L/\dd t)_{\text{体}}=(A\dot\omega_1,B\dot\omega_2,C\dot\omega_3)$。再算 $\boldsymbol\omega\times\vect L$ 的三个分量：
\[
(\boldsymbol\omega\times\vect L)_1=\omega_2(C\omega_3)-\omega_3(B\omega_2)=(C-B)\omega_2\omega_3,
\]
另两个循环。合起来就是\textbf{欧拉方程}：
\[
\boxed{
\begin{aligned}
A\dot\omega_1+(C-B)\omega_2\omega_3&=M_1,\\
B\dot\omega_2+(A-C)\omega_3\omega_1&=M_2,\\
C\dot\omega_3+(B-A)\omega_1\omega_2&=M_3 .
\end{aligned}}
\]
无外力矩（$M_k=0$）时，这三条就是刚体自由转动的全部动力学。$A,B,C$ 是体坐标系里的三个主转动惯量，$\omega_k$ 是角速度在体坐标系的三个分量，$M_k$ 是外力矩的体坐标分量。

\subhdr{5.3\quad 对称陀螺：$A=B\ne C$}
\label{subsec:top}

取 $A=B$、$\vect M=\vect0$。第三条方程立刻给 $C\dot\omega_3=0$，即
\[
\omega_3=\text{常量}.
\]
前两条变成
\[
\dot\omega_1=-\Omega\,\omega_2,\qquad \dot\omega_2=\Omega\,\omega_1,
\qquad \Omega\equiv\frac{C-A}{A}\,\omega_3 .
\]
这里 $\Omega$ 是一个带符号的标量（横向转动的角速度），不要与前面粗体的角速度矢量 $\boldsymbol\Omega$ 混淆；$\Omega$ 的符号决定横向分量是顺着还是逆着 $\vect e_3$ 转。这两条不借助微分方程理论也能读懂。把横向部分看成一个二维矢量 $\vect w_\perp=(\omega_1,\omega_2)$，上式说
\[
\frac{\dd}{\dd t}\begin{pmatrix}\omega_1\\\omega_2\end{pmatrix}
=\Omega\begin{pmatrix}-\omega_2\\\omega_1\end{pmatrix},
\]
即横向“速度”与 $\vect w_\perp$ 垂直、大小是 $|\Omega|\,|\vect w_\perp|$。先看大小：把上式与 $(\omega_1,\omega_2)$ 点乘，
\[
\omega_1\dot\omega_1+\omega_2\dot\omega_2
=\Omega(-\omega_1\omega_2+\omega_2\omega_1)=0,
\]
所以 $\omega_1^2+\omega_2^2=\omega_\perp^2$ 是常数。再看方向：若 $\Omega>0$，速度总是把 $(\omega_1,\omega_2)$ 朝逆时针方向转 $90^\circ$；所以它在以原点为圆心、半径 $\omega_\perp$ 的圆上匀速转，角速度就是 $\Omega$（$\Omega<0$ 时反向）。

把这句话写成公式：
\[
\omega_1=\omega_\perp\cos(\Omega t+\varphi_0),\qquad
\omega_2=\omega_\perp\sin(\Omega t+\varphi_0),
\]
其中 $\omega_\perp\ge0$ 是圆半径，$\varphi_0$ 是 $t=0$ 时横向矢量与 $\vect e_1$ 的夹角。直接代回可以当场核对：
\[
\dot\omega_1=-\Omega\omega_\perp\sin(\Omega t+\varphi_0)=-\Omega\omega_2,
\qquad
\dot\omega_2=\Omega\omega_\perp\cos(\Omega t+\varphi_0)=\Omega\omega_1 .
\]
给定初始横向矢量 $\vect w_\perp(0)=(\omega_1(0),\omega_2(0))$，两个常数就定下来了（$\omega_\perp>0$ 时；若 $\omega_\perp=0$，横向分量恒为零，$\varphi_0$ 可任取）：
\[
\omega_\perp=\sqrt{\omega_1(0)^2+\omega_2(0)^2},
\qquad
\cos\varphi_0=\frac{\omega_1(0)}{\omega_\perp},\quad
\sin\varphi_0=\frac{\omega_2(0)}{\omega_\perp}.
\]
例如 $t=0$ 时 $\vect w_\perp$ 沿 $\vect e_1$，则 $\varphi_0=0$，解是 $\omega_1=\omega_\perp\cos\Omega t$、$\omega_2=\omega_\perp\sin\Omega t$；若沿 $\vect e_2$，则 $\varphi_0=\pi/2$。这样从任意给定的初始横向矢量都能写出解、也能代回核对，不需要预先学过“简谐振子”。同一个数学结构在其它问题里常叫“简谐振子”；本稿不需要那个名称，也不依赖它的结论。

\textbf{结论：}角速度矢量的横向分量在体坐标系里绕体轴 $3$ 匀速转动，角速度 $\Omega$（可正可负，负号表示反向）、大小 $\omega_\perp$ 不变，$\omega_3$ 也不变，所以 $|\boldsymbol\omega|$ 不变。严格说，这里是 $\boldsymbol\omega$ 在体坐标系里绕 $\vect e_3$ 转，不是陀螺本身在空间中的进动；后者还要再换回固定系。所谓\textbf{进动}，就是转轴或角速度方向绕着另一条轴慢慢转。

\textbf{几何图像。}$\vect L=(A\omega_1,A\omega_2,C\omega_3)$。因为 $\vect L$ 的前两个分量只是 $\boldsymbol\omega$ 的前两个分量乘同一个常数 $A$，所以 $\boldsymbol\omega$、$\vect L$、$\vect e_3$ \textbf{三个矢量始终共面}；这个平面绕 $\vect e_3$ 以 $\Omega$ 转。这是对称陀螺最有用的图像。

\begin{center}
\begin{tikzpicture}[scale=1.6,>={Stealth[length=1.7mm]}]
  \coordinate (O) at (0,0);
  \coordinate (E3) at (0,1.6);
  \coordinate (W) at (0.937,1.040);
  \coordinate (L) at (0.876,1.217);
  \draw[ax] (O)--(E3) node[above]{\scriptsize $\vect e_3$};
  \draw[inkgray,thin,densely dashed] (0.937,1.040)--(0.937,0);
  \draw[inkgray,thin,densely dashed] (0.876,1.217)--(0,1.217);
  \draw[vv] (O)--(W) node[right]{\scriptsize $\boldsymbol\omega$};
  \draw[violet,very thick,->,>={Stealth[length=1.7mm]}] (O)--(L) node[above left]{\scriptsize $\vect L$};
  \draw[inkgray,->] (0.30,1.42) arc (20:160:0.30);
  \node[inkgray] at (0.44,1.52) {\scriptsize $\Omega$};
  \fill (O) circle (0.7pt);
\end{tikzpicture}\\[2pt]
\figcap{\mbox{\textbf{图 8}（自绘）：}对称陀螺的 $\boldsymbol\omega$、$\vect L$、$\vect e_3$ 共面。图中 $A<C$（扁盘状）的取法使 $\vect L$ 比 $\boldsymbol\omega$ 更靠近 $\vect e_3$；这个平面绕 $\vect e_3$ 以 $\Omega$ 转，$\boldsymbol\omega$ 的横向分量在体坐标系里以同一 $\Omega$ 匀速转圈。本图只画三者的共面关系与转向，不画该横向分量扫出的那个圆。}
\end{center}

\begin{center}
\fcolorbox{violet}{bandviolet}{\parbox{0.94\linewidth}{\centering\vspace{3pt}
\small $C>A$（扁盘状，如地球）时 $\Omega>0$，$\boldsymbol\omega$ 的横向分量与 $\vect L$ 绕 $\vect e_3$ 同向转；$C<A$（细长杆状）时反向。\vspace{3pt}}}
\end{center}

\textbf{教学示范。}取对称陀螺 $A=B=4$ kg·m$^2$、$C=5$ kg·m$^2$，自转 $\omega_3=6$ rad/s，无外力矩。则
\[
\Omega=\frac{C-A}{A}\omega_3=\frac{5-4}{4}\times6=1.5\ \text{rad/s},
\qquad
T=\frac{2\pi}{\Omega}=\frac{2\pi}{1.5}=\frac{4\pi}{3}\approx4.19\ \text{s}.
\]
即 $\boldsymbol\omega$ 的横向分量每 $4.19$ 秒绕体轴转一圈，$\omega_3$ 与 $|\boldsymbol\omega|$ 都不变。$C>A$，属于“扁盘状”，$\Omega>0$，横向分量沿 $+\vect e_3$ 的正向转。

\subhdr{5.4\quad 回到 \S4.6：维持对角转动需要多大力矩}
\label{subsec:torque}

把“$\vect L$ 与 $\boldsymbol\omega$ 不平行”代进 $\vect M=\boldsymbol\omega\times\vect L$，就得到“维持这种转动需要多大扭矩”。\S4.6 已经对 $m=12$、$a=1,b=2,c=3$ 的长方体算过：$|\vect M|=\sqrt{98}/3\approx3.30$ N·m，方向沿 $(1/3)(-5,8,-3)$。

一般地，若 $\boldsymbol\omega$ 在体坐标系里保持不动，则
\[
\vect M=\boldsymbol\omega\times\vect L,\qquad
|\vect M|=|\boldsymbol\omega|\,|\vect L|\sin\theta,
\]
其中 $\theta$ 是 $\boldsymbol\omega$ 与 $\vect L$ 的夹角；方向垂直于 $\boldsymbol\omega$ 与 $\vect L$ 张成的平面。\textbf{自由转动时没有这个力矩，所以 $\boldsymbol\omega$ 必须自己动起来}——这正是 \S5.3 那条进动的来源。反过来说：\textbf{若 $\vect M=\vect0$ 且转动是稳定的，则 $\boldsymbol\omega\times\vect L=\vect0$，即 $\vect L\parallel\boldsymbol\omega$，$\boldsymbol\omega$ 必须沿某根主轴。}这就是“只有绕主轴的匀速转动可以无扭矩维持”的并矢解释。

\hdr{§6}{落点三：应力张量}
\label{sec:stress}

\subhdr{6.1\quad 面力是法向的函数：柯西应力与商法则}
\label{subsec:cauchy}

在连续体里（把物体看成连续分布的物质，而不是一个个离散质点）过一点取一个截面，法向为单位矢量 $\vect n$。截面上单位面积受到的力记作 $\vect t(\vect n)$，叫\textbf{面力矢量}（traction）。它描述的是“法向为 $\vect n$ 的那一侧材料，对另一侧单位面积的作用力”。

\textbf{柯西应力定理说：$\vect t$ 是 $\vect n$ 的线性函数。}为什么？取一个微小的四面体：三个面分别垂直于 $\vect e_1,\vect e_2,\vect e_3$，第四个面的法向是 $\vect n$。设第四个面的面积为 $\dd S$，三个坐标面的面积分别是 $\dd S\,n_1,\dd S\,n_2,\dd S\,n_3$（面积投影：一个面的面积乘上它与坐标面的夹角的余弦，正好是法向分量）。这里取 $\vect n$ 所在卦限的四面体，使 $n_1,n_2,n_3\ge0$；若 $\vect n$ 的分量有负号，就取对应卦限的坐标面，下面的线性式仍成立。四面体静止时合力为零；体积力正比于体积（三阶小量），当四面体缩小时可以忽略，只剩下四个面上的面力：
\[
\vect t(\vect n)\dd S+\vect t(-\vect e_1)\dd S\,n_1+\vect t(-\vect e_2)\dd S\,n_2+\vect t(-\vect e_3)\dd S\,n_3=\vect0 .
\]
用 $\vect t(-\vect e_j)=-\vect t(\vect e_j)$（作用与反作用），除 $\dd S$ 得
\[
\vect t(\vect n)=n_1\vect t(\vect e_1)+n_2\vect t(\vect e_2)+n_3\vect t(\vect e_3).
\]
\textbf{这就是线性性，而且是“对任意 $\vect n$ 都成立”的线性性。}把 $\vect t(\vect e_j)$ 的三个分量记作 $\sigma_{ij}$：
\[
\vect t(\vect e_j)=\sigma_{ij}\vect e_i,
\qquad
\boxed{\;\vect t(\vect n)=\ten\sigma\cdot\vect n,\qquad t_i=\sigma_{ij}n_j\;}
\]
其中 $\sigma_{ij}$ 读作“\textbf{法向为 $\vect e_j$ 的面上、沿 $\vect e_i$ 方向的应力分量}”。第一个指标是力的方向，第二个指标是面的法向——这个顺序在 \S6.4 的散度里会再出现一次。

\begin{center}
\begin{tikzpicture}[scale=1.65,>={Stealth[length=1.7mm]}]
  \coordinate (O) at (0,0);
  \coordinate (A) at (2.0,0);
  \coordinate (B) at (0,1.55);
  \coordinate (C) at (-0.55,-0.55);
  \draw[inkgray,thin] (O)--(A);
  \draw[inkgray,thin] (O)--(B);
  \draw[inkgray,thin] (O)--(C);
  \draw[violet,thick] (A)--(B)--(C)--cycle;
  \fill[bandviolet,opacity=0.45] (A)--(B)--(C)--cycle;
  \draw[vv] (O)--(A) node[below right]{\scriptsize $\vect e_1$};
  \draw[vv] (O)--(B) node[above right]{\scriptsize $\vect e_2$};
  \draw[vv] (O)--(C) node[below left]{\scriptsize $\vect e_3$};
  \node[inkgray,align=center,fill=white,inner sep=1.5pt] at (1.55,0.42) {\scriptsize 斜面：法向 $\vect n$，\\ 面积 $\dd S$};
\end{tikzpicture}\\[2pt]
\figcap{\mbox{\textbf{图 9}（自绘）：}柯西四面体。三个坐标面上的面力与斜面（法向 $\vect n$，面积 $\dd S$）上的面力平衡；坐标面的面积是 $\dd S n_j$。四面体缩小时体积力是高阶小量，所以 $\vect t(\vect n)=\sum_j n_j\vect t(\vect e_j)$。图中只标出斜面的面与法向名称，不把透视投影的短线当作真实法向；$\vect n$ 的具体方向由所取截面决定。}
\end{center}

\textbf{为什么它一定是张量？}因为 $\vect t$ 与 $\vect n$ 都是矢量，而“对任意 $\vect n$ 成立 $\vect t=\ten\sigma\vect n$”正是 \S2.7 商法则的前提，所以 $\ten\sigma$ 自动是二阶张量。\textbf{这一步不是假设，是被逼出来的}——这也是 \S2.7 存在的理由。

\subhdr{6.2\quad 对称性来自角动量守恒}
\label{subsec:stresssym}

对微元取力矩平衡（不设体力矩；准静态就是每个时刻都可当作平衡处理，加速度的影响可忽略）。考虑一个小立方体，边长 $h$。作用在法向为 $\vect e_1$ 的面上的切应力 $\sigma_{21}$，与作用在法向为 $\vect e_2$ 的面上的切应力 $\sigma_{12}$，分别产生绕第三个轴的力矩。

先把符号算清楚。$+\vect e_1$ 面上的力是 $\sigma_{21}h^2\vect e_2$，作用点在 $(h/2)\vect e_1$；$-\vect e_1$ 面上的力是 $-\sigma_{21}h^2\vect e_2$，作用点在 $-(h/2)\vect e_1$。这一对力对中心的力矩是
\[
\Bigl(\frac h2\vect e_1\Bigr)\times(\sigma_{21}h^2\vect e_2)
+\Bigl(-\frac h2\vect e_1\Bigr)\times(-\sigma_{21}h^2\vect e_2)
=\sigma_{21}h^3\vect e_3 .
\]
同理，$+\vect e_2$ 面上的力是 $\sigma_{12}h^2\vect e_1$，$-\vect e_2$ 面上的力是 $-\sigma_{12}h^2\vect e_1$，这一对力的力矩是
\[
\Bigl(\frac h2\vect e_2\Bigr)\times(\sigma_{12}h^2\vect e_1)
+\Bigl(-\frac h2\vect e_2\Bigr)\times(-\sigma_{12}h^2\vect e_1)
=-\sigma_{12}h^3\vect e_3 .
\]
两对力矩之和是 $(\sigma_{21}-\sigma_{12})h^3\vect e_3$。立方体的转动惯量正比于 $h^5$，角加速度有限时惯性力矩是 $O(h^5)$；让 $h\to0$，$h^3$ 项必须为零，所以 $\sigma_{21}=\sigma_{12}$。对另外两个轴重复同样的论证，得到
\[
\sigma_{ij}=\sigma_{ji}.
\]
三个独立等式给出 $\sigma_{12}=\sigma_{21}$、$\sigma_{13}=\sigma_{31}$、$\sigma_{23}=\sigma_{32}$。所以应力张量也\textbf{对称}，\S1.4 的谱定理同样适用：三个\textbf{主应力} $\sigma_1,\sigma_2,\sigma_3$、三个互相垂直的\textbf{主平面}（这些截面上剪应力为零）。主应力就是 $\ten\sigma$ 的特征值，主平面的法向就是特征向量。

\begin{center}
\begin{tikzpicture}[scale=1.7,>={Stealth[length=1.7mm]}]
  \draw[inkgray,thick] (-1,-1) rectangle (1,1);
  \draw[vv] (1,0)--(1,0.72) node[right]{\scriptsize $\sigma_{21}$};
  \draw[vv] (-1,0)--(-1,-0.72) node[left]{\scriptsize $-\sigma_{21}$};
  \draw[vv] (0,1)--(0.72,1) node[above]{\scriptsize $\sigma_{12}$};
  \draw[vv] (0,-1)--(-0.72,-1) node[below]{\scriptsize $-\sigma_{12}$};
  \draw[inkgray,thin] (0,0) circle (0.07);
  \fill (0,0) circle (0.02);
  \node[inkgray] at (0.35,-0.62) {\scriptsize $+\vect e_3$ 出纸面};
\end{tikzpicture}\\[2pt]
\figcap{\mbox{\textbf{图 10}（自绘）：}小立方体在 $xy$ 平面内的截面示意。$+\vect e_1$ 面上的切应力 $\sigma_{21}$ 沿 $+\vect e_2$，$-\vect e_1$ 面上的对应力反向；$+\vect e_2$ 面上的 $\sigma_{12}$ 沿 $+\vect e_1$，$-\vect e_2$ 面上的对应力反向。两对力对中心的力矩都沿 $\vect e_3$（出纸面）但方向相反，平衡给出 $\sigma_{21}=\sigma_{12}$。图中只画 $xy$ 截面，$z$ 方向厚度为 $h$。}
\end{center}

\textbf{一个与坐标无关的数。}三个主应力之和
\[
\sigma_1+\sigma_2+\sigma_3=\tr\ten\sigma
\]
是与坐标无关的（\S0.5）；它对应“静水压强”那种不随截面变化的成分，$p=-\frac13\tr\ten\sigma$（拉伸为正时，压力取负号）。把各向同性的压力部分减掉，剩下的部分叫偏斜部分，它才决定材料怎么变形。这一段不需要计算，只要知道有这层区分即可。

\subhdr{6.3\quad 教学示范：斜截面上的面力分解}
\label{subsec:oblique}

取单轴拉伸
\[
\ten\sigma=\sigma_0\,\vect e_1\otimes\vect e_1
\qquad(\sigma_0>0),
\]
即只有 $\sigma_{11}=\sigma_0$，其余分量为零。主轴就是 $\vect e_1,\vect e_2,\vect e_3$，主应力是 $\sigma_0,0,0$。

在 $xy$ 面内取与 $\vect e_1$ 成 $\theta=30^\circ$ 的截面，法向
\[
\vect n=\cos\theta\,\vect e_1+\sin\theta\,\vect e_2=\frac{\sqrt3}{2}\vect e_1+\frac12\vect e_2 .
\]
面力：
\[
\vect t=\ten\sigma\cdot\vect n
=\sigma_0\,\vect e_1(\vect e_1\cdot\vect n)
=\sigma_0\cos\theta\,\vect e_1
=\frac{\sqrt3}{2}\sigma_0\,\vect e_1 .
\]
它仍然沿 $\vect e_1$，但斜截面把它切成两个分量。法向分量：
\[
t_n=\vect t\cdot\vect n=\sigma_0\cos^2\theta=\frac34\sigma_0,
\]
法向矢量是 $t_n\vect n$。切向矢量是 $\vect t-t_n\vect n$，其大小
\[
\tau=\sigma_0\sin\theta\cos\theta=\frac{\sqrt3}{4}\sigma_0,
\]
方向是 $(\sin\theta,-\cos\theta)=(1/2,-\sqrt3/2)$（沿截面方向）。核对：$|\vect t|^2=\frac34\sigma_0^2$，$t_n^2+\tau^2=\frac{9}{16}\sigma_0^2+\frac{3}{16}\sigma_0^2=\frac34\sigma_0^2$，一致。这里只算了 $30^\circ$ 这一个截面；换一个角度，法向与切向的比例会变，怎样找最大剪应力是 P10 的内容。

\begin{center}
\begin{tikzpicture}[scale=1.65,>={Stealth[length=1.7mm]}]
  % 左：试件与斜截面
  \begin{scope}
    \fill[bandcream] (0,-0.6) rectangle (2.4,0.6);
    \draw[violet,thick] (0,-0.6) rectangle (2.4,0.6);
    \draw[very thick,violet!60!black] (0.9,0.5196)--(1.5,-0.5196);
    \draw[vv] (1.2,0)--(1.7369,0.31) node[above right]{\scriptsize $\vect n$};
    \draw[vv] (1.2,0)--(2.25,0) node[below right]{\scriptsize $\vect t=\ten\sigma\cdot\vect n$};
    \draw[vv] (2.4,0)--(3.0,0) node[right]{\scriptsize $+\sigma_0\vect e_1$};
    \draw[vv] (0,0)--(-0.6,0) node[left]{\scriptsize $-\sigma_0\vect e_1$};
    \node at (1.2,-0.95) {\scriptsize 单轴拉伸：两端向外拉};
  \end{scope}
  % 右：t 沿 n（30°）与截面方向分解
  \begin{scope}[xshift=6.1cm]
    \draw[ax] (-0.26,-0.15)--(1.55,0.895) node[above right]{\scriptsize $\vect n$};
    \draw[vv] (0,0)--(1.55,0) node[below right]{\scriptsize $\vect t$};
    \draw[vviolet] (0,0)--(1.1626,0.6712);
    \draw[inkgray,very thick,->,>={Stealth[length=1.7mm]}] (0,0)--(0.3875,-0.6712);
    \draw[inkgray,densely dotted] (1.1626,0.6712)--(1.55,0);
    \draw[inkgray,densely dotted] (0.3875,-0.6712)--(1.55,0);
    \draw[inkgray,thin] (0.1039,0.0600)--(0.1639,-0.0439)--(0.0600,-0.1039);
    \node[violet,above left] at (1.08,0.70) {\scriptsize $\tfrac34\sigma_0$};
    \node[inkgray,below right] at (0.36,-0.72) {\scriptsize $\tfrac{\sqrt3}{4}\sigma_0$};
    \draw[->] (0.45,0) arc (0:30:0.45);
    \node at (0.60,0.20) {\scriptsize $30^\circ$};
  \end{scope}
\end{tikzpicture}\\[2pt]
\figcap{\mbox{\textbf{图 11}（自绘，示意；各向量长度不按同一比例）：}单轴拉伸 $\ten\sigma=\sigma_0\vect e_1\otimes\vect e_1$。左：右端面的外法向是 $+\vect e_1$，面力 $\vect t(+\vect e_1)=\ten\sigma\cdot\vect e_1=+\sigma_0\vect e_1$；左端面的外法向是 $-\vect e_1$，面力 $\vect t(-\vect e_1)=\ten\sigma\cdot(-\vect e_1)=-\sigma_0\vect e_1$。两个向量都指向试件外侧，所以是“拉”而不是“压”。$30^\circ$ 截面（法向 $\vect n$）上的面力 $\vect t=\ten\sigma\cdot\vect n$ 仍沿 $+\vect e_1$。右：把这个 $\vect t$ 分成法向 $\tfrac34\sigma_0$ 与切向 $\tfrac{\sqrt3}{4}\sigma_0$ 两个分量；灰色虚线是 $\vect n$ 方向的参考轴（$30^\circ$），不是原来的 $x$ 轴。\textbf{斜面上原来看不出有剪应力，换个截面就出来了}——这就是 $\ten\sigma\cdot\vect n$ 的全部内容。}
\end{center}

\subhdr{6.4\quad 平衡方程：应力变化怎样变成净力}
\label{subsec:equilibrium}

\textbf{偏导在做什么。}应力张量在物体里逐点不同，是一个\textbf{张量场} $\ten\sigma(\vect x)$：每个点给一个 $3\times3$ 表。$\partial_j=\partial/\partial x_j$ 的意思是：\textbf{只让第 $j$ 个坐标变，另外两个坐标固定}，看分量 $\sigma_{ij}$ 随它怎么变。例如 $\partial_1\sigma_{11}$ 是“沿 $x_1$ 方向走一小步，$\sigma_{11}$ 增加多少除以这一步”，$x_2,x_3$ 不动。

\textbf{为什么它的和是净面力。}取一个边长为 $h$ 的小立方体，考虑法向为 $\pm\vect e_j$ 的一对面上、沿 $\vect e_i$ 方向的力。$+\vect e_j$ 面的面力是 $\sigma_{ij}(\vect x+h\vect e_j)$，面积 $h^2$；$-\vect e_j$ 面的面力是 $-\sigma_{ij}(\vect x)$，面积也是 $h^2$。两块加起来，沿 $\vect e_i$ 的净力是
\[
\bigl[\sigma_{ij}(\vect x+h\vect e_j)-\sigma_{ij}(\vect x)\bigr]h^2
\;\approx\;h^3\,\partial_j\sigma_{ij}.
\]
除以体积 $h^3$，就得到\textbf{单位体积的净面力}的第 $i$ 个分量 $\partial_j\sigma_{ij}$。对 $j=1,2,3$ 求和，就是 $\nabla\cdot\ten\sigma$。所以“散度”不是又一个符号，而是“三个方向各查一次应力差，再相加”。

\textbf{一个能算的例子。}设
\[
\ten\sigma(x_1)=\frac{\sigma_0 x_1}{L}\,\vect e_1\otimes\vect e_1
\qquad(\sigma_0>0,\;L>0),
\]
即只有 $\sigma_{11}=\sigma_0x_1/L$，沿 $x_1$ 线性增大。则
\[
(\nabla\cdot\ten\sigma)_1=\partial_1\sigma_{11}=\frac{\sigma_0}{L},
\qquad
(\nabla\cdot\ten\sigma)_2=(\nabla\cdot\ten\sigma)_3=0 .
\]
也就是说，每单位体积有一个沿 $+\vect e_1$ 的净面力 $\sigma_0/L$：$+x_1$ 那一侧的面拉得比 $-x_1$ 那一侧更用力。要在静力学里平衡它，需要体力
\[
\vect f=-\nabla\cdot\ten\sigma=\Bigl(-\frac{\sigma_0}{L},0,0\Bigr),
\]
即一个沿 $-\vect e_1$ 的均匀体力。反过来，若整个应力分量处处相同，即 $\sigma_{11}(\vect x)\equiv\sigma_{\mathrm c}$（$\sigma_{\mathrm c}$ 是常数，不再是上面那个随 $x_1$ 线性增大的场），则 $\partial_1\sigma_{11}=0$，内部没有净面力，只有端面上的外力——这就是普通的均匀单轴拉伸。

\textbf{平衡方程。}把面力沿闭合曲面加起来，用高斯定理换成体积分。对张量场 $\ten\sigma$，分量形式是
\[
\oint_S\sigma_{ij}n_j\dd S=\int_V\partial_j\sigma_{ij}\dd V,
\qquad
(\nabla\cdot\ten\sigma)_i=\partial_j\sigma_{ij},
\]
其中 $\vect n$ 是闭合曲面的外法向，$\partial_j=\partial/\partial x_j$。与体力 $\vect f$ 平衡得（体力是作用在体积上的力，如重力、惯性力，$\vect f$ 是单位体积的力）
\[
\nabla\cdot\ten\sigma+\vect f=\vect0\quad\text{（静力学）},
\qquad
\nabla\cdot\ten\sigma+\vect f=\rho\vect a\quad\text{（动力学）},
\]
其中 $\rho$ 是质量密度、$\vect a$ 是加速度。写得开就是三个方向的方程：
\begin{align*}
\partial_1\sigma_{11}+\partial_2\sigma_{12}+\partial_3\sigma_{13}+f_1&=0,\\
\partial_1\sigma_{21}+\partial_2\sigma_{22}+\partial_3\sigma_{23}+f_2&=0,\\
\partial_1\sigma_{31}+\partial_2\sigma_{32}+\partial_3\sigma_{33}+f_3&=0.
\end{align*}
每个方程只含一个方向的力分量，所以叫“三个方向的力平衡”。

\textbf{高斯定理在这里做什么。}标量形式是 $\oint_S f n_j\dd S=\int_V\partial_j f\dd V$；把 $f$ 换成 $\sigma_{ij}$（固定 $i$）逐分量使用即可。一维类比是 $\int_a^b f'(x)\dd x=f(b)-f(a)$：内部点成对抵消，只剩端点。严格证明要用中值定理控制“近似”，本稿不展开；这里只需要知道它把“面上的力”换成了“体内每一点的散度”。

\warnbox{\textbf{指标顺序要盯住。}按 \S6.1 的约定，$\sigma_{ij}$ 的第一个指标是\textbf{力的方向}，第二个是\textbf{面的法向}。面力是 $t_i=\sigma_{ij}n_j$，所以高斯定理缩并的是\textbf{第二个指标}：$\partial_j\sigma_{ij}$。若 $\ten\sigma$ 对称，$\partial_j\sigma_{ji}$ 与它逐分量相等；若 $\ten\sigma$ 不对称（例如有体力矩的介质），两者一般\textbf{不同}，而且不是简单差一个整体符号——它们是两个不同的缩并，不能互相替换。$\nabla\cdot\ten\sigma$ 是\textbf{矢量}（缩并一次），$\nabla\ten\sigma$ 才是三阶张量；两者差一个点，差一个阶。}

\hdr{§7}{独立练习}
\label{sec:exercises}

下面十二道题覆盖本稿的全部操作：并矢运算与不变量（P1）、惯量张量（P2、P4）、欧拉方程与对称陀螺（P3、P12）、并矢的左右作用与张量判定（P5、P6）、迹与对称化（P7、P8）、惯量的构造（P9）、应力（P10）、旋转算子（P11）。
\textbf{用法：}先自己在纸上做，不看附录；做完再翻附录 A 对详解。正文只给题干与考点，不给关键步骤或结果。P4 用到 P2 的 $\ten J$，若还没做 P2，请先做 P2（或按 P2 的公式自己算一遍）。

\practice{P1}{自编题，承旧稿例 1}
{设 $\ten D=4\vect e_1\otimes\vect e_1+\vect e_1\otimes\vect e_2+\vect e_2\otimes\vect e_1+4\vect e_2\otimes\vect e_2+2\vect e_3\otimes\vect e_3$，$\vect x=\vect e_1+2\vect e_2+3\vect e_3$。求 $\ten D\cdot\vect x$、$\tr\ten D$、$\det\ten D$、$\ten D:\ten D$、$\ten D$ 的特征值与一组主轴，并写出谱分解。}
{并矢的分量、右作用、迹/行列式/双点积，以及对称张量的特征值与主轴。}{ans:P1}

\practice{P2}{自编题，承旧稿例 2}
{均匀长方体质量 $m=5$ kg，边长 $a=2$、$b=3$、$c=4$ m，质心在原点，棱沿坐标轴。求它的转动惯量张量 $\ten J$；再把它压成薄板（$c\to0$，质量仍为 $m$），验证垂直轴定理。}
{长方体惯量张量公式、薄板极限与垂直轴定理。}{ans:P2}

\practice{P3}{自编题，承旧稿例 3}
{对称陀螺 $A=B=2$ kg·m$^2$、$C=3$ kg·m$^2$，自转 $\omega_3=10$ rad/s，无外力矩。求横向分量绕体轴的角速度与周期，并说明转向。}
{欧拉方程在 $A=B$ 时的解、进动角速度与周期。}{ans:P3}

\practice{P4}{自编题，承旧稿例 4}
{用 P2 得到的 $\ten J$。求沿与三条棱成等角的方向 $\vect n=(1,1,1)/\sqrt3$ 的转动惯量 $I(\vect n)=\vect n\cdot\ten J\cdot\vect n$（这个方向不是 P2 长方体的空间对角线）；再取 $\boldsymbol\omega=\vect n$（$|\boldsymbol\omega|=1$），验证 $\boldsymbol\omega\cdot\vect L=2K$。}
{二次型沿方向的值、角动量与转动动能的关系。}{ans:P4}

\practice{P5}{自编概念题，承旧稿附录 B 第 1 问}
{判断并说明：(a) $(\vect a\otimes\vect b)\cdot\vect c$ 与 $\vect a(\vect b\cdot\vect c)$ 是否相同？(b) $(\vect a\otimes\vect b)\cdot\vect c$ 与 $(\vect a\cdot\vect b)\vect c$ 是否相同？}
{并矢右作用与两种“先点乘再乘矢量”的区别。}{ans:P5}

\practice{P6}{自编概念题，承旧稿附录 B 第 2 问}
{给出一组九个排成方阵的数，怎样判断它是不是二阶张量？举一个“在旋转下像矢量、但坐标系允许平移时就不像矢量”的例子。}
{张量的变换律与“对象/账本”的区别。}{ans:P6}

\practice{P7}{自编概念题，承旧稿附录 B 第 3 问}
{设 $\ten R$ 是 $3\times3$ 正交矩阵，即 $\ten R^{\mathsf T}\ten R=\ten R\ten R^{\mathsf T}=\unitd$；$\ten T$ 是任意 $3\times3$ 实矩阵。用 $\tr(\ten A\ten B)=\tr(\ten B\ten A)$ 写出两步，证明 $\tr(\ten R^{\mathsf T}\ten T\ten R)=\tr\ten T$，并说明每一步用了什么性质。}
{迹的循环性与正交换基下的不变性。}{ans:P7}

\practice{P8}{自编概念题，承旧稿附录 B 第 4 问}
{二次型 $\vect x\cdot\ten D\cdot\vect x$ 为什么只与 $\ten D$ 的对称部分有关？举一个 $2\times2$ 的数字例子验证。}
{指标重命名、对称化与二次型。}{ans:P8}

\practice{P9}{自编概念题，承旧稿附录 B 第 5 问}
{$J_{12}=-\int xy\dd m$ 里的负号从哪里来？如果把它丢掉，矩阵还对称吗？谱定理还成立吗？算出的主转动惯量还是真实的吗？}
{惯量张量的构造、对称性与主值的物理意义。}{ans:P9}

\practice{P10}{自编概念题，承旧稿附录 B 第 6 问}
{单轴拉伸 $\ten\sigma=\sigma_0\vect e_1\otimes\vect e_1$。最大剪应力多大？出现在哪个截面？为什么这个结果不需要背？}
{面力 $\vect t=\ten\sigma\cdot\vect n$、斜截面分解与最大剪应力。}{ans:P10}

\practice{P11}{自编题}
{设 $S(t)$ 是绕 $\vect e_3$ 轴转 $\varphi(t)=\alpha t$ 的旋转算子（$\alpha$ 为常数）。(a) 写出 $S$ 在 $\{\vect e_1,\vect e_2,\vect e_3\}$ 下的矩阵；(b) 验证 $S^{\mathsf T}S=\unitd$；(c) 用 $\boldsymbol\Omega^{\times}=(\dd S/\dd t)S^{\mathsf T}$ 求 $\boldsymbol\Omega$。}
{旋转算子的矩阵、正交性与角速度。}{ans:P11}

\practice{P12}{自编题}
{刚体的三个主转动惯量为 $A=1$、$B=2$、$C=3$ kg·m$^2$。某瞬时在体坐标系中 $\boldsymbol\omega=(1,1,1)$ rad/s，外力矩 $\vect M=\vect0$。用欧拉方程求此刻 $\dot\omega_1,\dot\omega_2,\dot\omega_3$。}
{欧拉方程的分量形式与陀螺项。}{ans:P12}

\hdr{§8}{易错点、速查与三句话}
\label{sec:review}

\warnbox{\textbf{八个高频错误}\\[2pt]
1. \textbf{把张量当成“九个数字”}：九个数字要按变换律 $T'_{ij}=R_{ki}R_{lj}T_{kl}$ 变才叫张量（\S2.7）。不检验就套用谱定理，是错的源头。\par
2. \textbf{把 $\vect a\otimes\vect b$ 当成 $\vect a\cdot\vect b$}：前者是算子（九个分量），后者是标量。中间是 $\otimes$ 还是 $\cdot$，差一个层级。\par
3. \textbf{左右的顺序混用}：作为算子，$\ten D\cdot\vect x$ 与 $\vect x\cdot\ten D$ 一般不同；它们对\textbf{所有} $\vect x$ 相同，当且仅当 $\ten D$ 对称（\S2.3）。力学里 $\ten J$、$\ten\sigma$ 对称，所以“碰巧”怎么写都对；换到一般旋转算子 $S$ 上，左右一混通常出错。\par
4. \textbf{直接写 $\dd\vect L/\dd t=\ten J\cdot\dot{\boldsymbol\omega}$}：只在 $\ten J$ 于该系为常量时成立；体坐标系里求的是体导数，固定系里还要补 $\boldsymbol\omega\times\vect L$（\S5.1）。\par
5. \textbf{把惯量张量的非对角元符号写反}：$J_{12}=-\int xy\dd m$，负号不能丢；丢了以后主轴会算错（P9）。\par
6. \textbf{把 $\nabla\cdot\ten\sigma$ 的指标缩并错}：按“力方向在前、法向在后”的约定，面力是 $t_i=\sigma_{ij}n_j$，所以散度是 $\partial_j\sigma_{ij}$；对称时与 $\partial_j\sigma_{ji}$ 相同，非对称时不同（\S6.4）。\par
7. \textbf{把特征值当成“唯一的解”}：三个特征值不变，但\textbf{特征向量的分量随基变}（\S1.6）。说“主轴方向”必须连着坐标系或刚体说；重根时还要说明在哪个特征空间里选。\par
8. \textbf{把换基矩阵的转置写反}：$\ten R$ 的列是新基在旧基下的分量时，矢量的新坐标是 $\vect x'=\ten R^{\mathsf T}\vect x$，二阶张量是 $\ten T'=\ten R^{\mathsf T}\ten T\ten R$（\S0.4）；不是 $\ten R\vect x$。}

\begin{center}
\tabtop
\begin{tabular}{@{}p{0.32\linewidth}p{0.58\linewidth}@{}}
\thead{要做的事} & \thead{用什么} \\
\multicolumn{2}{@{}l@{}}{\color{petal}\rule{\linewidth}{0.5pt}}\\
判断是不是张量 & 查变换律 $T'_{ij}=R_{ki}R_{lj}T_{kl}$；或对已知矢量的线性关系用商法则（\S2.7）\\[3pt]
把“矢量 $\to$ 矢量”写成式子 & $\ten D=D_{ij}\vect e_i\otimes\vect e_j$，作用规则 $(\vect a\otimes\vect b)\cdot\vect x=\vect a(\vect b\cdot\vect x)$\\[3pt]
判断是否对称 & 看 $D_{ij}=D_{ji}$；反对称部分等价于一个轴矢量（叉乘，\S2.6）\\[3pt]
求主轴与主值 & 解 $\det(\ten D-\lambda\unitd)=0$；对称时主轴正交，$\ten D=\sum_k\lambda_k\vect n_k\otimes\vect n_k$（\S1.4、\S2.5）\\[3pt]
验算 & $I_1=\tr$、$I_2$、$I_3=\det$；惯量张量还有 $A+B\ge C$（\S1.6、\S4.3）\\[3pt]
转动惯量 & $\ten J=\int(r^2\unitd-\vect r\otimes\vect r)\dd m$；平行轴 $\ten J_O=\ten J_C+m(R^2\unitd-\vect R\otimes\vect R)$（\S4.1、\S4.5）\\[3pt]
角动量与动能 & $\vect L=\ten J\cdot\boldsymbol\omega$，$K=\tfrac12\boldsymbol\omega\cdot\ten J\cdot\boldsymbol\omega$，$\boldsymbol\omega\cdot\vect L=2K$（\S4.2）\\[3pt]
体坐标系动力学 & $(\dd\vect L/\dd t)_{\text{体}}+\boldsymbol\omega\times\vect L=\vect M$，即欧拉方程（\S5.2）\\[3pt]
面力 & $\vect t(\vect n)=\ten\sigma\cdot\vect n$；$\sigma_{ij}=\sigma_{ji}$；主应力 = 特征值（\S6.1--\S6.2）\\
\end{tabular}
\tabbot
\end{center}

\flowerrule

\subhdr{三句话总结}
\begin{enumerate}[label=\arabic*.,leftmargin=1.6em]
\item \textbf{张量是对象，矩阵是账本。}换基时账本变、对象不变；迹、行列式、特征值是账本里不随基变的数。
\item \textbf{并矢只是二阶张量的另一种写法。}作用规则 $(\vect a\otimes\vect b)\cdot\vect x=\vect a(\vect b\cdot\vect x)$ 一条就够；左右作用作为算子一般不对称，只有张量对称时才对所有输入相同；双点积吃掉两个指标、谱分解是三个投影并矢相加。
\item \textbf{理论力学里三处落点用的是同一台机器。}惯量张量给出 $\vect L=\ten J\cdot\boldsymbol\omega$，欧拉方程是它在体坐标系里的求导，应力张量是“面力对法向线性”的商法则产物。
\end{enumerate}

\textbf{最短主线。}如果时间有限，按 \S0.1--\S0.7 $\to$ \S1.1--\S1.5 $\to$ \S2.1--\S2.7 $\to$ \S3.1--\S3.6 $\to$ \S4.1--\S4.5 $\to$ \S5.1--\S5.3 $\to$ \S6.1--\S6.4 读；这些段落里没有把后面要用的定义藏进选读。标“选读”的段落跳过以后，主线仍然成立。

\clearpage
\appendix
\hdr{附录 A}{练习详解}
\label{app:solutions}

\answerstart{P1}{自编题，承旧稿例 1}
{设 $\ten D=4\vect e_1\otimes\vect e_1+\vect e_1\otimes\vect e_2+\vect e_2\otimes\vect e_1+4\vect e_2\otimes\vect e_2+2\vect e_3\otimes\vect e_3$，$\vect x=\vect e_1+2\vect e_2+3\vect e_3$。求 $\ten D\vect x$、$\tr\ten D$、$\det\ten D$、$\ten D:\ten D$、特征值与一组主轴，并写出谱分解。}

先把九个系数按行列摆好（第 $i$ 行第 $j$ 列是 $\vect e_i\otimes\vect e_j$ 的系数）：
\[
\ten D=\begin{pmatrix}4&1&0\\1&4&0\\0&0&2\end{pmatrix}.
\]
作用：$(\ten D\vect x)_i=\sum_jD_{ij}x_j$。第一行 $4\cdot1+1\cdot2+0=6$，第二行 $1\cdot1+4\cdot2=9$，第三行 $0+0+2\cdot3=6$，所以
\[
\ten D\vect x=(6,9,6).
\]
因为 $\ten D$ 对称，左作用 $\vect x\cdot\ten D$ 给同一个结果（\S2.3）。

迹与行列式：$\tr\ten D=4+4+2=10$；按第三行展开，$\det\ten D=2(16-1)=30$。
双点积：$\ten D:\ten D=\sum_{i,j}D_{ij}^2=16+1+1+16+4=38$。

特征方程由 \S0.5：$-\lambda^3+10\lambda^2-I_2\lambda+30=0$，其中 $I_2=\frac12(10^2-38)=31$。即
\[
\lambda^3-10\lambda^2+31\lambda-30=0 .
\]
试根 $\lambda=2$：$8-40+62-30=0$，成立。分解
\[
(\lambda-2)(\lambda^2-8\lambda+15)=(\lambda-2)(\lambda-3)(\lambda-5)=0,
\]
所以特征值是 $2,3,5$。

主轴：$\lambda=5$ 时
\[
(\ten D-5\unitd)=\begin{pmatrix}-1&1&0\\1&-1&0\\0&0&-3\end{pmatrix},
\]
前两行都给出 $-x+y=0$，即 $x=y$；第三行给 $-3z=0$，即 $z=0$。所以方向是 $(1,1,0)$，单位化 $\vect n_1=(\vect e_1+\vect e_2)/\sqrt2$。代回核验：
\[
(\ten D-5\unitd)\begin{pmatrix}1\\1\\0\end{pmatrix}
=\begin{pmatrix}-1+1\\1-1\\0\end{pmatrix}=\vect0 .
\]
$\lambda=3$ 时 $(\ten D-3\unitd)$ 的前两行给 $x+y=0$，取 $\vect n_2=(\vect e_1-\vect e_2)/\sqrt2$（也可以取相反方向，差一个符号）；$\lambda=2$ 时前两行给 $2x+y=0$ 与 $x+2y=0$，只能 $x=y=0$，第三行不限制 $z$，所以 $\vect n_3=\vect e_3$。谱分解
\[
\ten D=5\,\vect n_1\otimes\vect n_1+3\,\vect n_2\otimes\vect n_2+2\,\vect e_3\otimes\vect e_3 .
\]
验算（\S1.6）：$5+3+2=10=\tr$，$5\cdot3\cdot2=30=\det$，$25+9+4=38=\ten D:\ten D$，都对。

\answerstart{P2}{自编题，承旧稿例 2}
{均匀长方体 $m=5$ kg，$a=2,b=3,c=4$ m，质心在原点，棱沿坐标轴。求 $\ten J$；再压成薄板（$c\to0$，质量仍为 $m$），验证垂直轴定理。}

由 \S4.4 直接代入：
\[
\begin{aligned}
\ten J&=\frac{m}{12}\operatorname{diag}(b^2+c^2,\;a^2+c^2,\;a^2+b^2)\\
&=\frac{5}{12}\operatorname{diag}(25,\;20,\;13)
=\operatorname{diag}\Bigl(\frac{125}{12},\frac{100}{12},\frac{65}{12}\Bigr)\\
&\approx\operatorname{diag}(10.417,\;8.333,\;5.417)\ \text{kg·m}^2 .
\end{aligned}
\]
验算三角形不等式：$10.417+8.333=18.750\ge5.417$；$10.417+5.417=15.834\ge8.333$；$8.333+5.417=13.750\ge10.417$，三个循环都成立（\S4.3）。

压成薄板（$c=0$，质量仍为 $5$ kg）后
\[
\ten J_{\text{板}}=\frac{m}{12}\operatorname{diag}(b^2,\;a^2,\;a^2+b^2)
=\frac{5}{12}\operatorname{diag}(9,\;4,\;13)
=\operatorname{diag}(3.750,\;1.667,\;5.417).
\]
于是 $J_{zz}=65/12$，而 $J_{xx}+J_{yy}=(45+20)/12=65/12$，二者相等，垂直轴定理成立。

\answerstart{P3}{自编题，承旧稿例 3}
{对称陀螺 $A=B=2$ kg·m$^2$、$C=3$ kg·m$^2$，自转 $\omega_3=10$ rad/s，无外力矩。求横向分量绕体轴的角速度与周期，并说明转向。}

由 \S5.3，
\[
\Omega=\frac{C-A}{A}\omega_3=\frac{3-2}{2}\times10=5\ \text{rad/s},
\qquad
T=\frac{2\pi}{\Omega}=\frac{2\pi}{5}\approx1.257\ \text{s}.
\]
即 $\boldsymbol\omega$ 的横向分量每 $1.257$ 秒绕体轴转一圈，$\omega_3$ 与 $|\boldsymbol\omega|$ 都不变。因为 $C>A$，$\Omega>0$，所以横向分量沿 $+\vect e_3$ 的正向（右手定则）转；在体坐标系里看，$\boldsymbol\omega$ 的端点绕 $\vect e_3$ 画圆，$\vect L$ 与 $\boldsymbol\omega$、$\vect e_3$ 始终共面（\S5.3 的图 8）。

\answerstart{P4}{自编题，承旧稿例 4}
{用 P2 的 $\ten J$。求沿与三条棱成等角的方向 $\vect n=(1,1,1)/\sqrt3$ 的转动惯量 $I(\vect n)=\vect n\cdot\ten J\cdot\vect n$（该方向不是 P2 长方体的空间对角线 $(2,3,4)/\sqrt{29}$）；再取 $\boldsymbol\omega=\vect n$，验证 $\boldsymbol\omega\cdot\vect L=2K$。}

P2 已得
\[
\ten J=\operatorname{diag}\Bigl(\frac{125}{12},\frac{100}{12},\frac{65}{12}\Bigr).
\]
沿单位方向 $\vect n$ 的转动惯量是二次型在 $\vect n$ 上的值（\S1.5 的瑞利商）：
\[
I(\vect n)=\vect n\cdot\ten J\vect n
=\frac{1}{3}\Bigl(\frac{125}{12}+\frac{100}{12}+\frac{65}{12}\Bigr)
=\frac{1}{3}\cdot\frac{290}{12}
=\frac{145}{18}\approx8.056\ \text{kg·m}^2 .
\]
它落在最小主值 $65/12\approx5.417$ 与最大主值 $125/12\approx10.417$ 之间，符合 $\lambda_{\min}\le I\le\lambda_{\max}$。

验证：取 $\boldsymbol\omega=\vect n$（$|\boldsymbol\omega|=1$），则 $K=\frac12\boldsymbol\omega\cdot\ten J\boldsymbol\omega=\frac12I=145/36$；
而
\[
\vect L=\ten J\boldsymbol\omega=\frac{1}{\sqrt3}\Bigl(\frac{125}{12},\frac{100}{12},\frac{65}{12}\Bigr),
\qquad
\boldsymbol\omega\cdot\vect L=\frac{1}{3}\cdot\frac{290}{12}=\frac{145}{18}=2K .
\]
最后这个等式是恒等式（不依赖是否在主轴系）：$\boldsymbol\omega\cdot(\ten J\boldsymbol\omega)=\boldsymbol\omega\cdot\ten J\cdot\boldsymbol\omega=2K$。

\answerstart{P5}{自编概念题，承旧稿附录 B 第 1 问}
{判断并说明：(a) $(\vect a\otimes\vect b)\cdot\vect c$ 与 $\vect a(\vect b\cdot\vect c)$ 是否相同？(b) $(\vect a\otimes\vect b)\cdot\vect c$ 与 $(\vect a\cdot\vect b)\vect c$ 是否相同？}

(a) 相同。按右作用规则 $(\vect a\otimes\vect b)\cdot\vect c=\vect a(\vect b\cdot\vect c)$；读法是“先把 $\vect c$ 投影到 $\vect b$ 上得到数 $\vect b\cdot\vect c$，再沿 $\vect a$ 放大”。

(b) 一般不同。左边是“投到 $\vect b$、沿 $\vect a$ 放大”；右边是“先算 $\vect a\cdot\vect b$ 这个数，再沿 $\vect c$ 放大”。取 $\vect a=\vect e_1$、$\vect b=\vect e_1$、$\vect c=\vect e_2$：
\[
(\vect e_1\otimes\vect e_1)\cdot\vect e_2=\vect e_1(\vect e_1\cdot\vect e_2)=\vect0,
\qquad
(\vect e_1\cdot\vect e_1)\vect e_2=\vect e_2 .
\]
一个为零，一个不为零，所以不相同。两者只有在特殊输入下才碰巧相等。例如让 $\vect c$ 沿 $\vect a$，即 $\vect c=k\vect a$：
\[
(\vect a\otimes\vect b)\cdot(k\vect a)=k\vect a(\vect b\cdot\vect a),
\qquad
(\vect a\cdot\vect b)(k\vect a)=k\vect a(\vect a\cdot\vect b),
\]
两边相同（数乘可以交换）。这说明“一般不同、特殊输入可以相同”，但不能因此把两个表达式当成同一个。

\answerstart{P6}{自编概念题，承旧稿附录 B 第 2 问}
{给出一组九个排成方阵的数，怎样判断它是不是二阶张量？举一个“在旋转下像矢量、但坐标系允许平移时就不像矢量”的例子。}

判据是\textbf{变换律}：对任意正交换基（$R_{ki}=\vect n_i\cdot\vect e_k$ 是新的正交基），检查
\[
T'_{ij}=R_{ki}R_{lj}T_{kl}
\]
是否成立。成立才是二阶张量；不成立只是九个数字。实际操作中，若已知一个物理关系 $\vect u=\ten T\vect v$ 对任意矢量 $\vect v$ 成立，且 $\vect u,\vect v$ 都是矢量，则可用商法则（\S2.7）直接判定 $\ten T$ 是张量，不必逐项验证。

例子：$\delta_{ij}$ 与 $a_ib_j$ 都是张量，前者已在 \S2.7 验过，后者就是并矢（\S2.2）。\textbf{位置矢量} $(x_1,x_2,x_3)$ 在旋转下按 $\vect x'=\ten R^{\mathsf T}\vect x$ 变，所以对旋转来说它是矢量；但若允许坐标系平移 $\vect x\to\vect x+\vect a$，分量会多出一个常数项，不再是齐次的，所以对“平移 + 旋转”这个更大的群它不是矢量。\textbf{“是不是张量”必须连着允许哪些变换说}；本稿只允许旋转与反射，在这个群里位置矢量是矢量。

\answerstart{P7}{自编概念题，承旧稿附录 B 第 3 问}
{设 $\ten R$ 是 $3\times3$ 正交矩阵，即 $\ten R^{\mathsf T}\ten R=\ten R\ten R^{\mathsf T}=\unitd$；$\ten T$ 是任意 $3\times3$ 实矩阵。用 $\tr(\ten A\ten B)=\tr(\ten B\ten A)$ 写出两步，证明 $\tr(\ten R^{\mathsf T}\ten T\ten R)=\tr\ten T$，并说明每一步用了什么性质。}

\textbf{先把用到的两个工具说清楚。}第一个是迹的定义：对 $3\times3$ 矩阵 $\ten A$，
\[
\tr\ten A=A_{11}+A_{22}+A_{33}=\sum_{i=1}^{3}A_{ii}.
\]
第二个是迹的循环性。题目已经把这条性质作为工具给出；这里把它写清楚，是为了让第一步有依据：
\[
\tr(\ten A\ten B)=\tr(\ten B\ten A)
\]
对任意两个同阶方阵 $\ten A,\ten B$ 成立。证明是一次指标求和：
\[
\tr(\ten A\ten B)
=\sum_{i=1}^{3}(\ten A\ten B)_{ii}
=\sum_{i=1}^{3}\sum_{j=1}^{3}A_{ij}B_{ji}
=\sum_{j=1}^{3}\sum_{i=1}^{3}B_{ji}A_{ij}
=\sum_{j=1}^{3}(\ten B\ten A)_{jj}
=\tr(\ten B\ten A).
\]
\clearpage
\answercontinue{P7}{P7 的迹循环性已写出；本页接续把 $\ten R^{\mathsf T}\ten T$ 看作整体，用循环性与正交性完成证明，并说明每一步用的性质。}
这里只交换了两个有限和的次序，并对同一个乘积换了一种分组方式；没有用到对称、可逆或正交。$\ten A,\ten B$ 只要是同阶方阵即可。

\textbf{现在证明题目。}把 $\ten R^{\mathsf T}\ten T$ 看成一个整体，记作 $\ten A$，把 $\ten R$ 记作 $\ten B$。第一步用循环性：
\[
\tr\bigl((\ten R^{\mathsf T}\ten T)\ten R\bigr)
=\tr\bigl(\ten R(\ten R^{\mathsf T}\ten T)\bigr).
\]
这一步只用了 $\tr(\ten A\ten B)=\tr(\ten B\ten A)$，其中 $\ten A=\ten R^{\mathsf T}\ten T$，$\ten B=\ten R$。

第二步把右边的括号重新结合，并用正交性：
\[
\tr\bigl(\ten R(\ten R^{\mathsf T}\ten T)\bigr)
=\tr\bigl((\ten R\ten R^{\mathsf T})\ten T\bigr)
=\tr(\unitd\ten T)
=\tr\ten T .
\]
第一个等号用了矩阵乘法的结合律，把 $\ten R$ 与 $\ten R^{\mathsf T}$ 先乘起来；第二个等号用了正交性 $\ten R\ten R^{\mathsf T}=\unitd$；第三个等号用了单位矩阵 $\unitd\ten T=\ten T$。把两步接起来就是
\[
\tr(\ten R^{\mathsf T}\ten T\ten R)
=\tr\bigl((\ten R^{\mathsf T}\ten T)\ten R\bigr)
=\tr\bigl(\ten R(\ten R^{\mathsf T}\ten T)\bigr)
=\tr\bigl((\ten R\ten R^{\mathsf T})\ten T\bigr)
=\tr\ten T .
\]

\textbf{每一步用了什么性质。}第一步用“迹的循环性”；随后用“矩阵乘法的结合律”；再用“$\ten R$ 正交，即 $\ten R\ten R^{\mathsf T}=\unitd$”；最后用“单位矩阵乘任何矩阵不变”。题目只要求“写出两步”，核心就是“循环一次，再把 $\ten R\ten R^{\mathsf T}$ 结合出来”；中间多写的等号只是把结合律写细，不改变逻辑。

\textbf{条件与含义。}这里 $\ten R$ 必须是方阵且正交；$\ten T$ 不必对称，也不必可逆。结论说明：若 $\ten T'=\ten R^{\mathsf T}\ten T\ten R$ 是同一线性映射在另一组正交归一基下的矩阵，那么 $\tr\ten T'=\tr\ten T$；也就是说，\textbf{迹在正交换基下不变}。但要注意，它不说明 $\ten T$ 的每个矩阵元都不变——换基时矩阵元一般会变，只有迹这样的组合不变。当 $\ten T$ 又是实对称矩阵、$\ten R$ 取 \S1.5 的主轴矩阵时，把 $\ten R^{\mathsf T}\ten T\ten R=\operatorname{diag}(\lambda_1,\lambda_2,\lambda_3)$ 代入，还能得到
\[
\tr\ten T=\lambda_1+\lambda_2+\lambda_3,
\]
这与 \S1.3 用迹抽查特征值的做法一致。

\answerstart{P8}{自编概念题，承旧稿附录 B 第 4 问}
{二次型 $\vect x\cdot\ten D\cdot\vect x$ 为什么只与 $\ten D$ 的对称部分有关？举一个 $2\times2$ 的数字例子验证。}

把二次型写成 $\sum_{i,j}D_{ij}x_ix_j$。交换 $i,j$ 的名字（只是换字母，值不变）：
\[
\sum_{i,j}D_{ij}x_ix_j=\sum_{i,j}D_{ji}x_jx_i=\sum_{i,j}D_{ji}x_ix_j .
\]
两式相加除二，得到 $\sum_{i,j}\frac12(D_{ij}+D_{ji})x_ix_j$，只含对称部分 $S_{ij}=\frac12(D_{ij}+D_{ji})$。所以反对称部分对二次型没有贡献。

数字例子：取 $\ten D=\begin{pmatrix}0&3\\0&0\end{pmatrix}$，则
\[
E=\begin{pmatrix}x&y\end{pmatrix}\begin{pmatrix}0&3\\0&0\end{pmatrix}\begin{pmatrix}x\\y\end{pmatrix}=3xy .
\]
\clearpage
\answercontinue{P8}{P8 已用 $\ten D=\begin{pmatrix}0&3\\0&0\end{pmatrix}$ 算到 $E=3xy$；本页接续算对称部分，并比较两个二次型。}
对称部分是 $\ten S=\begin{pmatrix}0&1.5\\1.5&0\end{pmatrix}$，它给出的二次型是
\[
E=\begin{pmatrix}x&y\end{pmatrix}\begin{pmatrix}0&1.5\\1.5&0\end{pmatrix}\begin{pmatrix}x\\y\end{pmatrix}=1.5xy+1.5yx=3xy,
\]
一样。注意原矩阵不对称，但二次型只“看见”它的对称部分。

\answerstart{P9}{自编概念题，承旧稿附录 B 第 5 问}
{$J_{12}=-\int xy\dd m$ 里的负号从哪里来？如果把它丢掉，矩阵还对称吗？谱定理还成立吗？算出的主转动惯量还是真实的吗？}

负号来自 \S4.1 的构造：$\ten J=\int(r^2\unitd-\vect r\otimes\vect r)\dd m$，而 $\vect r\otimes\vect r$ 的分量是 $x_ix_j$。对 $i=1,j=2$，$r^2\delta_{12}=0$，所以
\[
J_{12}=\int(0-x_1x_2)\dd m=-\int xy\dd m .
\]
它不是“人为加的”，是 $-\vect r\otimes\vect r$ 那一项的直接结果。物理上它保证“沿 $\vect r$ 方向转时惯量为零”这类事实。

如果只把 $J_{12}=J_{21}$ 的负号去掉（即把它们写成 $+\int xy\dd m$，其余分量不动），矩阵仍然是对称的，\S1.4 的谱定理也仍然成立（对称矩阵总能正交对角化），但主转动惯量会变。举一个让非对角元真正起作用的例子：一个质点 $m$ 在 $\vect r=(1,1,1)$。真实的
\[
\ten J=m(r^2\unitd-\vect r\otimes\vect r)
=m\begin{pmatrix}2&-1&-1\\-1&2&-1\\-1&-1&2\end{pmatrix},
\]
它的三个主值是 $0,3m,3m$：沿 $\vect r$ 方向（质点就在这条直线上）惯量为零，垂直于 $\vect r$ 的任意方向为 $3m$。若把 $J_{12}=J_{21}=-m$ 错写成 $+m$，得到
\[
\ten J'=m\begin{pmatrix}2&1&-1\\1&2&-1\\-1&-1&2\end{pmatrix},
\]
它仍对称、仍可正交对角化。注意 $\ten J'$ 带整体因子 $m$：把它写成 $\ten J'=m\ten A$，其中
\[
\ten A=\begin{pmatrix}2&1&-1\\1&2&-1\\-1&-1&2\end{pmatrix}.
\]
\clearpage
\answercontinue{P9}{P9 已写出真实 $\ten J$、错误 $\ten J'$ 与 $\ten A$ 的矩阵；本页接续由 $\ten A$ 的特征方程求 $\ten J'$ 的主值，并与真实主值比较。}
整体因子只把特征值乘上 $m$、不改变特征向量，所以 $\ten A$ 的特征方程是 $(\lambda-1)^2(\lambda-4)=0$，$\ten J'$ 的主值就是 $m,m,4m$（等价地，直接解 $\det(\ten J'-\lambda\unitd)=0$ 也会得到 $\lambda=m,m,4m$）。对比 $0,3m,3m$，可见“对称”只保证谱定理可用，不保证主值正确。负号是构造式 $-\vect r\otimes\vect r$ 的一部分，丢了它就换了另一个张量。（另一种更粗的错误是把整个构造写成 $r^2\unitd+\vect r\otimes\vect r$；那会连对角元一起改，后果更大。）\textbf{所以“对称”不等于“正确”；物理内容要靠构造式保证。}

\answerstart{P10}{自编概念题，承旧稿附录 B 第 6 问}
{单轴拉伸 $\ten\sigma=\sigma_0\vect e_1\otimes\vect e_1$。最大剪应力多大？出现在哪个截面？为什么这个结果不需要背？}

取法向与 $\vect e_1$ 成 $\theta$ 的截面（在 $\vect e_1\vect e_2$ 面内），$\vect n=\cos\theta\,\vect e_1+\sin\theta\,\vect e_2$。面力
\[
\vect t=\ten\sigma\cdot\vect n=\sigma_0\cos\theta\,\vect e_1 .
\]
法向分量 $t_n=\vect t\cdot\vect n=\sigma_0\cos^2\theta$；切向分量大小
\[
\tau=\sqrt{|\vect t|^2-t_n^2}
=\sigma_0\sin\theta\cos\theta
=\frac{\sigma_0}{2}\sin2\theta .
\]
所以 $\tau$ 在 $\sin2\theta=1$ 时最大，即 $2\theta=90^\circ$、$\theta=45^\circ$，最大值
\[
\tau_{\max}=\frac{\sigma_0}{2}.
\]
出现在法向与 $\vect e_1$ 成 $45^\circ$ 的截面上（例如 $\vect n=(\vect e_1+\vect e_2)/\sqrt2$ 或 $(\vect e_1-\vect e_2)/\sqrt2$；绕 $\vect e_1$ 转一圈的任一这样的截面都一样）。在 $45^\circ$ 截面上，法向应力也正好是 $\sigma_0/2$。

\textbf{为什么不需要背：}只要写出 $\vect t=\ten\sigma\cdot\vect n$，再把它分成法向与切向，$\tau=\sigma_0\sin\theta\cos\theta$ 就出来了；最大值是 $\sin2\theta$ 的最大值，不需要记特殊角度。这个例子也解释了“拉伸试件常沿 $45^\circ$ 断裂”：那个截面上剪应力最大。

\answerstart{P11}{自编题}
{设 $S(t)$ 是绕 $\vect e_3$ 轴转 $\varphi(t)=\alpha t$ 的旋转算子（$\alpha$ 为常数）。(a) 写出 $S$ 的矩阵；(b) 验证 $S^{\mathsf T}S=\unitd$；(c) 用 $\boldsymbol\Omega^{\times}=(\dd S/\dd t)S^{\mathsf T}$ 求 $\boldsymbol\Omega$。}

(a) 由 \S3.2，绕 $\vect e_3$ 转 $\varphi$ 的矩阵是
\[
S(t)=
\begin{pmatrix}
\cos\alpha t&-\sin\alpha t&0\\
\sin\alpha t&\cos\alpha t&0\\
0&0&1
\end{pmatrix}.
\]
它的第 $1$ 列是 $\vect n'_1=(\cos\alpha t,\sin\alpha t,0)$，第 $2$ 列是 $\vect n'_2=(-\sin\alpha t,\cos\alpha t,0)$，第 $3$ 列是 $\vect e_3$。

(b) 记 $c=\cos\alpha t$、$s=\sin\alpha t$。直接乘：
\[
S^{\mathsf T}S=
\begin{pmatrix}
c&s&0\\-s&c&0\\0&0&1
\end{pmatrix}
\begin{pmatrix}
c&-s&0\\s&c&0\\0&0&1
\end{pmatrix}
=
\begin{pmatrix}
c^2+s^2&-cs+sc&0\\
-sc+cs&s^2+c^2&0\\
0&0&1
\end{pmatrix}
=\unitd .
\]
\clearpage
\answercontinue{P11}{P11 已写出 $S$ 的矩阵并验证 $S^{\mathsf T}S=\unitd$；本页接续完成 $SS^{\mathsf T}=\unitd$ 的核对，并做 (c) 的求导，记 $c=\cos\alpha t$、$s=\sin\alpha t$。}
同理 $SS^{\mathsf T}=\unitd$，所以 $S$ 是正交算子（\S3.3）。
(c) 求导：
\[
\frac{\dd S}{\dd t}=\alpha
\begin{pmatrix}
-\sin\alpha t&-\cos\alpha t&0\\
\cos\alpha t&-\sin\alpha t&0\\
0&0&0
\end{pmatrix}.
\]
于是
\[
\frac{\dd S}{\dd t}S^{\mathsf T}
=\alpha
\begin{pmatrix}
-s&-c&0\\ c&-s&0\\0&0&0
\end{pmatrix}
\begin{pmatrix}
c&s&0\\-s&c&0\\0&0&1
\end{pmatrix}
=\alpha
\begin{pmatrix}
0&-1&0\\1&0&0\\0&0&0
\end{pmatrix}
=\alpha\,\vect e_3^{\times}.
\]
最后一个等号因为
\[
\vect e_3^{\times}=
\begin{pmatrix}0&-1&0\\1&0&0\\0&0&0\end{pmatrix}.
\]
所以 $\boldsymbol\Omega^{\times}=\alpha\vect e_3^{\times}$，即
\[
\boldsymbol\Omega=\alpha\,\vect e_3 .
\]
这与“绕 $\vect e_3$ 以角速度 $\alpha$ 转”一致；方向沿转轴，大小是转角的变化率。

\answerstart{P12}{自编题}
{刚体 $A=1$、$B=2$、$C=3$ kg·m$^2$。某瞬时在体坐标系中 $\boldsymbol\omega=(1,1,1)$ rad/s，外力矩 $\vect M=\vect0$。用欧拉方程求此刻 $\dot\omega_1,\dot\omega_2,\dot\omega_3$。}

由 \S5.2 的欧拉方程（$M_k=0$）：
\[
\begin{aligned}
A\dot\omega_1+(C-B)\omega_2\omega_3&=0,\\
B\dot\omega_2+(A-C)\omega_3\omega_1&=0,\\
C\dot\omega_3+(B-A)\omega_1\omega_2&=0 .
\end{aligned}
\]
代入 $A=1,B=2,C=3$、$\omega_1=\omega_2=\omega_3=1$：
\begin{gather*}
1\cdot\dot\omega_1+(3-2)\cdot1\cdot1=0 \quad\Longrightarrow\quad \dot\omega_1=-1,\\
2\cdot\dot\omega_2+(1-3)\cdot1\cdot1=0 \quad\Longrightarrow\quad 2\dot\omega_2-2=0 \quad\Longrightarrow\quad \dot\omega_2=1,\\
3\cdot\dot\omega_3+(2-1)\cdot1\cdot1=0 \quad\Longrightarrow\quad 3\dot\omega_3+1=0 \quad\Longrightarrow\quad \dot\omega_3=-\frac13 .
\end{gather*}
所以
\[
\dot{\boldsymbol\omega}=\Bigl(-1,\;1,\;-\frac13\Bigr)\ \text{rad/s}^2 .
\]
核对：把三个结果代回三条方程，左边分别是 $1\cdot(-1)+1=0$、$2\cdot1-2=0$、$3\cdot(-1/3)+1=0$，都对。注意这是一个瞬时的变化率；下一时刻 $\boldsymbol\omega$ 变了，右边的陀螺项也会变。

\hdr{附录 B}{指标与记号速查}
\label{app:notation}

\aside{求和约定：同一项里重复出现两次的指标自动对 $1,2,3$ 求和。本稿正文常写求和号，附录与推导里用约定，避免“缩并了哪个指标”看错。
$\varepsilon_{ijk}$ 的取值见 \S0.2：偶排列 $+1$、奇排列 $-1$、有重复 $0$。RIT Whelan notes 10 的对偶张量定义为 $\ast A=-\vect a^{\times}$（符号相反），引用时先核对矩阵里每个分量的符号。}

\begin{center}
\begin{minipage}{\linewidth}
\small
\renewcommand{\arraystretch}{1.05}
\tabtop
\begin{tabular}{@{}ll@{}}
\thead{对象} & \thead{指标写法} \\
\multicolumn{2}{@{}l@{}}{\color{petal}\rule{0.55\linewidth}{0.5pt}}\\
并矢 & $D_{ij}=a_ib_j$ \\
右作用 & $(\ten D\cdot\vect x)_i=D_{ij}x_j$ \\
左作用 & $(\vect x\cdot\ten D)_j=x_iD_{ij}$ \\
张量乘张量 & $(\ten A\cdot\ten B)_{ij}=A_{ik}B_{kj}$ \\
双点积 & $\ten A:\ten B=A_{ij}B_{ij}$ \\
迹 & $\tr\ten D=D_{ii}$ \\
二次型 & $E=a_iD_{ij}a_j$ \\
变换律 & $D'_{ij}=R_{ki}R_{lj}D_{kl}$ \\
叉乘 & $(\vect a\times\vect b)_i=\varepsilon_{ijk}a_jb_k$ \\
对偶张量 & $(\vect a^{\times})_{ij}=\varepsilon_{ikj}a_k$ \\
转动惯量 & $J_{ij}=\int(r^2\delta_{ij}-x_ix_j)\dd m$ \\
散度 & $(\nabla\cdot\ten D)_i=\partial_jD_{ij}$ \\
\end{tabular}
\tabbot
\end{minipage}
\end{center}

\hdr{附录 C}{来源对应、补充与自补推导}
\label{app:sources}

\textbf{一句话：}本稿以课件 p8--p12 为主要来源，p15 的刚体速度关系作为已用前提；惯量张量部分与 RIT Whelan notes 10 的 pp.2--14 对照；欧拉方程、对称陀螺、应力与输运定理在课件对应页没有展开，本稿按自己的推导补出，并在下面逐项标清。旧稿只作为主题范围、题目身份与可用讲解的参考，不是事实依据。

\subhdr{C.1\quad 课件原页与本稿位置}

\begin{center}
\tabtop
\setlength{\tabcolsep}{9pt}%
\begin{tabular}{@{}p{0.10\linewidth}p{0.39\linewidth}p{0.39\linewidth}@{}}
\thead{课件页} & \thead{原页实际内容} & \thead{本稿位置} \\
\multicolumn{3}{@{}l@{}}{\color{petal}\rule{\linewidth}{0.5pt}}\\
p8 & 正交基、坐标 $a_k=a\cdot n_k$、求和约定、点积、叉乘、$\varepsilon_{ijk}$、对偶张量 $a^{\times}$ 及矩阵 & \S0.1、\S0.2、\S2.6、附录 B \\
p9 & 张量积（Dyad）、左右作用、转置、一般张量、线性算子 $b=Ta$、二次型 $E=aTa$、单位张量 & \S1.1、\S2.1--\S2.5、\S2.8 \\
p10 & 固定基与随体基、半径矢量两种展开、$S(t)=\sum_k n'_k\otimes n_k$ & \S3.1--\S3.2 \\
p11 & $S^{\mathsf T}S=SS^{\mathsf T}=\unitd$、保点积、绕轴转 $\varphi$ 的显式公式、无穷小旋转 & \S3.3--\S3.5 \\
p12 & $\boldsymbol\Omega^{\times}=(\dd S/\dd t)S^{\mathsf T}$、$\dd n'_k/\dd t=\boldsymbol\Omega\times n'_k$、无穷小情形 $\boldsymbol\Omega=\dd\varphi/\dd t$ & \S3.5、\S5.1 \\
p15 & 总结 $v=V+\boldsymbol\Omega\times r$；原点平移时 $V'=V+\boldsymbol\Omega\times a$、$\boldsymbol\Omega$ 不变；瞬时转轴 & \S3.6、\S4.1 \\
p13--p18 & 欧拉角、E8.01/E8.02、HW8.01--HW8.03 & 不属于本专题，未扩入；需要时按课件与运动学讲义处理 \\
\end{tabular}
\tabbot
\end{center}

\textbf{旧元数据的一处纠正。}旧稿的来源说明曾把“对称陀螺两方程”和“Gauss 并矢版”挂到课件 p12；核对原页后，p12 只有旋转算子的求导与角速度定义，没有这两项。本稿因此把欧拉方程、对称陀螺、应力平衡都标为“本稿自补推导”，不伪称课件原有。这一纠正也写进随附的 \texttt{source-map.json}。

\subhdr{C.2\quad 旧稿主题、练习与本稿位置}

\begin{center}
\tabtop
\begin{tabular}{@{}p{0.30\linewidth}p{0.54\linewidth}@{}}
\thead{旧稿主题} & \thead{本稿位置与处理} \\
\multicolumn{2}{@{}l@{}}{\color{petal}\rule{\linewidth}{0.5pt}}\\
旧 \S1--\S2（矩阵、换基、迹、行列式、特征多项式） & 重组为 \S0.2--\S0.6；换基改为“三步推导 + 核对”，并明确 $\ten R$ 的列约定 \\
旧 \S3（二次型、椭球、特征值、谱定理） & 重组为 \S1；谱定理证明从选读提升为正文，补 $2\times2$ 引理与正交补降维 \\
旧 \S4（阶、张量积、指标、变换律、商法则、上下标） & \S2.1、\S2.2、\S2.7；上下标移入选读（\S2.8 末） \\
旧 \S5（并矢、左右作用、点积与双点积、单位与投影、转置与对称分解、对偶张量、缩并表、谱分解） & \S2.3--\S2.6、\S2.8；缩并表作为查阅表；谱分解与 \S1.3 的例子对上 \\
旧 \S6（三记号字典） & \S2.8，保留并补充“并矢不是矩阵”的说明 \\
旧 \S7（惯量张量） & \S4，补出半正定性、三角形不等式、对称轴判断与平行轴推导；惯量椭球标为选读 \\
旧 \S8（输运定理、欧拉方程、对称陀螺） & \S5，补出并矢版输运定理的逐项推导与两个核对；对称陀螺教学示范换数字 \\
旧 \S9（应力张量） & \S6，补出柯西四面体论证、切应力力矩平衡、指标顺序与平衡方程 \\
旧 \S10（四道例题） & 改为独立练习 P1--P4，详解在附录 A；教学示范另换数字，避免提前曝光 \\
旧 \S11（易错与速查） & \S8，改为八条，并与正文交叉引用 \\
旧 \S12（缺口清单） & 本附录 C.3、C.4；机器可读版在 \texttt{source-map.json} \\
旧附录 A（指标速查） & 附录 B \\
旧附录 B（自测六问） & 独立练习 P5--P10，答案移入附录 A \\
旧稿练习 P11、P12 & 保留题号与题型，分别落在 \S3、\S5 的操作上，详解在附录 A \\
\end{tabular}
\tabbot
\end{center}

\subhdr{C.3\quad 本稿自补推导（课件未展开）}

以下各项不是课件原页的现成结论，而是本稿按已有定义与定理补出的推导；标“对照”的只表示与 Whelan 讲义的主题对应，不表示照抄。

\gapitem{惯量张量的完整构造；对照 Whelan notes 10 pp.2--14}{从 $\vect v=\boldsymbol\omega\times\vect r$ 与 $\vect L=\int\vect r\times\vect v\dd m$ 出发，用三矢量叉乘公式得到 $\ten J=\int(r^2\unitd-\vect r\otimes\vect r)\dd m$。课件 p9 只把二阶张量读成线性算子与二次型，没有给出这个构造。}

\gapitem{平行轴定理}{本稿在 \S4.5 展开 $\ten J_O$，用质心定义 $\int\vect r'\dd m=\vect0$ 消去交叉项；Whelan pp.13--14 有同一主题的对照，但本稿的每一步都按自己的记号重写。}

\gapitem{输运定理（矢量版与并矢版）}{课件 p12 只给 $\dd\vect n'_k/\dd t=\boldsymbol\Omega\times\vect n'_k$，没有给固定系/体系的导数关系。本稿在 \S5.1 逐项求导，并用单并矢与单位张量两个核对。}

\gapitem{欧拉方程与对称陀螺}{课件 p12 没有欧拉方程，也没有对称陀螺的两条方程。本稿在 \S5.2 从角动量定理与输运公式推出欧拉方程，在 \S5.3 对 $A=B\ne C$、$\vect M=\vect0$ 解出横向匀速转动。这是本稿自补，不归给课件。}

\gapitem{柯西应力、应力对称性与平衡方程}{课件没有应力章。本稿在 \S6.1 用四面体面积投影与力平衡得到 $\vect t=\ten\sigma\cdot\vect n$，在 \S6.2 用小立方体切应力力矩平衡得到 $\sigma_{ij}=\sigma_{ji}$，在 \S6.4 用小立方体净面力与高斯定理得到 $\nabla\cdot\ten\sigma+\vect f=\vect0$。Whelan notes 10 也没有应力章，故只标“本稿推导”。}

\subhdr{C.4\quad 选读清单与“跳过检查”}

本稿的选读只有三处：
\begin{enumerate}[label=(\arabic*),leftmargin=1.6em]
\item \S2.8 末尾的上下标/协变逆变说明；
\item \S4.5 末尾的惯量椭球；
\item \S6.4 末尾的高斯定理说明。
\end{enumerate}
\textbf{跳过检查：}遮住这三处，主线仍然读得通——\S2.8 的字典在选读之前，\S3 之后只用下标；\S4.6 与 \S5 不用惯量椭球；\S6.4 的平衡方程在选读之前，P12 与 P10 不依赖高斯定理的证明。三处都在末尾写了“跳过后从哪里继续”，没有把后文要用的定义藏在选读里。

\callout{\textbf{交付说明。}本稿的数值结果、矩阵乘法、特征值、惯量、欧拉方程与应力分解由随附的 \texttt{verify\_handout.py} 复核，输出见 \texttt{verification.txt}；静态检查用 \texttt{check\_handout.py}，编译日志见 \texttt{handout.log}。\textbf{讲义交付不等于学习者已掌握，也不等于外部审核已通过}；页面目视验收与用户试读反馈是另外两件事，见交付报告。本稿没有收到用户对本次整本的实际试读反馈，报告里不编造。}

\flowerrule

{\small\color{inkgray}本稿为 2026-11-11 的 round-01 独立首稿：主题范围保留，章节、例子、图与表述按本轮依赖重新组织；旧稿与已验收片段只作为可用材料，不原样拼接。}

\end{document}

